%-----------------------------------------------------------------------------
% Master Prefix Router for Randomized Algorithms Class Notes & Book
%
% Dispatches to:
%   - prefix_lwarp.tex when warping/HTML mode is active (lwarpmk html or \warpMode)
%   - prefix_4ht.tex when TeX4ht/make4ht defines \HCode
%   - prefix_latex.tex for standard LaTeX / XeLaTeX / LuaLaTeX PDF builds
%-----------------------------------------------------------------------------

% Include only once - second time... igmore
\ifdefined\PrefixLoadedAlready
  \endinput
\fi

\def\PrefixLoadedAlready{}


\makeatletter
\def\input@path{{styles/}}
\makeatother


\makeatletter
\def\LWR@isWarp{0}%
\ifx\warpMode\undefined\else\def\LWR@isWarp{1}\fi
\expandafter\ifx\csname opt@lwarp.sty\endcsname\relax\else\def\LWR@isWarp{1}\fi

\if1\LWR@isWarp
  \def\warpMode{1}%
  \IfFileExists{styles/prefix_lwarp.tex}{%-----------------------------------------------------------------------------
% Master Prefix File for Lwarp (HTML Web Generation)
% Compatible with LuaLaTeX and Lwarp
%-----------------------------------------------------------------------------

\RequirePackage{iftex}
\ifLuaTeX
  % LuaLaTeX verified
\else
  \PackageError{lwarp-gate}{%
    prefix_lwarp requires LuaLaTeX! Please run with lualatex.%
  }{HTML generation via lwarp must be run with LuaLaTeX.}
  \stop
\fi

% Always use memoir
\documentclass[12pt,oneside,fleqn]{memoir}

% Load lwarp immediately after documentclass
\usepackage[mathjax]{lwarp}

% Lwarp settings: split at chapter level (0) with section names
\setcounter{FileDepth}{0}
\booltrue{FileSectionNames}
\setcounter{SideTOCDepth}{1}

\makeatletter
% Multi-page SVG resolving hook for Ipe drawings with page=N
\def\LWR@current@figpage{}
\define@key{igraph}{page}{\def\LWR@current@figpage{#1}}
% Perpage compatibility for auxiliary files
\providecommand{\pp@pagectr}[4]{}
\AtBeginDocument{
  \ifdef{\LWR@nullfonts}{%
    \appto\LWR@nullfonts{%
      \def\xspace{ }%
      \def\Turan{Turán}%
      \def\Lovasz{Lovász}%
      \def\Erdos{Erdős}%
      \def\Holder{Hölder}%
      \def\Matousek{Matoušek}%
    }%
  }{}%
  \ifdef{\LWR@replaceEPSSVG}{
    \appto\LWR@replaceEPSSVG{%
      \ifx\LWR@current@figpage\@empty\else
        \edef\LWR@tempone{\LWR@tempone_p\LWR@current@figpage}%
        \def\LWR@current@figpage{}%
      \fi
    }%
  }{}%
}
\makeatother

%-----------------------------------------------------------------------------
% Typography, Fonts & Math Packages
%-----------------------------------------------------------------------------
\usepackage{mathtools}
\PassOptionsToPackage{warnings-off={mathtools-overbracket,mathtools-colon}}{unicode-math}
\usepackage{fontspec}
\usepackage{unicode-math}

% New Computer Modern OpenType text and math fonts
\setmainfont{NewCM10-Book.otf}[
  BoldFont          = NewCM10-Bold.otf,
  ItalicFont        = NewCM10-BookItalic.otf,
  BoldItalicFont    = NewCM10-BoldItalic.otf,
  SlantedFont       = NewCM10-Book.otf,
  SlantedFeatures   = {FakeSlant = 0.25},
  BoldSlantedFont   = NewCM10-Bold.otf,
  BoldSlantedFeatures = {FakeSlant = 0.25},
  SmallCapsFeatures = {Letters = SmallCaps}
]
\setsansfont{NewCMSans10-Book.otf}[
  BoldFont          = NewCMSans10-Bold.otf,
  ItalicFont        = NewCMSans10-BookOblique.otf,
  BoldItalicFont    = NewCMSans10-BoldOblique.otf
]
\setmonofont{NewCMMono10-Book.otf}[
  BoldFont          = NewCMMono10-Bold.otf,
  ItalicFont        = NewCMMono10-BookItalic.otf,
  BoldItalicFont    = NewCMMono10-BoldOblique.otf
]

\setmathfont{NewCMMath-Book.otf}
\setmathfont{NewCMMath-Bold.otf}[version=bold]

% OpenType math aliases and fixes
\providecommand{\square}{\mdlgwhtsquare}
\def\Box{\mdlgwhtsquare}
\def\Diamond{\mdlgwhtdiamond}
\def\blacksquare{\mdlgblksquare}
\def\leadsto{\rightsquigarrow}
\providecommand{\llbracket}{\lBrack}
\providecommand{\rrbracket}{\rBrack}
\DeclareTextCommand{\textcircled}{TU}[1]{%
  \mbox{\ooalign{\hfil\raise0.1ex\hbox{#1}\hfil\crcr\char"25EF}}%
}

\usepackage{alphabeta}
\usepackage[normalem]{ulem} % Provides \sout
\usepackage{mleftright}
\usepackage{microtype}
\usepackage{xspace}
\usepackage{xurl}
\usepackage[dvipsnames,table]{xcolor}
\usepackage{graphicx}
\usepackage{caption}
\usepackage{booktabs,array,colortbl,hhline,multirow,tabularx}
\usepackage[inline]{enumitem}
\usepackage[thicklines]{cancel}
\renewcommand\CancelColor{\color{purple}}
\usepackage{euscript}
\usepackage{marvosym}
\let\Cross\undefined
\let\Square\undefined
\usepackage{bbding}
\usepackage{epigraph}
\usepackage{makeidx}
\usepackage{perpage}
\usepackage{tikz}

%-----------------------------------------------------------------------------
% Boondox / mathcalb font declaration
%-----------------------------------------------------------------------------
\DeclareFontFamily{U}{BOONDOX-calo}{\skewchar\font=45 }
\DeclareFontShape{U}{BOONDOX-calo}{m}{n}{
  <-> s*[1.05] BOONDOX-r-calo}{}
\DeclareFontShape{U}{BOONDOX-calo}{b}{n}{
  <-> s*[1.05] BOONDOX-b-calo}{}
\DeclareMathAlphabet{\mathcalb}{U}{BOONDOX-calo}{m}{n}
\SetMathAlphabet{\mathcalb}{bold}{U}{BOONDOX-calo}{b}{n}
\DeclareMathAlphabet{\mathbcalb}{U}{BOONDOX-calo}{b}{n}

%-----------------------------------------------------------------------------
% Algorithm Environment (formerly salgorithm.sty)
%-----------------------------------------------------------------------------
\newcommand{\ProgWordX}[1]{{\small\textbf{#1}}{\ }}
\newcommand{\Do}{\ProgWordX{do}\xspace}
\newcommand{\End}{\ProgWordX{end}\xspace}
\newcommand{\Return}{\ProgWordX{return}\xspace}%
\newcommand{\Procbegin}{\ProgWordX{begin}}
\newcommand{\If}{\ProgWordX{if}\xspace}
\newcommand{\Then}{\ProgWordX{then}\xspace}
\newcommand{\Else}{\ProgWordX{else}\xspace}
\newcommand{\For}{\ProgWordX{for}\xspace}
\newcommand{\While}{{\small\bfseries while}\ }

\newcommand{\xbeginlgox}{\begin{minipage}{1in}\begin{tabbing}
           \quad\=\qquad\=\qquad\=\qquad\=\qquad\=\qquad\=\qquad\=\kill}
\newcommand{\xendlgox}{\end{tabbing}\end{minipage}}

\newenvironment{algorithmEnv}{%
   \begin{tabular}{|l|}\hline\xbeginlgox\bigskip}%
   {\xendlgox\\\hline\end{tabular}}

\newenvironment{program}{%
   \begin{minipage}{4.0in}%
   \begin{tabbing}%
       \ \ \ \ \=\ \ \ \=\ \ \ \ \=\ \ \ \ \=\ \ \ \ \=%
      \ \ \ \ \=\ \ \ \ \=\ \ \ \ \=\ \ \ \ \=%
      \ \ \ \ \=\ \ \ \ \=\ \ \ \ \=\ \ \ \ \=\kill%
}{%
   \end{tabbing}%
   \end{minipage}%
}

\newcommand{\AlgorithmBox}[1]{%
   \fbox{\begin{program} #1 \end{program}}%
}

\definecolor{AlgorithmColor}{cmyk}{0.07,0.90,0,0.34}
\providecommand{\Algorithm}[1]{{\AlgorithmI{#1}\index{algorithm!#1@{\AlgorithmI{#1}}}}}
\newcommand{\CodeComment}[1]{\textcolor{blue}{\texttt{#1}}}

%-----------------------------------------------------------------------------
% Path & Directory Resolution
%-----------------------------------------------------------------------------
\newcommand{\RootDIR}{../}
\IfFileExists{../../../book.tex}{\renewcommand{\RootDIR}{../../../}}{}
\IfFileExists{../../book.tex}{\renewcommand{\RootDIR}{../../}}{}
\IfFileExists{../book.tex}{\renewcommand{\RootDIR}{../}}{}
\IfFileExists{./book.tex}{\renewcommand{\RootDIR}{}}{}

\providecommand{\File}[1]{#1}
\providecommand{\RootFile}[1]{\RootDIR #1}

%-----------------------------------------------------------------------------
% QED Duck Symbol
%-----------------------------------------------------------------------------
\newlength{\CharHeight}
\AtBeginDocument{\setlength{\CharHeight}{\fontcharht\font`X}}

\IfFileExists{styles/qed_duck.svg}{%
   \providecommand{\myqedsymbol}{%
      \reflectbox{\includegraphics[height=1.4\CharHeight]{styles/qed_duck}}%
   }%
}{%
   \providecommand{\myqedsymbol}{$\blacksquare$}%
}

%-----------------------------------------------------------------------------
% BibLaTeX Setup (Pure Biber, No BibTeX)
%-----------------------------------------------------------------------------
\usepackage[bibencoding=utf8,style=alphabetic,refsegment=chapter,
indexing=cite,
backend=biber]{biblatex}

\ExecuteBibliographyOptions{%
   maxnames = 40,%
   maxalphanames = 4,%
   maxcitenames = 3,%
   minalphanames = 3,%
   minnames = 1,%
   giveninits = true,%
   doi = false,%
   url = false,%
   isbn = false%
}

\renewbibmacro{in:}{}

\AtEveryBibitem{%
   \ifentrytype{inproceedings}{%
      \clearname{editor}%
      \clearlist{editor}%
      \clearfield{editor}%
      \clearlist{publisher}%
      \clearfield{series}%
      \clearlist{location}%
      \clearfield{publisher}%
   }{}%
   \ifentrytype{book}{%
      \clearname{publisher}%
   }{}%
}
\AtEveryCitekey{%
   \ifentrytype{inproceedings}{%
      \clearlist{publisher}%
   }{}%
   \clearname{editor}%
   \clearlist{location}%
}
\DeclareFieldFormat[inproceedings]{date}{}
\DeclareFieldFormat[inproceedings]{pages}{#1, \printfield{year}}
\DeclareFieldFormat[inproceedings]{booktitle}{\mkbibemph{#1,\nopunct}}

\DeclareFieldFormat{journaltitle}{\mkbibemph{#1}\isdot}

\DeclareFieldFormat[article]{number}{\mkbibparens{#1}}
\DeclareFieldFormat[article]{date}{}
\DeclareFieldFormat[article]{pages}{#1, \printfield{year}}
\DeclareFieldFormat[article]{issue}{\nopunct#1}
\renewbibmacro*{issue+date}{%
   \iffieldundef{issue}{}%
   {%
      \printtext{\unspace}%
      (\printfield{issue})%
   }%
   \iffieldundef{pages}%
   {\addcomma\printfield{year}}%
   {}%
}

\renewcommand{\bibpagespunct}{%
   \ifentrytype{article}{\addcolon}{\addcomma}\space%
}

\DeclareFieldFormat[article]{titlecase}{\MakeSentenceCase{#1}}
\DeclareFieldFormat[inproceedings]{titlecase}{\MakeSentenceCase{#1}}
\DeclareFieldFormat[incollection]{titlecase}{\MakeSentenceCase{#1}}
\DeclareFieldFormat[book]{titlecase}{\MakeSentenceCase{#1}}

\renewbibmacro*{journal}{%
   \iffieldundef{journaltitle} {} {\printtext[journaltitle]{%
         \printfield[myplain]{journaltitle}%
         \setunit{\subtitlepunct}%
         \printfield[myplain]{journalsubtitle}}}}
\DeclareFieldFormat{myplain}{#1}

\renewbibmacro*{booktitle}{%
   \ifboolexpr{ test {\iffieldundef{booktitle}} and test
      {\iffieldundef{booksubtitle}} } {} {\printtext[booktitle]{%
         \printfield[myplain]{booktitle}%
         \setunit{\subtitlepunct}%
         \printfield[myplain]{booksubtitle}}%
      \newunit}%
   \printfield{booktitleaddon}}

\renewbibmacro*{volume+number+eid}{%
   \setunit*{\addcomma\addspace}%
   \printfield{volume}%
   \printfield{number}%
   \setunit{\addcomma\space}%
   \printfield{eid}%
}

\newbibmacro{string+doi}[1]{%
   \iffieldundef{doi}{%
      \iffieldundef{url}{%
         #1%
      }{%
         \href{\thefield{url}}{#1}%
      }%
   }{%
      \href{http://dx.doi.org/\thefield{doi}}{#1}}%
}
\DeclareFieldFormat*{title}{\usebibmacro{string+doi}{\mkbibemph{#1}}}
\renewcommand*{\multicitedelim}{\addcomma\space}
\DeclareFieldFormat[inproceedings]{title}{\usebibmacro{string+doi}{{#1}}}
\DeclareFieldFormat[article]{title}{\usebibmacro{string+doi}{{#1}}}
\DeclareFieldFormat[incollection]{title}{\usebibmacro{string+doi}{{#1}}}
\DeclareFieldFormat[book]{titlecase}{{#1}}

\DeclareCiteCommand{\tcite}%
{\usebibmacro{prenote}}%
{\usebibmacro{citeindex}\usebibmacro{cite}}%
{\multicitedelim}%
{\usebibmacro{postnote}}

%-----------------------------------------------------------------------------
% Hyperref & Visual Formatting
%-----------------------------------------------------------------------------
\usepackage{hyperref}
\hypersetup{%
   breaklinks,%
   colorlinks=true,%
   linkcolor=[rgb]{0.45,0.0,0.0},%
   citecolor=[rgb]{0,0,0.45}%
}

\definecolor{DeepNavy}{HTML}{002366}
\definecolor{RoyalBlue}{HTML}{004E92}
\definecolor{SteelBlue}{HTML}{4682B4}
\definecolor{TealBlue}{HTML}{008080}
\definecolor{DarkSlate}{HTML}{2F4F4F}
\definecolor{DustyGrey}{HTML}{696969}

\makeatletter
\AtBeginDocument{\let\@elt\relax}
\makeatother

\headstyles{bringhurst}
\chapterstyle{dash}

\renewcommand{\chapnamefont}{\normalfont\huge\bfseries\color{DeepNavy}}
\renewcommand{\chapnumfont}{\normalfont\huge\bfseries\color{DeepNavy}}
\renewcommand{\chaptitlefont}{\normalfont\Huge\bfseries\color{DeepNavy}}

\counterwithin{footnote}{chapter}

\setsubparaheadstyle{\normalfont\normalsize\itshape\color{DustyGrey}}
\setlength{\beforechapskip}{-20pt} 
\setlength{\afterchapskip}{30pt}
\setsecnumformat{\csname the#1\endcsname.\quad}

%-----------------------------------------------------------------------------
% Standalone Chapter vs Master Book Modes
%-----------------------------------------------------------------------------
\newcommand{\Copyright}{This work is licensed under the Creative Commons Attribution-Noncommercial 3.0 License. To view a copy of this license, visit \url{http://creativecommons.org/licenses/by-nc/3.0/} or send a letter to Creative Commons, 171 Second Street, Suite 300, San Francisco, California, 94105, USA.}

\newcommand{\RealChapterEndInner}{%
   {%
      \begingroup%
      \emergencystretch=1em%
      \printbibliography[heading=subbibliography,segment=\therefsegment]{}%
      \endgroup%
   }%
}

\providecommand{\BookMode}[1]{}
\providecommand{\StandAloneMode}[1]{#1}

\ifx\bookMode\undefined
   % STANDALONE CHAPTER MODE
   \renewcommand{\StandAloneMode}[1]{#1}
   \renewcommand{\BookMode}[1]{}

   \newcommand{\ChapterNumPage}[2]{%
      \if\relax\detokenize{#1}\relax
         \setcounter{chapter}{0}%
      \else\ifnum#1>0
         \setcounter{chapter}{\numexpr#1-1\relax}%
      \else
         \setcounter{chapter}{0}%
      \fi\fi
      \if\relax\detokenize{#2}\relax
         \setcounter{page}{1}%
      \else\ifnum#2>0
         \setcounter{page}{#2}%
      \else
         \setcounter{page}{1}%
      \fi\fi
   }

   \DeclareRobustCommand{\setheadingcolor}{\color{RoyalBlue}}
   \setsecheadstyle{\setheadingcolor\normalfont\Large\bfseries}
   \setsubsecheadstyle{\normalfont\large\bfseries\color{SteelBlue}}
   \setsubsubsecheadstyle{\normalfont\normalsize\bfseries\color{TealBlue}}
   \setparaheadstyle{\normalfont\normalsize\bfseries\color{DarkSlate}}

   \DeclareDocumentCommand{\Chapter}{s o m}{%
      \IfBooleanTF{#1}{%
         \IfValueTF{#2}{\chapter*[]{#2}}{\chapter*[]{#3}}%
      }{%
         \IfValueTF{#2}{\chapter*[]{#2}}{\chapter[]{#3}}%
      }%
      \setcounter{footnote}{0}%
   }
   \newcommand{\ChapterPrefix}{}
   \newcommand{\ChapterGen}{\pagestyle{plain}\vspace*{-0.7cm}}
   \newcommand{\ChapterEnd}{%
      \RealChapterEndInner{}%
      \end{document}%
   }

   \renewcommand{\File}[1]{#1}
   \renewcommand{\RootFile}[1]{../#1}
\else
   % MASTER BOOK MODE
   \renewcommand{\StandAloneMode}[1]{}
   \renewcommand{\BookMode}[1]{#1}

   \newcommand{\ChapterNumPage}[2]{}

   \DeclareDocumentCommand{\Chapter}{s o m}{%
      \IfBooleanTF{#1}{%
         \IfValueTF{#2}{\chapter*[#2]{#3}}{\chapter*{#3}}%
         \ifx\Dir\undefined\else
            \typeout{CHAPTER_MAP: \Dir/\CurrentChapFile => chap: 0, page: \thepage}%
         \fi
      }{%
         \IfValueTF{#2}{\chapter[#2]{#3}}{\chapter{#3}}%
         \ifx\Dir\undefined\else
            \typeout{CHAPTER_MAP: \Dir/\CurrentChapFile => chap: \thechapter, page: \thepage}%
         \fi
      }%
   }
   \newcommand{\ChapterPrefix}{%
      \vspace{-1cm}%
      \subsection*{\small
         {598 - Class notes for Randomized Algorithms}\\
         {Sariel Har-Peled}\\
         {\today}}%
   }
   \newcommand{\ChapterGen}{}
   \newcommand{\ChapterEnd}{\RealChapterEndInner{}}

   \renewcommand{\File}[1]{\Dir/#1}
   \renewcommand{\RootFile}[1]{#1}

   \newcommand{\InputChap}[2]{%
      \def\Dir{#1}%
      \def\CurrentChapFile{#2}%
      \typeout{}%
      \typeout{------------------------------------------------------------}%
      \typeout{PROCESSING FILE: #1/#2}%
      \typeout{}%
      \input{#1/#2}%
   }
\fi

%-----------------------------------------------------------------------------
% Fragments & Quotes
%-----------------------------------------------------------------------------
\newcommand{\InputExt}[1]{%
   \IfFileExists{#1}%
   {\input{#1}}%
   {\IfFileExists{../#1}%
      {\input{../#1}}%
      {\IfFileExists{../../#1}%
         {\input{../../#1}}%
         {\typeout{ERROR: Unable to find #1}}%
      }%
   }%
}
\makeatletter
\def\FragNormKey#1{%
  \frag@norm@check#1def_\relax{#1}%
}
\def\frag@norm@check#1def_#2\relax#3{%
  \if\relax\detokenize{#1}\relax
    #3%
  \else
    def_#3%
  \fi
}
\makeatother

\newcommand{\FragmentName}[1]{\RootFile{fragment/\FragNormKey{#1}.tex}}

\newenvironment{fragment}[1]{}{}

\makeatletter
\AtBeginDocument{%
   \let\real@latex@label\label
   \def\label#1{%
      \global\@namedef{all@labels@seen@#1}{1}%
      \real@latex@label{#1}%
   }%
}

\DeclareDocumentCommand{\IncFragment}{m}{%
   \begingroup
      \def\label##1{%
         \@ifundefined{all@labels@seen@##1}{%
            \global\@namedef{all@labels@seen@##1}{1}%
            \real@latex@label{##1}%
         }{}%
      }%
      \IfFileExists{\FragmentName{#1}}%
         {\InputExt{\FragmentName{#1}}}%
         {\InputExt{\RootFile{fragment/#1.tex}}}%
   \endgroup
}
\makeatother

\newcommand{\Quote}[2]{%
   \SaveIndent
   \begin{center}%
      {\footnotesize
         \begin{quote}
             \RestoreIndent
             {#1}\\
             \hspace*{2cm}-- {#2.}
         \end{quote}%
      }%
      \index{quotation!#2}%
   \end{center}%
   \bigskip
}

\newcommand{\QuotePExt}[3]{%
   \bigskip%
   \hfill%
   \begin{minipage}{0.95\linewidth}%
       \footnotesize {#1}%

       \rule{\linewidth}{0.3mm}%
       \hfill #2, #3%
   \end{minipage}%
   \bigskip%
}

\epigraphfontsize{\small\itshape}
\setlength{\epigraphwidth}{0.6\textwidth}
\newcommand{\QuoteExt}[3]{%
   \vspace*{-1.6cm}%
   \epigraph{{\footnotesize #1}}{{#2}, {#3}}%
   \vspace*{-1.2cm}%
}

\providecommand{\delX}[1]{{\color{red}\sout{#1}}}
\providecommand{\newX}[1]{{\color{blue}#1}}
\providecommand{\chgY}[2]{{\color{red}\sout{#1}}{\color{blue}#2}}
\providecommand{\remX}[1]{{\color{purple}\textbf{[Remark: #1]}}}
\newcommand{\SpellIgnore}[1]{#1}
\newcommand{\EQ}{\mathrel{\;\;=\;\;}}
\newcommand{\LEQ}{\mathrel{\;\;\leq\;\;}}
\newcommand{\LT}{\mathrel{\;\;<\;\;}}
\newcommand{\GT}{\mathrel{\;\;<\;\;}}

\newcommand{\GEQ}{\mathrel{\;\;\geq\;\;}}

%-----------------------------------------------------------------------------
% Save a symbol that we know is going to get redefined.
\def\savesymbol#1{%
  \expandafter\let\expandafter\origsym\expandafter=\csname#1\endcsname
  \expandafter\let\csname orig#1\endcsname=\origsym
  \expandafter\let\csname#1\endcsname=\relax
}

% Restore a previously saved symbol, and rename the current one.
\def\restoresymbol#1#2{%
  \expandafter\let\expandafter\newsym\expandafter=\csname#2\endcsname
  \expandafter\global\expandafter\let\csname#1#2\endcsname=\newsym
  \expandafter\let\expandafter\origsym\expandafter=\csname orig#2\endcsname
  \expandafter\global\expandafter\let\csname#2\endcsname=\origsym
}

\usepackage{pictex}%
\savesymbol{triangle}%
\usepackage{tree}%
\restoresymbol{TreePKG}{triangle}%

% Definitions absorbed from sbook.sty
\providecommand{\MdPrc}{\nu}

\StartDefiningTabulars
\providecommand{\TwoFiguresBase}[4]{%
   \begin{tabular}{cc}
      \begin{minipage}{0.48\linewidth}
          \begin{center}
              {{{#1}}}
          \end{center}
      \end{minipage}
      &
      \begin{minipage}{0.48\linewidth}
          \begin{center}
              {{#3}}
          \end{center}
      \end{minipage}
      \\
      {#2} & {#4}
   \end{tabular}%
}

\providecommand{\TwoFigures}[6]{%
   \begin{tabular}{cc}
      \begin{minipage}{0.48\linewidth}
          \begin{center}
              {\IncludeGraphics[#1]{#2}}
          \end{center}
      \end{minipage}
      &
      \begin{minipage}{0.48\linewidth}
          \centerline{\IncludeGraphics[#4]{#5}}
      \end{minipage}
      \\
      {#3} & {#6}
   \end{tabular}%
}

\providecommand{\ThreeFigures}[6]{%
   \begin{tabular}{ccc}
        \begin{minipage}{0.30\linewidth}
            \IncludeGraphics[width=5.1cm]{#1}\\
        \end{minipage}
        &
        \begin{minipage}{0.30\linewidth}
            \IncludeGraphics[width=5.1cm]{#2}\\
        \end{minipage}
        &
        \begin{minipage}{0.30\linewidth}
            \IncludeGraphics[width=5.1cm]{#3}\\
        \end{minipage}
        \vspace{-1cm}
        \\
        {#4} & {#5}& {#6}
   \end{tabular}%
}
\StopDefiningTabulars

%-----------------------------------------------------------------------------
\providecommand{\remove}[1]{}
\providecommand{\NotAMSMode}[1]{#1}
\providecommand{\Solution}[1]{}
\providecommand{\Version}[1]{\hfill{Version: #1}}
\providecommand{\Distor}{\mathrm{dist}}
\providecommand{\Hardness}[1]{(h:#1)}
\providecommand{\QuoteOptExt}[3]{\QuoteExt{#1}{#2}{#3}}
\providecommand{\Hypercube}{\EuScript{H}}
\providecommand{\emphiC}[1]{\emphic{#1}{#1}}
\providecommand{\CoverTime}{\EuScript{C}}
\providecommand{\Nbr}{\Gamma}
\providecommand{\CommuteTime}[2]{\mathbf{CT}_{#1#2}}
\providecommand{\Resistance}{\mathbf{R}}
\providecommand{\EffectResist}[2]{\mathbf{R}_{#1#2}}
\providecommand{\NetworkE}{\EuScript{N}}
\providecommand{\trinom}[4]{\binom{#1}{#2 \;\; #3 \;\; #4}}
\providecommand{\Ranges}{\EuScript{R}}
\providecommand{\CClass}[1]{\ensuremath{\mathbf{#1}}}
\providecommand{\coCClass}[1]{\ensuremath{\mathrm{co-}}\CClass{#1}}
\providecommand{\CX}{\EuScript{C}}
\providecommand{\packet}{\mathbf{c}}
\providecommand{\packetA}{\mathbf{d}}
\providecommand{\mopt}{m_{\mathrm{opt}}}
\providecommand{\girth}{\mathop{\mathrm{girth}}}
\providecommand{\CR}{\mathop{\mathrm{cr}}}
\providecommand{\TRUE}{{\textcolor[named]{RedViolet}{\textsf{TRUE}}}\xspace}
\providecommand{\FALSE}{{\textcolor[named]{RedViolet}{\textsf{FALSE}}}\xspace}
\providecommand{\Part}[1]{\part{#1}}
\providecommand{\Section}[1]{\section{#1}}
\providecommand{\SubSection}[1]{\subsection{#1}}
\providecommand{\SubSubSection}[1]{\subsubsection{#1}}
\providecommand{\Exercise}[5]{%
\begin{exercise_h}[#1]%
   \exelab{#3}{}%
   #4%
\end{exercise_h}%
}
\providecommand{\si}[1]{#1}
\providecommand{\myParagraph}[1]{\noindent\textbf{#1}}
\providecommand{\DefName}[1]{#1}
\providecommand{\deflabel}[1]{#1}
\providecommand{\Description}[1]{}

%-----------------------------------------------------------------------------
% Mathematical Notations & Environments (from notations.sty)
%-----------------------------------------------------------------------------
%-----------------------------------------------------------------------------
% List environments (questions & compactenum)
%-----------------------------------------------------------------------------

\newlist{compactenumA}{enumerate}{5}%
\setlist[compactenumA]{topsep=0pt,itemsep=-1ex,partopsep=1ex,parsep=1ex,%
   label=(\Alph*)}%

\newlist{compactenumI}{enumerate}{5}%
\setlist[compactenumI]{topsep=0pt,itemsep=-1ex,partopsep=1ex,parsep=1ex,%
   label=(\Roman*)}%

\newlist{compactenuma}{enumerate}{5}%
\setlist[compactenuma]{topsep=0pt,itemsep=-1ex,partopsep=1ex,parsep=1ex,%
   label=(\alph*)}%

\newlist{compactitem}{enumerate}{5}%
\setlist[compactitem]{topsep=0pt,itemsep=-1ex,partopsep=1ex,parsep=1ex,%
   label=\textbullet}%

\newlist{compactenumi}{enumerate}{5}%
\setlist[compactenumi]{topsep=0pt,itemsep=-1ex,partopsep=1ex,parsep=1ex,%
   label=(\roman*)}%

\setcounter{secnumdepth}{5}

\newcommand{\HST}{\Term{HST}\xspace}
\newcommand{\Term}[1]{\textsf{#1}}
\newcommand{\TermI}[1]{\Term{#1}\index{#1@\Term{#1}}}

\DeclareMathAlphabet{\mathpzc}{OT1}{pzc}{m}{it}

\newcommand{\seg}{\mathsf{s}}

\newcommand{\Diam}{\mathrm{diam}}
\newcommand{\diam}{\mathrm{diam}}

\newcommand{\Copt}{C_{\mathrm{opt}}}
\newcommand{\CenterPnt}{\mathrm{cen}}

\newcommand{\parent}{\overline{\mathrm{p}}}
\newcommand{\parentX}[2][\!]{\parent\pth[#1]{#2}}

\newcommand{\FMS}{\EuScript{X}}  
\newcommand{\FMSB}{\EuScript{Y}}  
\newcommand{\FastCut}{\Algorithm{FastCut}\xspace}
\newcommand{\MinCut}{\Algorithm{MinCut}\xspace}
\newcommand{\MinCutRep}{\Algorithm{MinCutRep}\xspace}
\newcommand{\Contract}{\Algorithm{Contract}\xspace}

\newcommand{\Grid}{\mathrm{G}}
\newcommand{\ClosestPair}{\mathcal{CP}}
\newcommand{\SphereChar}{\ensuremath{\mathbb{S}}}
\newcommand{\Sphere}[1]{{\SphereChar^{(#1)}}}

\newcommand{\PntSet}{{\mathsf{P}}}%
\providecommand{\PntSetA}{{\mathsf{Q}}}%

\newcommand{\query}{\mathcalb{q}}

\newcommand{\pnt}{\mathsf{p}}
\newcommand{\pntA}{\mathsf{u}}
\newcommand{\pntB}{\mathsf{v}}
\newcommand{\pntC}{{\mathsf{x}}}
\newcommand{\pntD}{{\mathsf{y}}}
\newcommand{\pntE}{{\mathsf{z}}}

\newcommand{\pntCoord}{p}
\newcommand{\pntCoordA}{q}

%%%%%%%%%%%%%%%%%%%%%%%%%%%%%%%%%%%%%%%%%%%%%%%%%%%%%%%%%%%%%%%%%%%%% 
%%%%%%%%%%%%%%%%%%%%%%%%%%%%%%%%%%%%%%%%%%%%%%%%%%%%%%%%%%%%%%%%%%%%%
% Markov chain notations.

\newcommand{\MCSpc}{\mathtt{S}}
\newcommand{\MCPrb}{\mathtt{P}}
\newcommand{\MCPrbY}[2]{\MCPrb_{#1#2}}
\newcommand{\MCPrbM}{\mathbf{P}}%
\newcommand{\MCPrbT}[3]{\MCPrb_{#1#2}^{(#3)}}
\newcommand{\MCFPrbT}[3]{\mathtt{r}_{#1#2}^{(#3)}}
\newcommand{\MCVstProb}[2]{\mathtt{f}_{#1#2}}
\newcommand{\MCExSteps}[2]{\mathtt{h}_{#1#2}}
\newcommand{\MCPVec}[1]{\mathbf{q}^{(#1)}}
%%%%%%%%%%%%%%%%%%%%%%%%%%%%%%%%%%%%%%%%%%%%%%%%%%%%%%%%%%%%%%%%%%%%%

\newcommand{\Spread}{\Phi}
\newcommand{\VC}{\ensuremath{\mathsf{VC}}\xspace}
\newcommand{\MC}{\ensuremath{\mathsf{MC}}\xspace}
\newcommand{\hyperplane}{h}

\newcommand{\QuickSort}{\Algorithm{QuickSort}\xspace}
\newcommand{\Alg}{\Algorithm{Alg}\xspace}

\newcommand{\Graph}{\mathsf{G}}%
\newcommand{\GraphB}{\mathsf{H}}
\newcommand{\GraphC}{\mathsf{K}}
\newcommand{\GraphD}{\mathsf{F}}

\providecommand{\G}{}
\renewcommand{\G}{\Graph}

\newcommand{\RLP}{\ComplexityClass{RLP}\xspace}
\newcommand{\USTCON}{\ComplexityClass{USTCON}}
\newcommand{\STCON}{\overrightarrow{\ComplexityClass{STCON}}}

\newcommand{\Mtrx}{\mathsf{M}}
\newcommand{\TMtrx}{\mathsf{Q}}
\newcommand{\IMtrx}{\mathcal{I}}

\newcommand{\Expander}[3]{$\left[#1,#2,#3\right]$-expander}

\newcommand{\ChartX}[1]{\ensuremath{\chi_{#1}}}
\newcommand{\Field}{\mathbb{F}}
\newcommand{\LSpace}{\mathbb{G}}
\newcommand{\LD}[2]{\mathsf{L{D}}\pth{#1,#2}}
\newcommand{\SEVN}{\widehat{\lambda_2}}
\newcommand{\NX}[2]{#1\!\left[#2\right]}
\newcommand{\Entry}[2]{#1\!\left[#2\right]}
\newcommand{\NXExt}[3]{#1_{#2}\!\left[#3\right]}
\newcommand{\RangeI}[1]{\llbracket #1 \rrbracket}
\newcommand{\DialX}[1]{\gamma\pth{#1}}
\newcommand{\BiDialX}[1]{{\gamma}_2\pth{#1}}

\newcommand{\EV}[1]{\widehat{\lambda_{#1}}}
\newcommand{\TEV}[1]{\widehat{\lambda_{#1}}}
\newcommand{\EVal}{\widehat{\lambda}}
\newcommand{\Basis}{\mathcal{B}}

\newcommand{\ZProd}[2]{\ensuremath{#1 \, \text{\textcircled{z}} #2}}
\newcommand{\RProd}[2]{\ensuremath{#1 \, \text{\textcircled{r}} #2}}

\newcommand{\OneVec}[1]{\mathbf{1}^{#1}}
\newcommand{\DstrbV}[1]{\mathbf{\mathsf{q}}^{(#1)}}
\newcommand{\DstrbE}[2]{\mathsf{q}^{(#1)}_{#2}}

\newcommand{\DEdges}{\widetilde{\Edges}}

\newcommand{\VV}{\mathsf{V}}%
\newcommand{\EE}{\mathsf{E}}%

\newcommand{\Vertices}{\VV}%
\newcommand{\Edges}{\EE}%

\newcommand{\EdgesX}[1]{\Edges\pth{#1}}
\newcommand{\VerticesX}[1]{V\pth{#1}}

\newcommand{\StringX}[1]{\mathsf{s}\pth{#1}}

\newcommand{\BMtrx}{\mathsf{B}}
\newcommand{\WMtrx}{\mathsf{W}}
\newcommand{\NWMtrx}{\overline{\mathsf{W}}}
\newcommand{\Sequence}{\EuScript{S}}

\newcommand{\med}{\mathrm{med}}
\newcommand{\medX}[1]{\mathrm{med}\pth{#1}}

\newcommand{\medianC}{\nu}%

\newlength{\savedparindent}
\newcommand{\SaveIndent}{\setlength{\savedparindent}{\parindent}}
\newcommand{\RestoreIndent}{\setlength{\parindent}{\savedparindent}}

\newcommand{\NNN}{\mathbb{N}}%
\providecommand{\NN}{\mathbb{N}}%
\renewcommand{\Re}{\mathbb{R}}%
\newcommand{\reals}{\Re}%
\newcommand{\ZZ}{\mathbb{Z}}%

\newcommand{\level}{\mathop{\mathrm{level}}}

\providecommand{\TPDF}[2]{\texorpdfstring{#1}{#2}}

\newcommand{\Hint}[1]{\emph{[#1]}}
\newcommand{\Arr}{\mathop{\mathrm{\EuScript{A}}}}

\newcommand{\Array}{\mathcal{T}}%

\newcommand{\dirEdge}[2]{\left({{#1} \rightarrow {#2}} \right)}

%%%%%%%%%%%%%%%%%%%%%%%%%%%%%%%%%%%%%%%%%%%%%%%%%%%%%%%%%%%%%%%%%%%%%%
% Emphasize stuff...
\definecolor{blue25}{rgb}{0, 0, 11}
\newcommand{\emphic}[2]{%
   \textcolor{blue25}{%
      \bfseries{\emph{#1}}}%
   \index{#2}}
\newcommand{\emphi}[1]{\emphic{#1}{#1}}
\newcommand{\emphOnly}[1]{\textcolor{blue25}{\bfseries{\emph{#1}}}}

\definecolor{almostblack}{rgb}{0, 0, 0.3}
\providecommand{\emphw}[1]{{\textcolor{almostblack}{\emph{#1}}}}%

%
%%%%%%%%%%%%%%%%%%%%%%%%%%%%%%%%%%%%%%%%%%%%%%%%%%%%%%%%%%%%%%%%%%%%%%

\newcommand{\volX}[1]{\mathrm{vol}\pth{#1}}

\newcommand{\Interval}{\EuScript{I}}
\newcommand{\IntervalA}{\EuScript{I}}

\newcommand{\Exp}{\mathrm{Exp}}

\newcommand{\binsY}[2]{\sqcup_{\geq #1}\pth{#2}}
\newcommand{\ballsY}[2]{\mathcal{F}_{\geq #1}\pth{#2}}

\newcommand{\BLX}[1]{{{\sqcupplus{}}_{\geq #1}}}


\let\Ball\undefined
\newcommand{\Ball}{\mathbf{b}}
\newcommand{\hdim}{d}
\newcommand{\Domain}{\mathcal{D}}
\newcommand{\points}[1]{[\bfseries{#1 points}]}
\newcommand{\Event}{\EuScript{E}}%
\newcommand{\EventB}{\EuScript{B}}%
\newcommand{\EventF}{\EuScript{F}}%
\newcommand{\EventG}{\EuScript{G}}%

\newcommand{\distCmd}[1]{\left\| {#1}  \right\|}
\newcommand{\distY}[2]{\distCmd{#1 - #2}}

\newcommand{\NormalC}{\mathbf{N}}
\newcommand{\NormalY}[2]{\NormalC\pth{#1, #2}}%
\newcommand{\NormalDistOneDim}{\NormalC}
\newcommand{\NormalDistX}[1]{\NormalC^{#1}}
\DeclareMathOperator{\lca}{lca}
\newcommand{\RangeSpace}{\mathsf{S}}%
\newcommand{\RangeSpaceA}{\mathsf{T}}

\newcommand{\range}{\mathbf{r}}
\newcommand{\rangeA}{\mathbf{r}'}

\newcommand{\BadProb}{\varphi}

\newcommand{\EntropyChar}{\mathbb{H}}
\newcommand{\Entropy}[1]{\EntropyChar\pth{#1}}
\newcommand{\EntropyX}[1]{\EntropyChar\pth[]{#1}}

\newcommand{\Matrix}{\mathsf{M}}
\newcommand{\MatrixB}{\mathsf{B}}
\newcommand{\MatrixC}{\mathsf{C}}
\newcommand{\MatrixD}{\mathsf{D}}

\renewcommand{\th}{th\xspace}

\newcommand{\Family}{\mathcal{F}}
\newcommand{\FamilyB}{\mathcal{G}}
\newcommand{\FamilyC}{\mathcal{H}}
\newcommand{\FMSC}{\EuScript{Y}\index{metric space}}

\newcommand{\median}{\overline{\mathsf{m}}}
\newcommand{\QuickSelect}{\Algorithm{Quick{}Select}\xspace}

\providecommand{\row}{\mathbf{r}}%
\newcommand{\rowc}{{r}}%

\newcommand{\vecb}{\mathbf{b}}

\newcommand{\vecA}{\mathbf{v}}
\newcommand{\vecB}{\mathbf{u}}
\newcommand{\vecRand}{\mathbf{r}}

\newcommand{\vecH}{\mathbf{h}}
\newcommand{\vecV}{\mathbf{v}}

\newcommand{\vecX}{\mathbf{x}}
\newcommand{\vecY}{\mathbf{y}}

\newcommand{\DistBinom}[2]{\mathrm{Bin}\pth{#1,#2}}

\newcommand{\DistGeom}[1]{\mathrm{Geom}\pth{#1}}

\newcommand{\Sample}{\mathsf{R}}

\newcommand{\lcm}{\mathop{\mathrm{lcm}}}
\newcommand{\bdiv}{\mathop{\mathrm{div}}}%
\newcommand{\divides}{\mid}%
\newcommand{\poly}{\mathrm{poly}}%

\newcommand{\Group}{\mathcal{G}} 
\newcommand{\GroupA}{\mathcal{H}} 

\newcommand{\algGCD}{\Algorithm{EuclidGCD}\xspace}%

\newcommand{\gOp}{\times}%
\newcommand{\gIdent}{\mathsf{i}}%

\newcommand{\ordX}[2][\!]{\mathop{\mathrm{o{r}d}}\pth[#1]{#2}}

\newcommand{\lSym}[2]{%
   \left(
       \ds {#1}\divides{#2}%
   \right) %
}
\newcommand{\jSym}[2]{\left\llbracket%
       \ds {#1}\divides{#2}%
       \right\rrbracket%
}
\newcommand{\Line}{\ell}

\newcommand{\LevelChar}{\mathbb{L}}
\newcommand{\AtMostKLevelChar}{\LevelChar_{\leq k}}
\newcommand{\AtMostKLevel}[1]{\LevelChar_{\leq k} \pth{ #1}}

\newcommand{\LineSet}{\mathsf{L}}%
\newcommand{\ArrX}[1]{\EuScript{A}\pth{#1}}%

\DeclareRobustCommand{\Matousek}{Matoušek\xspace}
\DeclareRobustCommand{\Erdos}{Erdős\xspace}
\DeclareRobustCommand{\Lovasz}{Lovász\xspace}

\newcommand{\ProblemC}[1]{{\textsf{\textcolor{red}{#1}\index{problem!#1}}}}
\DefineNamedColor{named}{OliveGreen} {cmyk}{0.64,0,0.95,0.40}
\providecommand{\ComplexityClass}[1]{{{\textcolor[named]{OliveGreen}{%
            \textsc{#1}}}}}
\providecommand{\NPHard}{{\ComplexityClass{NP-Hard}}%
   \index{NP!hard}\xspace}
\providecommand{\NPComplete}{\ComplexityClass{NP-Complete}%
   \index{NP!complete}\xspace}
\providecommand{\APXHard}{\ComplexityClass{APX-Hard}%
   \index{APX!Hard}\xspace}
\providecommand{\NPC}{\ComplexityClass{NPC}%
   \index{NP!complete}\xspace}
\providecommand{\NP}{\ComplexityClass{NP}%
   \index{NP}\xspace}
\providecommand{\POLYT}{\ComplexityClass{P}\xspace}

\newcommand{\envFrame}[1]{\begin{center}\fbox{\begin{minipage}[c]{0.75\textwidth}{#1}\end{minipage}}\end{center}}

\newcommand{\mysqrt}[1]{\ts\!\!\sqrt{#1}}
\newcommand{\SegSet}{\mathsf{S}}
\newcommand{\ArrVD}[1]{\Arr^{|} \pth{#1}}
\newcommand{\RIC}{\textsf{RIC}\index{random incremental
      construction}\xspace}
\newcommand{\ArrB}{\EuScript{B}}

\newcommand{\TreeOpt}{\EuScript{T}_\mathrm{opt}}
\newcommand{\TreeOptA}{\EuScript{T}_\mathrm{opt}'}

\newcommand{\DepthExt}[2]{\mathrm{depth}_{#1}\pth{#2}}
\newcommand{\depthX}[1]{\mathrm{depth}\pth{#1}}

\newcommand{\TreeA}{\EuScript{U}}
\newcommand{\TreeB}{\EuScript{V}}

\newcommand{\Trap}{\sigma}
\newcommand{\TrapA}{\sigma'}
\newcommand{\TrapB}{\sigma''}

\newcommand{\SampleA}{\mathsf{R}'}

\newcommand{\AllRegions}{\mathcal{T}}
\newcommand{\RegionsSet}[1]{\mathcal{F}\pth{#1}}
\newcommand{\RegionsHeavy}[2]{\mathcal{F}_{\!\geq#1}\pth{#2}}
\newcommand{\NumOfRegionsZero}[1]{\ExChar \!f\pth{#1}}
\newcommand{\NumOfRegions}[2]{\ExChar\! f_{\!\geq #1}\pth[]{#2}}
\newcommand{\ObjSet}{\mathsf{S}}

\newcommand{\WeightX}[2][\!]{\omega \pth[#1]{#2}}
\newcommand{\Edge}{\mathsf{e}}

\newcommand{\Face}{\mathsf{f}}

\newcommand{\DefSet}[1]{D\pth{#1}}
\newcommand{\KillSet}[1]{K\pth{#1}}

\newcommand{\CSProbFull}[4]{
   \,\rho_{\!\!\stackrel{\rule[-.005cm]{0cm}{0.1cm}}{#1,
       #4}}\pth{\MakeSBig #2,#3}}

\newcommand{\disk}{D}
\newcommand{\diskA}{D'}

\newcommand{\IntervalSet}{\widehat{\mathcal{I}}}

\newcommand{\GroundSet}{\textsf{X}}
\newcommand{\FGroundSet}{\mathsf{x}}
\newcommand{\RangeSet}{{\mathcal{R}}}
\newcommand{\RangeSetA}{\mathcal{R}'}
\newcommand{\MeasureChar}{\overline{m}}
\newcommand{\Measure}[1]{\MeasureChar\pth{#1}}
\newcommand{\sMeasure}[1]{\overline{s}\pth{#1}}

\newcommand{\SetA}{S}%
\newcommand{\SetB}{T}%
\newcommand{\SetC}{N}%

\newcommand{\VCProj}[2]{#1_{|#2}}
\newcommand{\VCDim}{\mathrm{dim}_{\VC}}
\newcommand{\Dim}{\delta}
\newcommand{\DimA}{\delta'}
\newcommand{\PntSetB}{{{\mathsf{Q}}}}
\newcommand{\PntSetC}{{{\mathsf{R}}}}
\newcommand{\Growth}[2]{{\EuScript{G}}_{#1}\pth{#2}}
\newcommand{\ShatterFunc}[1]{\pi_{#1}}
\newcommand{\Distribution}{\EuScript{D}}
\newcommand{\Error}{\EuScript{E}}

\newcommand{\bd}{{\partial}}%
\newcommand{\eps}{{\varepsilon}}%

\newcommand{\clmlab}[1]{\label{claim:#1}}
\newcommand{\clmref}[1]{\HLink{Claim}{claim:#1}}

\newcommand{\CHX}[1]{{\mathcal{CH}}\pth{#1}}

\newcommand{\ts}{\hspace{0.6pt}}

\def\qedsymbol{\myqedsymbol}%

\newcommand{\qedsign}{\hfill{\hfill\myqedsymbol}}
\newcommand{\qedsignBig}{\hfill{\hfill\rule{4mm}{4mm}}}
\newcommand{\QED}{{\hfill{\hfill\qedsymbol}}}

\newcommand{\num}{\alpha}

\newcommand{\NEW}[1]{#1}

\StartDefiningTabulars
\newcommand{\ThreeFiguresExt}[9]{%
   \begin{tabular}{ccc}
     \begin{minipage}{0.31\linewidth}
         \begin{center}
             {\IncludeGraphics[#1]{#2}}
         \end{center}
     \end{minipage}
     &
       \begin{minipage}{0.31\linewidth}
           \centerline{\IncludeGraphics[#4]{#5}}
       \end{minipage}
     &
       \begin{minipage}{0.31\linewidth}
           \centerline{\IncludeGraphics[#7]{#8}}
       \end{minipage}
     \\
     ~\\[-0.2cm]
     {#3} & {#6} & {#9}
   \end{tabular}}
\StopDefiningTabulars

\newcommand{\IncludeGraphics}[2][]{%
   \typeout{}%
   \typeout{Graphics: #2}%
   \typeout{\ includegraphics[#1]{#2}}%
   \includegraphics[#1]{#2}
   \typeout{}%
}

%    -----------------------------------------------------------------
%    -----------------------------------------------------------------
%    -----------------------------------------------------------------
%    Handling references
%    -----------------------------------------------------------------
%    -----------------------------------------------------------------
%    -----------------------------------------------------------------
\newcommand{\HLinkShort}[2]{\hyperref[#2]{#1\ref*{#2}}}
\newcommand{\HLink}[2]{\hyperref[#2]{#1~\ref*{#2}}}
\newcommand{\HLinkPage}[2]{\hyperref[#2]{#1~\ref*{#2}%
      $_\text{p\pageref{#2}}$}}

\newcommand{\HLinkSuffix}[3]{\hyperref[#2]{#1\ref*{#2}{#3}}}
\newcommand{\HLinkPageSuffix}[3]{\hyperref[#2]{#1\ref*{#2}%
      #3$_\text{p\pageref{#2}}$}}

\newcommand{\figlab}[1]{\label{fig:#1}}
\newcommand{\figref}[1]{\HLink{Figure}{fig:#1}}

\newcommand{\exelab}[1]{\label{exer:#1}}
\newcommand{\exeref}[1]{Exercise~\ref{exer:#1}}

\providecommand{\deflab}[1]{\label{def:#1}}
\newcommand{\defref}[1]{\HLink{Definition}{def:#1}}
\newcommand{\defrefY}[2]{\hyperref[def:#2]{#1}}
\newcommand{\defrefpage}[1]{\HLinkPage{Definition}{def:#1}}

\newcommand{\etal}{\textit{et~al.}\xspace}

\newcommand{\obslab}[1]{\label{observation:#1}}
\newcommand{\obsref}[1]{\HLink{Observation}{observation:#1}}
\newcommand{\obsrefY}[2]{\hyperref[observation:#2]{#1}}

\newcommand{\exmlab}[1]{\label{example:#1}}
\newcommand{\exmref}[1]{\HLink{Example}{example:#1}}

\newcommand{\remlab}[1]{\label{rem:#1}}
\newcommand{\remref}[1]{\HLink{Remark}{rem:#1}}%

\newcommand{\chaplab}[1]{\label{chap:#1}}

\newcommand{\itemlab}[1]{\label{item:#1}}
\newcommand{\itemref}[1]{\HLinkSuffix{}{item:#1}{}}

\newcommand{\tedlab}[1]{\label{tedium:#1}}
\newcommand{\tedref}[1]{\HLink{Tedium}{tedium:#1}}%

\newcommand{\factlab}[1]{\label{fact:#1}}%

\newcommand{\alglab}[1]{\label{algorithm:#1}}

\newcommand{\net}{\mathcal{N}}%
      
\newcommand{\lemlab}[1]{\label{lemma:#1}}
\newcommand{\lemref}[1]{\HLink{Lemma}{lemma:#1}}%
\newcommand{\lemrefY}[2]{\hyperref[lemma:#2]{#1}}
\newcommand{\lemrefpage}[1]{\HLinkPage{Lemma}{lemma:#1}}
\newcommand{\lemrefshort}[1]{\HLinkShort{L}{lemma:#1}}

\newcommand{\eqlab}[1]{\label{equation:#1}}
\newcommand{\Eqrefbare}[1]{\HLinkSuffix{(}{equation:#1}{)}}
\newcommand{\Eqref}[1]{\HLinkSuffix{Eq.~(}{equation:#1}{)}}
\newcommand{\Eqrefpage}[1]{\HLinkPageSuffix{Eq.~(}{equation:#1}{)}}

\newcommand{\seclab}[1]{\label{sec:#1}}
\newcommand{\secref}[1]{\HLink{Section}{sec:#1}}
\newcommand{\secrefY}[2]{\hyperref[sec:#2]{#1}}
\newcommand{\secrefpage}[1]{\HLinkPage{Section}{sec:#1}}

\newcommand{\thmlab}[1]{{\label{theo:#1}}}
\newcommand{\thmref}[1]{\HLink{Theorem}{theo:#1}}
\newcommand{\thmrefY}[2]{\hyperref[theo:#2]{#1}}
\newcommand{\thmrefpage}[1]{\HLinkPage{Theorem}{theo:#1}}

\newcommand{\corlab}[1]{\label{cor:#1}}
\newcommand{\corref}[1]{\HLink{Corollary}{cor:#1}}%

%    -----------------------------------------------------------------
%    -----------------------------------------------------------------
%    -----------------------------------------------------------------
%    Handling references
%    -----------------------------------------------------------------
%    -----------------------------------------------------------------
%    -----------------------------------------------------------------

\newcommand{\ceil}[1]{\left\lceil {#1} \right\rceil}
\newcommand{\floor}[1]{\left\lfloor {#1} \right\rfloor}
\newcommand{\brc}[1]{\left\{ {#1} \right\}}
\newcommand{\pth}[2][\!]{\mleft({#2}\mright)}%
\newcommand{\brk}[1]{\mleft[{#1}\mright]}%
\newcommand{\pbrc}[2][\!\!]{#1\left[ {#2} \bigr. \right]}
\newcommand{\pbrcS}[1]{\left[ {#1} \right]}

\newcommand{\dotprod}[2]{{#1}\cdot{#2}}
\newcommand{\DotProd}[2]{\permut{{#1},{#2}}}
\newcommand{\permut}[1]{\left\langle {#1} \right\rangle}

\newcommand{\MakeBig}{}
\newcommand{\MakesBig}{}
\newcommand{\MakeSBig}{}
\newcommand{\MakeVBig}{}

\newcommand{\DistChar}{\mathsf{d}}%
\newcommand{\DistSetY}[2]{\DistChar\pth{#1,#2}}
\newcommand{\distG}[2]{\DistChar_\Graph\pth{#1, #2}}
\newcommand{\DistMtrCharX}[1]{\DistChar_{#1}}

\newcommand{\dY}[2]{\DistChar\pth{#1, #2}}%
\newcommand{\dGY}[2]{\DistChar_\Graph\pth{#1, #2}}%

\newcommand{\cset}{\mathcal{C}}%

\newcommand{\MTR}{\mathcal{M}}
\newcommand{\DistMtr}[2]{\DistChar\pth{#1,#2}}

\newcommand{\dist}[1]{\left\| {#1} \right\|}
\newcommand{\norm}[1]{\left\| {#1} \right\|}

\providecommand{\ds}{\displaystyle}
\newcommand{\sep}[1]{\,\left|\, {#1} \bigr.\right.}%
\newcommand{\sepw}[1]{\left|\, {#1} \right.}%

\newcommand{\ProbExt}[2]{\mathop{\mathbf{Pr}}_{#1}\pbrcx{#2}}
\newcommand{\ProbLTR}{\mathbb{P}}%

\newcommand{\Prob}[1]{\mathop{\ProbLTR} \mleft[ #1 \mright]}%
\newcommand{\ProbCond}[2]{\mathop{\ProbLTR}\!\left[%
       #1 \;\middle\vert\; #2 \right]}
\newcommand{\ProbY}[2]{\ProbLTR^{}_{\!#1}\pbrc{#2}}

\newcommand{\ExChar}{\mathbb{E}}%
\newcommand{\ExSym}{\mathop{\ExChar}}%
\newcommand{\Ex}[2][\!]{\ExSym#1\pbrcx{#2}}
\newcommand{\ExExt}[2]{\ExSym_{#1}\!\pbrcx{#2}}
\newcommand{\ExExtS}[2]{\ExSym_{\substack{#1}}\!\pbrcx{#2}}
\newcommand{\ExCond}[2]{\ExSym\!\left[%
       #1 \;\middle\vert\; #2 \right]}

\newcommand{\VarSymb}{\mathop{\mathbb{V}}}
\newcommand{\Var}[1]{\mathop{\mathbb{V}}\!\pbrcx{#1}}

\newcommand{\pbrcx}[1]{\left[ {#1} \right]}%
\newcommand{\nfrac}[2]{{#1}/{#2}}
\newcommand{\cardin}[1]{\left| {#1} \right|}%

\newcommand\CCancel[2][black]{{\renewcommand\CancelColor{\color{#1}}\cancel{#2}}}

\newcommand{\IncGraphPage}[4][]{%
   \IncludeGraphics[page=#4,#1]{\File{#2/#3}}%
}

\newenvironment{enumEx}{
   \begin{compactenumA}
    }{\end{compactenumA}}

\newcommand{\cl}{{\mathrm{cl}}}

\usepackage[amsmath,thmmarks]{ntheorem}
\theoremseparator{.}%

\newtheorem{theorem:unnumbered}{Theorem}[section]
\newtheorem{theorem}{Theorem}[section]

\newtheorem{observation}[theorem]{Observation}
\newtheorem{conjecture}[theorem]{Conjecture}
\newtheorem{lemma}[theorem]{Lemma}
\newtheorem{corollary}[theorem]{Corollary}
\newtheorem{claim}[theorem]{Claim}
\newtheorem{assumption}[theorem]{Assumption}
\newtheorem{tedium}[theorem]{Tedium}

\theoremstyle{plain}%
\theoremheaderfont{\sffamily}
\theorembodyfont{\upshape}%

\newtheorem{defn}[theorem]{Definition}
\newtheorem{definition}[theorem]{Definition}

\newtheorem*{remarkunnumbered}{Remark}%
\newenvironment{remark_unnumbered}{\begin{remarkunnumbered}}{\end{remarkunnumbered}}
\newtheorem*{remarks}{Remarks}%
\newtheorem{remark}[theorem]{Remark}%
\newtheorem{fact}[theorem]{Fact}%
\newtheorem{example}[theorem]{Example}
\newtheorem{exercise}[theorem]{Exercise}
\newtheorem{problem}[theorem]{Problem}
\newtheorem{algorithm}[theorem]{Algorithm}%
\newtheorem{xca}[theorem]{Exercise}
\newtheorem{openproblem}[theorem]{Open Problem}
\newtheorem{exerciseh}[theorem]{Exercise}
\newenvironment{exercise_h}[1][]{\begin{exerciseh}[#1]}{\end{exerciseh}}

% Proof environments
\theoremheaderfont{\em}%
\theorembodyfont{\upshape}%
\theoremstyle{nonumberplain}
\theoremseparator{}
\theoremsymbol{\myqedsymbol}
\newtheorem{proof}{Proof:}
\newtheorem{proofextin}{Proof}
\newtheorem{altproof}{Alternative proof:}

\theoremsymbol{}
\newtheorem{proofnoqed}{Proof:}

\theoremsymbol{\myqedsymbol}
\theorempreskip{0pt}%
\theorempostskip{0pt}%
\newtheorem{proofof}{Proof of\!}%

\newenvironment{proofExt}[1]{{\emph{#1:}}}{\hfill{\hfill\myqedsymbol}}
\newenvironment{proofnoproof}{{}}{{\hfill{\hfill\myqedsymbol}}}

%%%%%%%%%%%%%%%%%%%%%%%%%%%%%%%%%%%%%%%%%%%%%%%%%%%%%%%%%%%%%%%%%%%%
%%%%%%%%%%%%%%%%%%%%%%%%%%%%%%%%%%%%%%%%%%%%%%%%%%%%%%%%%%%%%%%%%%%%
% Macros to embed figures in text

\newcommand{\EBRY}[2]{#1_{\langle#2\rangle}}%

\newcommand{\NiceFrac}[2]{%
   {\raisebox{0.2em}{\ensuremath{#1}}\!\!\!}%
   {\rotatebox[origin=c]{-20}{\Big/}}% Rotated /
   {\! \raisebox{-0.2em}{\ensuremath{#2}} }}%

% end - Macros to embed figures in text
%%%%%%%%%%%%%%%%%%%%%%%%%%%%%%%%%%%%%%%%%%%%%%%%%%%%%%%%%%%%%%%%%%%% 
%%%%%%%%%%%%%%%%%%%%%%%%%%%%%%%%%%%%%%%%%%%%%%%%%%%%%%%%%%%%%%%%%%%%

%%%%%%%%%%%%%%%%%%%%%%%%%%%%%%%%%%%%%%%%%%%%%%%%%%%%%%%%%%%%%%%%%%
% Color in math
\providecommand{\Mh}[1]{{#1}}%
\providecommand{\IfColorMath}[2]{#2}%

\IfFileExists{.latex_printer_friendly}{\def\GenPrinterVer{1}}{}%

\ifx\GenPrinterVer\undefined
   \IfFileExists{.latex_color}{\def\GenColorMath{1}}{}
\else
   \renewcommand{\IfPrinterVer}[2]{#1}%
\fi

\ifx\GenColorMath\undefined
\else
\renewcommand{\IfColorMath}[2]{#1}%
\renewcommand{\Mh}[1]{{\textcolor{red}{#1}}}%
\fi
% Color in math - end
%%%%%%%%%%%%%%%%%%%%%%%%%%%%%%%%%%%%%%%%%%%%%%%%%%%%%%%%%%%%%%%%%%

\newcommand{\HSetY}[2]{H\pth{#1, #2}}%
\newcommand{\CHSetY}[2]{C\pth{#1, #2}}%
\DeclareRobustCommand{\Holder}{Hölder\xspace}
\newcommand{\OldOm}{{\mathcal{O}}}
\newcommand{\NewOm}{{\mathcal{N}}}

%%%%%%%%%%%%%%%%%%%%%%%%%%%%%%%%%%%%%%%%%%%%%%%%%%%%%%%%%%%%%%%%%%%%%%%
%---------------------------------------------------------------
\newcommand{\IntRange}[1]{\left\llbracket #1 \right\rrbracket}
\newcommand{\IRX}[1]{\left\llbracket #1 \right\rrbracket}%
\newcommand{\IRY}[2]{\left\llbracket #1:#2 \right\rrbracket}

%---------------------------------------------------------------
%%%%%%%%%%%%%%%%%%%%%%%%%%%%%%%%%%%%%%%%%%%%%%%%%%%%%%%%%%%%%%%%%%%%%%%

\newcommand{\obj}{f}%
\newcommand{\Wg}{\omega}%
\newcommand{\TotalW}{W}%

\newcommand{\rankY}[2]{{#1}_{ \permut{#2}}}%%
\newcommand{\areaX}[1]{\mathrm{area}\pth{#1}}%%

\newcommand{\mybinom}[2]%
{\Bigl(\begin{array}{@{}c@{}}#1\\#2\end{array}\Bigr)}

\providecommand{\Set}[2]{\left\{ #1 \;\middle\vert\; #2 \right\}}
\newcommand*{\biga}[1]{\vcenter{\hbox{\scalebox{1.5}{\ensuremath#1}}}}
\newcommand{\CNF}{\TermI{CNF}\xspace}%
\newcommand{\Tree}{T}%

\DeclareSymbolFont{sfletters}{OML}{cmbrm}{m}{it}

\DeclareMathSymbol{\salpha}{\mathord}{sfletters}{"0B}

%%%%%%%%%%%%%%%%%%%%%%%%%%%%%%%%%%%%%%%%%%%%%%%%%%%%%%%%%%%%%%%%%%%%

\providecommand{\ProblemInd}[2]{\textsf{\textcolor{red}{#1}\index{Problem!#2}}}
\newcommand{\NPProblemInd}[4]{
   \begin{center}
      \vspace{0.6cm}
       \fbox{
          \begin{minipage}{0.9\linewidth}
              \vspace{-0.6cm} \hspace{-0.35cm}
              {\large \bfseries \ProblemInd{#1}{#2}}
              \vspace{+0.2cm}%

              \noindent{\textcolor{blue}{\bfseries \textsf{Instance}}:} {#3}
              
              \noindent{\textcolor{blue}{\bfseries \textsf{Question}:}} {#4}
          \end{minipage}
       }
   \end{center}
   \smallskip
   }

\newcommand{\NPProblemOpt}[4]{\NPProblemInd{#1}{#2}{#3}{#4}}

%%%%%%%%%%%%%%%%%%%%%%%%%%%%%%%%%%%%%%%%%%%%%%%%%%%%%%%%%%%%%%%%%%%%

\newcommand{\Hastad}{H{\aa}stad\xspace}

\newcommand{\ANNbr}{\Term{ANN}\index{approximate nearest neighbor}\xspace}
\newcommand{\LSH}{\TermI{L{S}H}\xspace}

\newcommand{\DS}{\EuScript{D}} % DataStructure
\newcommand{\DSC}{\mathsf{D}} % DataStructure

\newcommand{\distH}[2]{\DistChar_{H}\pth{#1, #2}}
\newcommand{\distHY}[2]{\DistChar_{H}\pth{#1, #2}}

\newcommand{\DSAprxNearNbr}{\ensuremath{\DataStructure{{\approx%
         \text{Near{}}}}}%
   \index{near neighbor!data-structure!approximate}%
}
\newcommand{\DataStructure}[1]{\ensuremath{\DS_{#1}}}
\newcommand{\DistNN}[2]{\DistChar\pth{#1,#2}}
\newcommand{\pntc}{{p}}%

\newcommand{\seq}{{M}}%
\newcommand{\cTimes}{\beta}%

\newcommand{\DSTimes}{L}%
\newcommand{\rad}{r}%

\newcommand{\HypercubeX}[1]{\Hypercube^{#1}}
\newcommand{\kFunctions}[2]{\mathrm{combine}\pth{#1,#2}}
\newcommand{\WEqualChar}{\Bumpeq}
\newcommand{\WEqual}[2]{{#1}\WEqualChar{#2}}
\newcommand{\FamilyAmp}{\FamilyC_{\WEqualChar}}
\newcommand{\PSPACE}{\ComplexityClass{PSPACE}\xspace}%

\newcommand{\PS}{P}%

\newcommand{\PSetFar}{\ensuremath{\PntSet^{}_{\geq}\xspace}}%
\newcommand{\DA}{D}%
\newcommand{\uniqX}[1]{\mathrm{uniq}\pth{#1}}%
\newcommand{\diffX}[1]{\mathop{}\!\mathrm{d}#1}

\newcommand{\Frq}{F}%
\newcommand{\frq}{f}%
\newcommand{\STRM}{\mathcal{S}}%
\newcommand{\strm}{{s}}%

\DeclareRobustCommand{\Turan}{Turán\xspace}
\newcommand{\fGrowX}[1]{G^{}_{\Dim}\pth{#1}}
\newcommand{\CompressFunc}{\mathtt{Compress}}
\newcommand{\RSA}{\Term{RSA}\xspace}
\newcommand{\MST}{\Term{MST}\xspace}
\providecommand{\CClassPP}{\CClass{P{}P}\xspace}
\providecommand{\ID}{\mathrm{id}}

\newcommand{\EncodeFunc}{\text{\ttfamily{Enc}}}
\newcommand{\DecodeFunc}{\text{\ttfamily{Dec}}}
\newcommand{\discX}[1]{\mathrm{disc}\pth{#1}}%

\newcommand{\propertyC}{\psi}%
\newcommand{\propertyX}[1]{\propertyC\pth{#1}}%
\newcommand{\numC}{\Gamma}%

\newcommand{\degreeX}[1]{\mathrm{degree}\pth{#1}}

\newcommand{\ind}{m}

\providecommand{\nequiv}{\not\equiv}
\providecommand{\sqcupplus}{\sqcup^+}

\newcommand{\AlgorithmI}[1]{{\protect\color{RedViolet}{#1}}}

\usepackage{marginnote}
\usepackage{ragged2e} 
% Short, subtle snark macro
\setlength{\marginparwidth}{1.5in}
\setlength{\marginparsep}{10pt}
\newcommand{\snark}[1]{%
   \marginnote{\raggedright\scriptsize\sffamily\color{gray}#1}%
}


\MakePerPage{footnote} % Resets footnote counter on every page

\allowdisplaybreaks % Allows page breaks inside multiline math environments

%-----------------------------------------------------------------------------
% Bibliography & Document Begin
%-----------------------------------------------------------------------------
\IfFileExists{\jobname.bib}{%
   \addbibresource{\jobname.bib}%
}{%
   \addbibresource{shortcuts.bib}%
   \addbibresource{geometry.bib}%
}
\makeindex
\begin{document}
}{\input{prefix_lwarp}}%
\else
  \ifdefined\HCode
    \IfFileExists{styles/prefix_4ht.tex}{%-----------------------------------------------------------------------------
% TeX4ht Prefix File for Randomized Algorithms Class Notes
% Kept separate from prefix_latex.tex so HTML-specific simplifications cannot
% affect the standard PDF or lwarp builds.
%-----------------------------------------------------------------------------

\RequirePackage{iftex}
\ifPDFTeX
  \PackageError{engine-guard}{This document requires LuaLaTeX or XeLaTeX}{%
    Please recompile using LuaLaTeX or XeLaTeX instead of pdfLaTeX.%
  }
  \stop
\fi

\ifdefined\HCode\else
  \PackageError{tex4ht-gate}{%
    Attempting to use prefix_4ht.tex outside TeX4ht.%
  }{Compile through make4ht, or use prefix_latex.tex for PDF generation.}
  \stop
\fi

% Always use memoir
\documentclass[12pt,oneside,fleqn]{memoir}

%-----------------------------------------------------------------------------
% Typography, Fonts & Math Packages
%-----------------------------------------------------------------------------
\usepackage{mathtools}
\PassOptionsToPackage{warnings-off={mathtools-overbracket,mathtools-colon}}{unicode-math}
\usepackage{fontspec}
\usepackage{unicode-math}

% New Computer Modern OpenType text and math fonts
\setmainfont{NewCM10-Book.otf}[
  BoldFont          = NewCM10-Bold.otf,
  ItalicFont        = NewCM10-BookItalic.otf,
  BoldItalicFont    = NewCM10-BoldItalic.otf,
  SlantedFont       = NewCM10-Book.otf,
  SlantedFeatures   = {FakeSlant = 0.25},
  BoldSlantedFont   = NewCM10-Bold.otf,
  BoldSlantedFeatures = {FakeSlant = 0.25},
  SmallCapsFeatures = {Letters = SmallCaps}
]
\setsansfont{NewCMSans10-Book.otf}[
  BoldFont          = NewCMSans10-Bold.otf,
  ItalicFont        = NewCMSans10-BookOblique.otf,
  BoldItalicFont    = NewCMSans10-BoldOblique.otf
]
\setmonofont{NewCMMono10-Book.otf}[
  BoldFont          = NewCMMono10-Bold.otf,
  ItalicFont        = NewCMMono10-BookItalic.otf,
  BoldItalicFont    = NewCMMono10-BoldOblique.otf
]

\setmathfont{NewCMMath-Book.otf}
\setmathfont{NewCMMath-Bold.otf}[version=bold]

% OpenType math aliases and fixes
\providecommand{\square}{\mdlgwhtsquare}
\def\Box{\mdlgwhtsquare}
\def\Diamond{\mdlgwhtdiamond}
\def\blacksquare{\mdlgblksquare}
\def\leadsto{\rightsquigarrow}
\providecommand{\llbracket}{\lBrack}
\providecommand{\rrbracket}{\rBrack}
\DeclareTextCommand{\textcircled}{TU}[1]{%
  \mbox{\ooalign{\hfil\raise0.1ex\hbox{#1}\hfil\crcr\char"25EF}}%
}

\usepackage{alphabeta}
\usepackage[normalem]{ulem} % Provides \sout
\usepackage{mleftright}
\usepackage{microtype}
\usepackage{xspace}
\usepackage{xurl}
\usepackage[dvipsnames,table]{xcolor}
\usepackage{graphicx}
\usepackage{caption}
\usepackage{booktabs,array,colortbl,hhline,multirow,tabularx}
\usepackage[inline]{enumitem}
\usepackage[thicklines]{cancel}
\renewcommand\CancelColor{\color{purple}}
\usepackage{euscript}
\usepackage{marvosym}
\let\Cross\undefined
\let\Square\undefined
\usepackage{bbding}
\usepackage{epigraph}
\usepackage{makeidx}
\usepackage{perpage}
\usepackage{tikz}

%-----------------------------------------------------------------------------
% Boondox / mathcalb font declaration
%-----------------------------------------------------------------------------
\DeclareFontFamily{U}{BOONDOX-calo}{\skewchar\font=45 }
\DeclareFontShape{U}{BOONDOX-calo}{m}{n}{
  <-> s*[1.05] BOONDOX-r-calo}{}
\DeclareFontShape{U}{BOONDOX-calo}{b}{n}{
  <-> s*[1.05] BOONDOX-b-calo}{}
\DeclareMathAlphabet{\mathcalb}{U}{BOONDOX-calo}{m}{n}
\SetMathAlphabet{\mathcalb}{bold}{U}{BOONDOX-calo}{b}{n}
\DeclareMathAlphabet{\mathbcalb}{U}{BOONDOX-calo}{b}{n}

%-----------------------------------------------------------------------------
% Algorithm Environment (formerly salgorithm.sty)
%-----------------------------------------------------------------------------
\newcommand{\ProgWordX}[1]{{\small\textbf{#1}}{\ }}
\newcommand{\Do}{\ProgWordX{do}\xspace}
\newcommand{\End}{\ProgWordX{end}\xspace}
\newcommand{\Return}{\ProgWordX{return}\xspace}%
\newcommand{\Procbegin}{\ProgWordX{begin}}
\newcommand{\If}{\ProgWordX{if}\xspace}
\newcommand{\Then}{\ProgWordX{then}\xspace}
\newcommand{\Else}{\ProgWordX{else}\xspace}
\newcommand{\For}{\ProgWordX{for}\xspace}
\newcommand{\While}{{\small\bfseries while}\ }

\newcommand{\xbeginlgox}{\begin{minipage}{1in}\begin{tabbing}
           \quad\=\qquad\=\qquad\=\qquad\=\qquad\=\qquad\=\qquad\=\kill}
\newcommand{\xendlgox}{\end{tabbing}\end{minipage}}

\newenvironment{algorithmEnv}{%
   \begin{tabular}{|l|}\hline\xbeginlgox\bigskip}%
   {\xendlgox\\\hline\end{tabular}}

\newenvironment{program}{%
   \begin{minipage}{4.0in}%
   \begin{tabbing}%
       \ \ \ \ \=\ \ \ \=\ \ \ \ \=\ \ \ \ \=\ \ \ \ \=%
      \ \ \ \ \=\ \ \ \ \=\ \ \ \ \=\ \ \ \ \=%
      \ \ \ \ \=\ \ \ \ \=\ \ \ \ \=\ \ \ \ \=\kill%
}{%
   \end{tabbing}%
   \end{minipage}%
}

\newcommand{\AlgorithmBox}[1]{%
   \fbox{\begin{program} #1 \end{program}}%
}

\definecolor{AlgorithmColor}{cmyk}{0.07,0.90,0,0.34}
\providecommand{\Algorithm}[1]{{\AlgorithmI{#1}\index{algorithm!#1@{\AlgorithmI{#1}}}}}
\newcommand{\CodeComment}[1]{\textcolor{blue}{\texttt{#1}}}

%-----------------------------------------------------------------------------
% Path & Directory Resolution
%-----------------------------------------------------------------------------
\newcommand{\RootDIR}{../}
\IfFileExists{../../../book.tex}{\renewcommand{\RootDIR}{../../../}}{}
\IfFileExists{../../book.tex}{\renewcommand{\RootDIR}{../../}}{}
\IfFileExists{../book.tex}{\renewcommand{\RootDIR}{../}}{}
\IfFileExists{./book.tex}{\renewcommand{\RootDIR}{}}{}

\providecommand{\File}[1]{#1}
\providecommand{\RootFile}[1]{\RootDIR #1}

%-----------------------------------------------------------------------------
% QED Duck Symbol
%-----------------------------------------------------------------------------
\IfFileExists{styles/qed_duck.svg}{%
   \providecommand{\myqedsymbol}{%
      \HCode{<img class="qed-duck" src="styles/qed_duck.svg" alt="QED" />}%
   }%
}{%
   \providecommand{\myqedsymbol}{\HCode{<span class="qed-fallback">&#x25A0;</span>}}%
}

%-----------------------------------------------------------------------------
% BibLaTeX Setup (Pure Biber, No BibTeX)
%-----------------------------------------------------------------------------
\usepackage[bibencoding=utf8,style=alphabetic,refsegment=chapter,
indexing=cite,
backend=biber]{biblatex}

\ExecuteBibliographyOptions{%
   maxnames = 40,%
   maxalphanames = 4,%
   maxcitenames = 3,%
   minalphanames = 3,%
   minnames = 1,%
   giveninits = true,%
   doi = false,%
   url = false,%
   isbn = false%
}

\renewbibmacro{in:}{}

\AtEveryBibitem{%
   \ifentrytype{inproceedings}{%
      \clearname{editor}%
      \clearlist{editor}%
      \clearfield{editor}%
      \clearlist{publisher}%
      \clearfield{series}%
      \clearlist{location}%
      \clearfield{publisher}%
   }{}%
   \ifentrytype{book}{%
      \clearname{publisher}%
   }{}%
}
\AtEveryCitekey{%
   \ifentrytype{inproceedings}{%
      \clearlist{publisher}%
   }{}%
   \clearname{editor}%
   \clearlist{location}%
}
\DeclareFieldFormat[inproceedings]{date}{}
\DeclareFieldFormat[inproceedings]{pages}{#1, \printfield{year}}
\DeclareFieldFormat[inproceedings]{booktitle}{\mkbibemph{#1,\nopunct}}

\DeclareFieldFormat{journaltitle}{\mkbibemph{#1}\isdot}

\DeclareFieldFormat[article]{number}{\mkbibparens{#1}}
\DeclareFieldFormat[article]{date}{}
\DeclareFieldFormat[article]{pages}{#1, \printfield{year}}
\DeclareFieldFormat[article]{issue}{\nopunct#1}
\renewbibmacro*{issue+date}{%
   \iffieldundef{issue}{}%
   {%
      \printtext{\unspace}%
      (\printfield{issue})%
   }%
   \iffieldundef{pages}%
   {\addcomma\printfield{year}}%
   {}%
}

\renewcommand{\bibpagespunct}{%
   \ifentrytype{article}{\addcolon}{\addcomma}\space%
}

\DeclareFieldFormat[article]{titlecase}{\MakeSentenceCase{#1}}
\DeclareFieldFormat[inproceedings]{titlecase}{\MakeSentenceCase{#1}}
\DeclareFieldFormat[incollection]{titlecase}{\MakeSentenceCase{#1}}
\DeclareFieldFormat[book]{titlecase}{\MakeSentenceCase{#1}}

\renewbibmacro*{journal}{%
   \iffieldundef{journaltitle} {} {\printtext[journaltitle]{%
         \printfield[myplain]{journaltitle}%
         \setunit{\subtitlepunct}%
         \printfield[myplain]{journalsubtitle}}}}
\DeclareFieldFormat{myplain}{#1}

\renewbibmacro*{booktitle}{%
   \ifboolexpr{ test {\iffieldundef{booktitle}} and test
      {\iffieldundef{booksubtitle}} } {} {\printtext[booktitle]{%
         \printfield[myplain]{booktitle}%
         \setunit{\subtitlepunct}%
         \printfield[myplain]{booksubtitle}}%
      \newunit}%
   \printfield{booktitleaddon}}

\renewbibmacro*{volume+number+eid}{%
   \setunit*{\addcomma\addspace}%
   \printfield{volume}%
   \printfield{number}%
   \setunit{\addcomma\space}%
   \printfield{eid}%
}

\newbibmacro{string+doi}[1]{%
   \iffieldundef{doi}{%
      \iffieldundef{url}{%
         #1%
      }{%
         \href{\thefield{url}}{#1}%
      }%
   }{%
      \href{http://dx.doi.org/\thefield{doi}}{#1}}%
}
\DeclareFieldFormat*{title}{\usebibmacro{string+doi}{\mkbibemph{#1}}}
\renewcommand*{\multicitedelim}{\addcomma\space}
\DeclareFieldFormat[inproceedings]{title}{\usebibmacro{string+doi}{{#1}}}
\DeclareFieldFormat[article]{title}{\usebibmacro{string+doi}{{#1}}}
\DeclareFieldFormat[incollection]{title}{\usebibmacro{string+doi}{{#1}}}
\DeclareFieldFormat[book]{titlecase}{{#1}}

\DeclareCiteCommand{\tcite}%
{\usebibmacro{prenote}}%
{\usebibmacro{citeindex}\usebibmacro{cite}}%
{\multicitedelim}%
{\usebibmacro{postnote}}

%-----------------------------------------------------------------------------
% Hyperref & Visual Formatting
%-----------------------------------------------------------------------------
\usepackage{hyperref}
\hypersetup{%
   breaklinks,%
   colorlinks=true,%
   linkcolor=[rgb]{0.45,0.0,0.0},%
   citecolor=[rgb]{0,0,0.45}%
}

\definecolor{DeepNavy}{HTML}{002366}
\definecolor{RoyalBlue}{HTML}{004E92}
\definecolor{SteelBlue}{HTML}{4682B4}
\definecolor{TealBlue}{HTML}{008080}
\definecolor{DarkSlate}{HTML}{2F4F4F}
\definecolor{DustyGrey}{HTML}{696969}

\makeatletter
\AtBeginDocument{\let\@elt\relax}
\makeatother

\headstyles{bringhurst}
\chapterstyle{dash}

\renewcommand{\chapnamefont}{\normalfont\huge\bfseries\color{DeepNavy}}
\renewcommand{\chapnumfont}{\normalfont\huge\bfseries\color{DeepNavy}}
\renewcommand{\chaptitlefont}{\normalfont\Huge\bfseries\color{DeepNavy}}

\counterwithin{footnote}{chapter}

\setsubparaheadstyle{\normalfont\normalsize\itshape\color{DustyGrey}}
\setlength{\beforechapskip}{-20pt} 
\setlength{\afterchapskip}{30pt}
\setsecnumformat{\csname the#1\endcsname.\quad}

%-----------------------------------------------------------------------------
% Standalone Chapter vs Master Book Modes
%-----------------------------------------------------------------------------
\newcommand{\Copyright}{This work is licensed under the Creative Commons Attribution-Noncommercial 3.0 License. To view a copy of this license, visit \url{http://creativecommons.org/licenses/by-nc/3.0/} or send a letter to Creative Commons, 171 Second Street, Suite 300, San Francisco, California, 94105, USA.}

\newcommand{\RealChapterEndInner}{%
   {%
      \begingroup%
      \emergencystretch=1em%
      \printbibliography[heading=subbibliography,segment=\therefsegment]{}%
      \endgroup%
   }%
}

\providecommand{\BookMode}[1]{}
\providecommand{\StandAloneMode}[1]{#1}

\ifx\bookMode\undefined
   % STANDALONE CHAPTER MODE
   \renewcommand{\StandAloneMode}[1]{#1}
   \renewcommand{\BookMode}[1]{}

   \newcommand{\ChapterNumPage}[2]{%
      \if\relax\detokenize{#1}\relax
         \setcounter{chapter}{0}%
      \else\ifnum#1>0
         \setcounter{chapter}{\numexpr#1-1\relax}%
      \else
         \setcounter{chapter}{0}%
      \fi\fi
      \if\relax\detokenize{#2}\relax
         \setcounter{page}{1}%
      \else\ifnum#2>0
         \setcounter{page}{#2}%
      \else
         \setcounter{page}{1}%
      \fi\fi
   }

   \DeclareRobustCommand{\setheadingcolor}{\color{RoyalBlue}}
   \setsecheadstyle{\setheadingcolor\normalfont\Large\bfseries}
   \setsubsecheadstyle{\normalfont\large\bfseries\color{SteelBlue}}
   \setsubsubsecheadstyle{\normalfont\normalsize\bfseries\color{TealBlue}}
   \setparaheadstyle{\normalfont\normalsize\bfseries\color{DarkSlate}}

   \DeclareDocumentCommand{\Chapter}{s o m}{%
      \IfBooleanTF{#1}{%
         \IfValueTF{#2}{\chapter*[]{#2}}{\chapter*[]{#3}}%
      }{%
         \IfValueTF{#2}{\chapter*[]{#2}}{\chapter[]{#3}}%
      }%
      \setcounter{footnote}{0}%
   }
   \newcommand{\ChapterPrefix}{}
   \newcommand{\ChapterGen}{\pagestyle{plain}\vspace*{-0.7cm}}
   \newcommand{\ChapterEnd}{%
      \RealChapterEndInner{}%
      \end{document}%
   }

   \renewcommand{\File}[1]{#1}
   \renewcommand{\RootFile}[1]{../#1}
\else
   % MASTER BOOK MODE
   \renewcommand{\StandAloneMode}[1]{}
   \renewcommand{\BookMode}[1]{#1}

   \newcommand{\ChapterNumPage}[2]{}

   \DeclareDocumentCommand{\Chapter}{s o m}{%
      \IfBooleanTF{#1}{%
         \IfValueTF{#2}{\chapter*[#2]{#3}}{\chapter*{#3}}%
         \ifx\Dir\undefined\else
            \typeout{CHAPTER_MAP: \Dir/\CurrentChapFile => chap: 0, page: \thepage}%
         \fi
      }{%
         \IfValueTF{#2}{\chapter[#2]{#3}}{\chapter{#3}}%
         \ifx\Dir\undefined\else
            \typeout{CHAPTER_MAP: \Dir/\CurrentChapFile => chap: \thechapter, page: \thepage}%
         \fi
      }%
   }
   \newcommand{\ChapterPrefix}{%
      \vspace{-1cm}%
      \subsection*{\small
         {CS 574: Randomized Algorithms}\\
         {Sariel Har-Peled}\\
         {\today}}%
   } \newcommand{\ChapterGen}{}
   \newcommand{\ChapterEnd}{\RealChapterEndInner{}}

   \renewcommand{\File}[1]{\Dir/#1}
   \renewcommand{\RootFile}[1]{#1}

   \newcommand{\InputChap}[2]{%
      \def\Dir{#1}%
      \def\CurrentChapFile{#2}%
      \typeout{}%
      \typeout{------------------------------------------------------------}%
      \typeout{PROCESSING FILE: #1/#2}%
      \typeout{}%
      \input{#1/#2}%
   }
\fi

%-----------------------------------------------------------------------------
% Fragments & Quotes
%-----------------------------------------------------------------------------
% The PDF build writes fragment bodies to shared files and later includes them
% in standalone appendices.  A chapter webpage already contains each fragment
% at its original location, so TeX4ht only needs a transparent wrapper and can
% omit the repeated standalone copies.
\DeclareDocumentEnvironment{fragment}{m +b}{%
   #2%
}{}
\DeclareDocumentCommand{\IncFragment}{m}{}

\newcommand{\QuoteOpen}{}
\newcommand{\QuoteClose}{}
\newcommand{\QuoteWrite}[2]{}

\newcommand{\Quote}[2]{%
   \SaveIndent
   \QuoteWrite{#1}{#2}%
   \begin{center}%
      {\footnotesize
         \begin{quote}
             \RestoreIndent
             {#1}\\
             \hspace*{2cm}-- {#2.}
         \end{quote}%
      }%
      \index{quotation!#2}%
   \end{center}%
   \bigskip
}

\newcommand{\WebChapterQuote}[3]{%
   \ifvmode\IgnorePar\fi\EndP%
   \HCode{<figure class="chapter-quote"><blockquote>}%
      #1%
   \HCode{</blockquote><hr /><figcaption><span class="quote-author">}%
      #2%
   \HCode{</span><span class="quote-work">}%
      #3%
   \HCode{</span></figcaption></figure>}%
   \par%
}

\newcommand{\QuotePExt}[3]{\WebChapterQuote{#1}{#2}{#3}}

\epigraphfontsize{\small\itshape}
\setlength{\epigraphwidth}{0.6\textwidth}
\newcommand{\QuoteExt}[3]{\WebChapterQuote{#1}{#2}{#3}}

\providecommand{\delX}[1]{{\color{red}\sout{#1}}}
\providecommand{\newX}[1]{{\color{blue}#1}}
\providecommand{\chgY}[2]{{\color{red}\sout{#1}}{\color{blue}#2}}
\providecommand{\remX}[1]{{\color{purple}\textbf{[Remark: #1]}}}
\newcommand{\SpellIgnore}[1]{#1}
\newcommand{\EQ}{\mathrel{\;\;=\;\;}}
\newcommand{\LEQ}{\mathrel{\;\;\leq\;\;}}
\newcommand{\LT}{\mathrel{\;\;<\;\;}}
\newcommand{\GT}{\mathrel{\;\;<\;\;}}

\newcommand{\GEQ}{\mathrel{\;\;\geq\;\;}}

%-----------------------------------------------------------------------------
% Save a symbol that we know is going to get redefined.
\def\savesymbol#1{%
  \expandafter\let\expandafter\origsym\expandafter=\csname#1\endcsname
  \expandafter\let\csname orig#1\endcsname=\origsym
  \expandafter\let\csname#1\endcsname=\relax
}

% Restore a previously saved symbol, and rename the current one.
\def\restoresymbol#1#2{%
  \expandafter\let\expandafter\newsym\expandafter=\csname#2\endcsname
  \expandafter\global\expandafter\let\csname#1#2\endcsname=\newsym
  \expandafter\let\expandafter\origsym\expandafter=\csname orig#2\endcsname
  \expandafter\global\expandafter\let\csname#2\endcsname=\origsym
}

\usepackage{pictex}%
\savesymbol{triangle}%
\usepackage{tree}%
\restoresymbol{TreePKG}{triangle}%

% PicTeX's tree macros perform box arithmetic that TeX4ht cannot convert and
% can terminate the HTML document at the first tree.  These lightweight web
% forms preserve the two children and optional subtree label as HTML/CSS.
\renewcommand{\tree}[2]{%
   \HCode{<span class="web-binary-tree"><span class="tree-node"></span><span class="tree-children"><span class="tree-child">}%
   #1%
   \HCode{</span><span class="tree-child">}%
   #2%
   \HCode{</span></span></span>}%
}
\renewcommand{\TreePKGtriangle}[3][]{%
   \HCode{<span class="web-subtree-triangle"><span class="subtree-label">}%
   #1%
   \HCode{</span></span>}%
}

% Definitions absorbed from sbook.sty
\providecommand{\MdPrc}{\nu}

\providecommand{\TwoFiguresBase}[4]{%
   \begin{tabular}{cc}
      \begin{minipage}{0.48\linewidth}
          \begin{center}
              {{{#1}}}
          \end{center}
      \end{minipage}
      &
      \begin{minipage}{0.48\linewidth}
          \begin{center}
              {{#3}}
          \end{center}
      \end{minipage}
      \\
      {#2} & {#4}
   \end{tabular}%
}

\providecommand{\TwoFigures}[6]{%
   \begin{tabular}{cc}
      \begin{minipage}{0.48\linewidth}
          \begin{center}
              {\IncludeGraphics[#1]{#2}}
          \end{center}
      \end{minipage}
      &
      \begin{minipage}{0.48\linewidth}
          \centerline{\IncludeGraphics[#4]{#5}}
      \end{minipage}
      \\
      {#3} & {#6}
   \end{tabular}%
}

\providecommand{\ThreeFigures}[6]{%
   \begin{tabular}{ccc}
        \begin{minipage}{0.30\linewidth}
            \IncludeGraphics[width=5.1cm]{#1}\\
        \end{minipage}
        &
        \begin{minipage}{0.30\linewidth}
            \IncludeGraphics[width=5.1cm]{#2}\\
        \end{minipage}
        &
        \begin{minipage}{0.30\linewidth}
            \IncludeGraphics[width=5.1cm]{#3}\\
        \end{minipage}
        \vspace{-1cm}
        \\
        {#4} & {#5}& {#6}
   \end{tabular}%
}

%-----------------------------------------------------------------------------
\providecommand{\remove}[1]{}
\providecommand{\NotAMSMode}[1]{#1}
\providecommand{\Solution}[1]{}
\providecommand{\Version}[1]{\hfill{Version: #1}}
\providecommand{\Distor}{\mathrm{dist}}
\providecommand{\Hardness}[1]{(h:#1)}
\providecommand{\QuoteOptExt}[3]{\QuoteExt{#1}{#2}{#3}}
\providecommand{\Hypercube}{\EuScript{H}}
\providecommand{\CoverTime}{\EuScript{C}}
\providecommand{\Nbr}{\Gamma}
\providecommand{\CommuteTime}[2]{\mathbf{CT}_{#1#2}}
\providecommand{\Resistance}{\mathbf{R}}
\providecommand{\EffectResist}[2]{\mathbf{R}_{#1#2}}
\providecommand{\NetworkE}{\EuScript{N}}
\providecommand{\trinom}[4]{\binom{#1}{#2 \;\; #3 \;\; #4}}
\providecommand{\Ranges}{\EuScript{R}}
\providecommand{\CClass}[1]{\ensuremath{\mathbf{#1}}}
\providecommand{\coCClass}[1]{\ensuremath{\mathrm{co-}}\CClass{#1}}
\providecommand{\CX}{\EuScript{C}}
\providecommand{\packet}{\mathbf{c}}
\providecommand{\packetA}{\mathbf{d}}
\providecommand{\mopt}{m_{\mathrm{opt}}}
\providecommand{\girth}{\mathop{\mathrm{girth}}}
\providecommand{\CR}{\mathop{\mathrm{cr}}}
\providecommand{\TRUE}{{\textcolor[named]{RedViolet}{\textsf{TRUE}}}\xspace}
\providecommand{\FALSE}{{\textcolor[named]{RedViolet}{\textsf{FALSE}}}\xspace}
\providecommand{\Part}[1]{\part{#1}}
\providecommand{\Section}[1]{\section{#1}}
\providecommand{\SubSection}[1]{\subsection{#1}}
\providecommand{\SubSubSection}[1]{\subsubsection{#1}}
\providecommand{\Exercise}[5]{%
\begin{exercise_h}[#1]%
   \exelab{#3}{}%
   #4%
\end{exercise_h}%
}
\providecommand{\si}[1]{#1}
\providecommand{\myParagraph}[1]{\noindent\textbf{#1}}
\providecommand{\DefName}[1]{#1}
\providecommand{\deflabel}[1]{#1}
\providecommand{\Description}[1]{}

%-----------------------------------------------------------------------------
% Mathematical Notations & Environments (from notations.sty)
%-----------------------------------------------------------------------------
%-----------------------------------------------------------------------------
% List environments (questions & compactenum)
%-----------------------------------------------------------------------------

\newlist{compactenumA}{enumerate}{5}%
\setlist[compactenumA]{topsep=0pt,itemsep=-1ex,partopsep=1ex,parsep=1ex,%
   label=(\Alph*)}%

\newlist{compactenumI}{enumerate}{5}%
\setlist[compactenumI]{topsep=0pt,itemsep=-1ex,partopsep=1ex,parsep=1ex,%
   label=(\Roman*)}%

\newlist{compactenuma}{enumerate}{5}%
\setlist[compactenuma]{topsep=0pt,itemsep=-1ex,partopsep=1ex,parsep=1ex,%
   label=(\alph*)}%

\newlist{compactitem}{enumerate}{5}%
\setlist[compactitem]{topsep=0pt,itemsep=-1ex,partopsep=1ex,parsep=1ex,%
   label=\textbullet}%

\newlist{compactenumi}{enumerate}{5}%
\setlist[compactenumi]{topsep=0pt,itemsep=-1ex,partopsep=1ex,parsep=1ex,%
   label=(\roman*)}%

\setcounter{secnumdepth}{5}

\newcommand{\HST}{\TermI{HST}\xspace}
\newcommand{\Term}[1]{\textsf{#1}}
\newcommand{\TermI}[1]{\Term{#1}\index{#1@\Term{#1}}}

\DeclareMathAlphabet{\mathpzc}{OT1}{pzc}{m}{it}

\newcommand{\seg}{\mathsf{s}}

\newcommand{\Diam}{\mathrm{diam}}
\newcommand{\diam}{\mathrm{diam}}

\newcommand{\Copt}{C_{\mathrm{opt}}}
\newcommand{\CenterPnt}{\mathrm{cen}}

\newcommand{\parent}{\overline{\mathrm{p}}}
\newcommand{\parentX}[2][\!]{\parent\pth[#1]{#2}}

\newcommand{\FMS}{\EuScript{X}\index{metric space}}  
\newcommand{\FMSB}{\EuScript{Y}\index{metric space}}  
\newcommand{\FastCut}{\Algorithm{FastCut}\xspace}
\newcommand{\MinCut}{\Algorithm{MinCut}\xspace}
\newcommand{\MinCutRep}{\Algorithm{MinCutRep}\xspace}
\newcommand{\Contract}{\Algorithm{Contract}\xspace}

\newcommand{\Grid}{\mathrm{G}}
\newcommand{\ClosestPair}{\mathcal{CP}}
\newcommand{\SphereChar}{\ensuremath{\mathbb{S}}}
\newcommand{\Sphere}[1]{{\SphereChar^{(#1)}}}

\newcommand{\PntSet}{{\mathsf{P}}}%
\providecommand{\PntSetA}{{\mathsf{Q}}}%

\newcommand{\query}{\mathcalb{q}}

\newcommand{\pnt}{\mathsf{p}}
\newcommand{\pntA}{\mathsf{u}}
\newcommand{\pntB}{\mathsf{v}}
\newcommand{\pntC}{{\mathsf{x}}}
\newcommand{\pntD}{{\mathsf{y}}}
\newcommand{\pntE}{{\mathsf{z}}}

\newcommand{\pntCoord}{p}
\newcommand{\pntCoordA}{q}

%%%%%%%%%%%%%%%%%%%%%%%%%%%%%%%%%%%%%%%%%%%%%%%%%%%%%%%%%%%%%%%%%%%%% 
%%%%%%%%%%%%%%%%%%%%%%%%%%%%%%%%%%%%%%%%%%%%%%%%%%%%%%%%%%%%%%%%%%%%%
% Markov chain notations.

\newcommand{\MCSpc}{\mathtt{S}}
\newcommand{\MCPrb}{\mathtt{P}}
\newcommand{\MCPrbY}[2]{\MCPrb_{#1#2}}
\newcommand{\MCPrbM}{\mathbf{P}}%
\newcommand{\MCPrbT}[3]{\MCPrb_{#1#2}^{(#3)}}
\newcommand{\MCFPrbT}[3]{\mathtt{r}_{#1#2}^{(#3)}}
\newcommand{\MCVstProb}[2]{\mathtt{f}_{#1#2}}
\newcommand{\MCExSteps}[2]{\mathtt{h}_{#1#2}}
\newcommand{\MCPVec}[1]{\mathbf{q}^{(#1)}}
%%%%%%%%%%%%%%%%%%%%%%%%%%%%%%%%%%%%%%%%%%%%%%%%%%%%%%%%%%%%%%%%%%%%%

\newcommand{\Spread}{\Phi}
\newcommand{\VC}{\ensuremath{\mathsf{VC}}\xspace}
\newcommand{\MC}{\ensuremath{\mathsf{MC}}\xspace}
\newcommand{\hyperplane}{h}

\newcommand{\QuickSort}{\Algorithm{QuickSort}\xspace}
\newcommand{\Alg}{\Algorithm{Alg}\xspace}

\newcommand{\Graph}{\mathsf{G}}%
\newcommand{\GraphB}{\mathsf{H}}
\newcommand{\GraphC}{\mathsf{K}}
\newcommand{\GraphD}{\mathsf{F}}

\providecommand{\G}{}
\renewcommand{\G}{\Graph}

\newcommand{\RLP}{\ComplexityClass{RLP}\xspace}
\newcommand{\USTCON}{\ComplexityClass{USTCON}}
\newcommand{\STCON}{\overrightarrow{\ComplexityClass{STCON}}}

\newcommand{\Mtrx}{\mathsf{M}}
\newcommand{\TMtrx}{\mathsf{Q}}
\newcommand{\IMtrx}{\mathcal{I}}

\newcommand{\Expander}[3]{$\left[#1,#2,#3\right]$-expander}

\newcommand{\ChartX}[1]{\ensuremath{\chi_{#1}}}
\newcommand{\Field}{\mathbb{F}}
\newcommand{\LSpace}{\mathbb{G}}
\newcommand{\LD}[2]{\mathsf{L{D}}\pth{#1,#2}}
\newcommand{\SEVN}{\widehat{\lambda_2}}
\newcommand{\NX}[2]{#1\!\left[#2\right]}
\newcommand{\Entry}[2]{#1\!\left[#2\right]}
\newcommand{\NXExt}[3]{#1_{#2}\!\left[#3\right]}
\newcommand{\RangeI}[1]{\llbracket #1 \rrbracket}
\newcommand{\DialX}[1]{\gamma\pth{#1}}
\newcommand{\BiDialX}[1]{{\gamma}_2\pth{#1}}

\newcommand{\EV}[1]{\widehat{\lambda_{#1}}}
\newcommand{\TEV}[1]{\widehat{\lambda_{#1}}}
\newcommand{\EVal}{\widehat{\lambda}}
\newcommand{\Basis}{\mathcal{B}}

\newcommand{\ZProd}[2]{\ensuremath{#1 \, \text{\textcircled{z}} #2}}
\newcommand{\RProd}[2]{\ensuremath{#1 \, \text{\textcircled{r}} #2}}

\newcommand{\OneVec}[1]{\mathbf{1}^{#1}}
\newcommand{\DstrbV}[1]{\mathbf{\mathsf{q}}^{(#1)}}
\newcommand{\DstrbE}[2]{\mathsf{q}^{(#1)}_{#2}}

\newcommand{\DEdges}{\widetilde{\Edges}}

\newcommand{\VV}{\mathsf{V}}%
\newcommand{\EE}{\mathsf{E}}%

\newcommand{\Vertices}{\VV}%
\newcommand{\Edges}{\EE}%

\newcommand{\EdgesX}[1]{\Edges\pth{#1}}
\newcommand{\VerticesX}[1]{V\pth{#1}}

\newcommand{\StringX}[1]{\mathsf{s}\pth{#1}}

\newcommand{\BMtrx}{\mathsf{B}}
\newcommand{\WMtrx}{\mathsf{W}}
\newcommand{\NWMtrx}{\overline{\mathsf{W}}}
\newcommand{\Sequence}{\EuScript{S}}

\newcommand{\med}{\mathrm{med}}
\newcommand{\medX}[1]{\mathrm{med}\pth{#1}}

\newcommand{\medianC}{\nu}%

\newlength{\savedparindent}
\newcommand{\SaveIndent}{\setlength{\savedparindent}{\parindent}}
\newcommand{\RestoreIndent}{\setlength{\parindent}{\savedparindent}}

\newcommand{\NNN}{\mathbb{N}}%
\providecommand{\NN}{\mathbb{N}}%
\renewcommand{\Re}{\mathbb{R}}%
\newcommand{\reals}{\Re}%
\newcommand{\ZZ}{\mathbb{Z}}%

\newcommand{\level}{\mathop{\mathrm{level}}}

\providecommand{\TPDF}[2]{\texorpdfstring{#1}{#2}}

\newcommand{\Hint}[1]{\emph{[#1]}}
\newcommand{\Arr}{\mathop{\mathrm{\EuScript{A}}}}

\newcommand{\Array}{\mathcal{T}}%

\newcommand{\dirEdge}[2]{\left({{#1} \rightarrow {#2}} \right)}

%%%%%%%%%%%%%%%%%%%%%%%%%%%%%%%%%%%%%%%%%%%%%%%%%%%%%%%%%%%%%%%%%%%%%%
% Emphasize stuff...
\definecolor{blue25}{rgb}{0, 0, 11}
\newcommand{\emphic}[2]{%
%   \BADVERYBAD%
   \textcolor{blue25}{%
      \bfseries{\emph{#1}}}%
   \index{#2}}
\newcommand{\emphi}[1]{\emphic{#1}{#1}}
\newcommand{\emphOnly}[1]{\textcolor{blue25}{\bfseries{\emph{#1}}}}
\providecommand{\emphiC}[1]{\emphic{#1}{#1}}


\definecolor{almostblack}{rgb}{0, 0, 0.3}
\providecommand{\emphw}[1]{{\textcolor{almostblack}{\emph{#1}}}}%

%
%%%%%%%%%%%%%%%%%%%%%%%%%%%%%%%%%%%%%%%%%%%%%%%%%%%%%%%%%%%%%%%%%%%%%%

\newcommand{\volX}[1]{\mathrm{vol}\pth{#1}}

\newcommand{\Interval}{\EuScript{I}}
\newcommand{\IntervalA}{\EuScript{I}}

\newcommand{\Exp}{\mathrm{Exp}}

\let\Ball\undefined
\newcommand{\Ball}{\mathbf{b}}
\newcommand{\hdim}{d}
\newcommand{\Domain}{\mathcal{D}}
\newcommand{\points}[1]{[\bfseries{#1 points}]}
\newcommand{\Event}{\EuScript{E}}%
\newcommand{\EventB}{\EuScript{B}}%
\newcommand{\EventF}{\EuScript{F}}%
\newcommand{\EventG}{\EuScript{G}}%

\newcommand{\distCmd}[1]{\left\| {#1}  \right\|}
\newcommand{\distY}[2]{\distCmd{#1 - #2}}

\newcommand{\NormalC}{\mathbf{N}}
\newcommand{\NormalY}[2]{\NormalC\pth{#1, #2}}%
\newcommand{\NormalDistOneDim}{\NormalC}
\newcommand{\NormalDistX}[1]{\NormalC^{#1}}
\DeclareMathOperator{\lca}{lca}
\newcommand{\RangeSpace}{\mathsf{S}}%
\newcommand{\RangeSpaceA}{\mathsf{T}}

\newcommand{\range}{\mathbf{r}}
\newcommand{\rangeA}{\mathbf{r}'}

\newcommand{\BadProb}{\varphi}

\newcommand{\EntropyChar}{\mathbb{H}}
\newcommand{\Entropy}[1]{\EntropyChar\pth{#1}}
\newcommand{\EntropyX}[1]{\EntropyChar\pth[]{#1}}

\newcommand{\Matrix}{\mathsf{M}}
\newcommand{\MatrixB}{\mathsf{B}}
\newcommand{\MatrixC}{\mathsf{C}}
\newcommand{\MatrixD}{\mathsf{D}}

\renewcommand{\th}{th\xspace}

\newcommand{\Family}{\mathcal{F}}
\newcommand{\FamilyB}{\mathcal{G}}
\newcommand{\FamilyC}{\mathcal{H}}
\newcommand{\FMSC}{\EuScript{Y}\index{metric space}}

\newcommand{\median}{\overline{\mathsf{m}}}
\newcommand{\QuickSelect}{\Algorithm{Quick{}Select}\xspace}

\providecommand{\row}{\mathbf{r}}%
\newcommand{\rowc}{{r}}%

\newcommand{\vecb}{\mathbf{b}}

\newcommand{\vecA}{\mathbf{v}}
\newcommand{\vecB}{\mathbf{u}}
\newcommand{\vecRand}{\mathbf{r}}

\newcommand{\vecH}{\mathbf{h}}
\newcommand{\vecV}{\mathbf{v}}

\newcommand{\vecX}{\mathbf{x}}
\newcommand{\vecY}{\mathbf{y}}

\newcommand{\DistBinom}[2]{\mathrm{Bin}\pth{#1,#2}}

\newcommand{\DistGeom}[1]{\mathrm{Geom}\pth{#1}}

\newcommand{\Sample}{\mathsf{R}}

\newcommand{\lcm}{\mathop{\mathrm{lcm}}}
\newcommand{\bdiv}{\mathop{\mathrm{div}}}%
\newcommand{\divides}{\mid}%
\newcommand{\poly}{\mathrm{poly}}%

\newcommand{\Group}{\mathcal{G}} 
\newcommand{\GroupA}{\mathcal{H}} 

\newcommand{\algGCD}{\Algorithm{EuclidGCD}\xspace}%

\newcommand{\gOp}{\times}%
\newcommand{\gIdent}{\mathsf{i}}%

\newcommand{\ordX}[2][\!]{\mathop{\mathrm{o{r}d}}\pth[#1]{#2}}

\newcommand{\lSym}[2]{%
   \left(
       \ds {#1}\divides{#2}%
   \right) %
}
\newcommand{\jSym}[2]{\left\llbracket%
       \ds {#1}\divides{#2}%
       \right\rrbracket%
}
\newcommand{\Line}{\ell}

\newcommand{\LevelChar}{\mathbb{L}}
\newcommand{\AtMostKLevelChar}{\LevelChar_{\leq k}}
\newcommand{\AtMostKLevel}[1]{\LevelChar_{\leq k} \pth{ #1}}

\newcommand{\LineSet}{\mathsf{L}}%
\newcommand{\ArrX}[1]{\EuScript{A}\pth{#1}}%

\providecommand{\Matousek}{Matou{\v s}ek\xspace}
\providecommand{\Erdos}{Erd{\H o}s\xspace}
\providecommand{\Lovasz}{Lov\'asz\xspace}

\newcommand{\ProblemC}[1]{{\textsf{\textcolor{red}{#1}\index{problem!#1}}}}
\DefineNamedColor{named}{OliveGreen} {cmyk}{0.64,0,0.95,0.40}
\providecommand{\ComplexityClass}[1]{{{\textcolor[named]{OliveGreen}{%
            \textsc{#1}}}}}
\providecommand{\NPHard}{{\ComplexityClass{NP-Hard}}%
   \index{NP!hard}\xspace}
\providecommand{\NPComplete}{\ComplexityClass{NP-Complete}%
   \index{NP!complete}\xspace}
\providecommand{\APXHard}{\ComplexityClass{APX-Hard}%
   \index{APX!Hard}\xspace}
\providecommand{\NPC}{\ComplexityClass{NPC}%
   \index{NP!complete}\xspace}
\providecommand{\NP}{\ComplexityClass{NP}%
   \index{NP}\xspace}
\providecommand{\POLYT}{\ComplexityClass{P}\xspace}

\newcommand{\envFrame}[1]{\begin{center}\fbox{\begin{minipage}[c]{0.75\textwidth}{#1}\end{minipage}}\end{center}}

\newcommand{\mysqrt}[1]{\ts\!\!\sqrt{#1}}
\newcommand{\SegSet}{\mathsf{S}}
\newcommand{\ArrVD}[1]{\Arr^{|} \pth{#1}}
\newcommand{\RIC}{\textsf{RIC}\index{random incremental
      construction}\xspace}
\newcommand{\ArrB}{\EuScript{B}}

\newcommand{\TreeOpt}{\EuScript{T}_\mathrm{opt}}
\newcommand{\TreeOptA}{\EuScript{T}_\mathrm{opt}'}

\newcommand{\DepthExt}[2]{\mathrm{depth}_{#1}\pth{#2}}
\newcommand{\depthX}[1]{\mathrm{depth}\pth{#1}}

\newcommand{\TreeA}{\EuScript{U}}
\newcommand{\TreeB}{\EuScript{V}}

\newcommand{\Trap}{\sigma}
\newcommand{\TrapA}{\sigma'}
\newcommand{\TrapB}{\sigma''}

\newcommand{\SampleA}{\mathsf{R}'}

\newcommand{\AllRegions}{\mathcal{T}}
\newcommand{\RegionsSet}[1]{\mathcal{F}\pth{#1}}
\newcommand{\RegionsHeavy}[2]{\mathcal{F}_{\!\geq#1}\pth{#2}}
\newcommand{\NumOfRegionsZero}[1]{\ExChar \!f\pth{#1}}
\newcommand{\NumOfRegions}[2]{\ExChar\! f_{\!\geq #1}\pth[]{#2}}
\newcommand{\ObjSet}{\mathsf{S}}

\newcommand{\WeightX}[2][\!]{\omega \pth[#1]{#2}}
\newcommand{\Edge}{\mathsf{e}}

\newcommand{\Face}{\mathsf{f}}

\newcommand{\DefSet}[1]{D\pth{#1}}
\newcommand{\KillSet}[1]{K\pth{#1}}

\newcommand{\CSProbFull}[4]{
   \,\rho_{\!\!\stackrel{\rule[-.005cm]{0cm}{0.1cm}}{#1,
       #4}}\pth{\MakeSBig #2,#3}}

\newcommand{\disk}{D}
\newcommand{\diskA}{D'}

\newcommand{\IntervalSet}{\widehat{\mathcal{I}}}

\newcommand{\GroundSet}{\textsf{X}}
\newcommand{\FGroundSet}{\mathsf{x}}
\newcommand{\RangeSet}{{\mathcal{R}}}
\newcommand{\RangeSetA}{\mathcal{R}'}
\newcommand{\MeasureChar}{\overline{m}}
\newcommand{\Measure}[1]{\MeasureChar\pth{#1}}
\newcommand{\sMeasure}[1]{\overline{s}\pth{#1}}

\newcommand{\SetA}{S}%
\newcommand{\SetB}{T}%
\newcommand{\SetC}{N}%

\newcommand{\VCProj}[2]{#1_{|#2}}
\newcommand{\VCDim}{\mathrm{dim}_{\VC}}
\newcommand{\Dim}{\delta}
\newcommand{\DimA}{\delta'}
\newcommand{\PntSetB}{{{\mathsf{Q}}}}
\newcommand{\PntSetC}{{{\mathsf{R}}}}
\newcommand{\Growth}[2]{{\EuScript{G}}_{#1}\pth{#2}}
\newcommand{\ShatterFunc}[1]{\pi_{#1}}
\newcommand{\Distribution}{\EuScript{D}}
\newcommand{\Error}{\EuScript{E}}

\newcommand{\bd}{{\partial}}%
\newcommand{\eps}{{\varepsilon}}%

\newcommand{\clmlab}[1]{\label{claim:#1}}
\newcommand{\clmref}[1]{\HLink{Claim}{claim:#1}}

\newcommand{\CHX}[1]{{\mathcal{CH}}\pth{#1}}

\newcommand{\ts}{\hspace{0.6pt}}

\def\qedsymbol{\myqedsymbol}%

\newcommand{\qedsign}{\hfill{\hfill\myqedsymbol}}
\newcommand{\qedsignBig}{\hfill{\hfill\rule{4mm}{4mm}}}
\newcommand{\QED}{{\hfill{\hfill\qedsymbol}}}

\newcommand{\num}{\alpha}

\newcommand{\NEW}[1]{#1}

\newcommand{\ThreeFiguresExt}[9]{%
   \begin{tabular}{ccc}
     \begin{minipage}{0.31\linewidth}
         \begin{center}
             {\IncludeGraphics[#1]{#2}}
         \end{center}
     \end{minipage}
     &
       \begin{minipage}{0.31\linewidth}
           \centerline{\IncludeGraphics[#4]{#5}}
       \end{minipage}
     &
       \begin{minipage}{0.31\linewidth}
           \centerline{\IncludeGraphics[#7]{#8}}
       \end{minipage}
     \\
     ~\\[-0.2cm]
     {#3} & {#6} & {#9}
   \end{tabular}}

\newcommand{\IncludeGraphics}[2][]{%
   \typeout{}%
   \typeout{Graphics: #2}%
   \typeout{\ includegraphics[#1]{#2}}%
   \includegraphics[#1]{#2}
   \typeout{}%
}

%    -----------------------------------------------------------------
%    -----------------------------------------------------------------
%    -----------------------------------------------------------------
%    Handling references
%    -----------------------------------------------------------------
%    -----------------------------------------------------------------
%    -----------------------------------------------------------------
\newcommand{\HLinkShort}[2]{\hyperref[#2]{#1\ref*{#2}}}
\newcommand{\HLink}[2]{\hyperref[#2]{#1~\ref*{#2}}}
\newcommand{\HLinkPage}[2]{\hyperref[#2]{#1~\ref*{#2}%
      $_\text{p\pageref{#2}}$}}

\newcommand{\HLinkSuffix}[3]{\hyperref[#2]{#1\ref*{#2}{#3}}}
\newcommand{\HLinkPageSuffix}[3]{\hyperref[#2]{#1\ref*{#2}%
      #3$_\text{p\pageref{#2}}$}}

\newcommand{\figlab}[1]{\label{fig:#1}}
\newcommand{\figref}[1]{\HLink{Figure}{fig:#1}}

\newcommand{\exelab}[1]{\label{exer:#1}}
\newcommand{\exeref}[1]{Exercise~\ref{exer:#1}}

\providecommand{\deflab}[1]{\label{def:#1}}
\newcommand{\defref}[1]{\HLink{Definition}{def:#1}}
\newcommand{\defrefY}[2]{\hyperref[def:#2]{#1}}
\newcommand{\defrefpage}[1]{\HLinkPage{Definition}{def:#1}}

\newcommand{\etal}{\textit{et~al.}\xspace}

\newcommand{\obslab}[1]{\label{observation:#1}}
\newcommand{\obsref}[1]{\HLink{Observation}{observation:#1}}
\newcommand{\obsrefY}[2]{\hyperref[observation:#2]{#1}}

\newcommand{\exmlab}[1]{\label{example:#1}}
\newcommand{\exmref}[1]{\HLink{Example}{example:#1}}

\newcommand{\remlab}[1]{\label{rem:#1}}
\newcommand{\remref}[1]{\HLink{Remark}{rem:#1}}%

\newcommand{\chaplab}[1]{\label{chap:#1}}

\newcommand{\itemlab}[1]{\label{item:#1}}
\newcommand{\itemref}[1]{\HLinkSuffix{}{item:#1}{}}

\newcommand{\tedlab}[1]{\label{tedium:#1}}
\newcommand{\tedref}[1]{\HLink{Tedium}{tedium:#1}}%

\newcommand{\factlab}[1]{\label{fact:#1}}%

\newcommand{\alglab}[1]{\label{algorithm:#1}}

\newcommand{\net}{\mathcal{N}}%
      
\newcommand{\problab}[1]{\label{prob:#1}}%
\newcommand{\probref}[1]{\HLink{Problem}{prob:#1}}%
\newcommand{\probrefY}[2]{\hyperref[prob:#1]{#2}}

\newcommand{\lemlab}[1]{\label{lemma:#1}}
\newcommand{\lemref}[1]{\HLink{Lemma}{lemma:#1}}%
\newcommand{\lemrefY}[2]{\hyperref[lemma:#2]{#1}}
\newcommand{\lemrefpage}[1]{\HLinkPage{Lemma}{lemma:#1}}
\newcommand{\lemrefshort}[1]{\HLinkShort{L}{lemma:#1}}

\newcommand{\eqlab}[1]{\label{equation:#1}}
\newcommand{\Eqrefbare}[1]{\HLinkSuffix{(}{equation:#1}{)}}
\newcommand{\Eqref}[1]{\HLinkSuffix{Eq.~(}{equation:#1}{)}}
\newcommand{\Eqrefpage}[1]{\HLinkPageSuffix{Eq.~(}{equation:#1}{)}}

\newcommand{\seclab}[1]{\label{sec:#1}}
\newcommand{\secref}[1]{\HLink{Section}{sec:#1}}
\newcommand{\secrefY}[2]{\hyperref[sec:#2]{#1}}
\newcommand{\secrefpage}[1]{\HLinkPage{Section}{sec:#1}}

\newcommand{\thmlab}[1]{{\label{theo:#1}}}
\newcommand{\thmref}[1]{\HLink{Theorem}{theo:#1}}
\newcommand{\thmrefY}[2]{\hyperref[theo:#2]{#1}}
\newcommand{\thmrefpage}[1]{\HLinkPage{Theorem}{theo:#1}}

\newcommand{\corlab}[1]{\label{cor:#1}}
\newcommand{\corref}[1]{\HLink{Corollary}{cor:#1}}%

%    -----------------------------------------------------------------
%    -----------------------------------------------------------------
%    -----------------------------------------------------------------
%    Handling references
%    -----------------------------------------------------------------
%    -----------------------------------------------------------------
%    -----------------------------------------------------------------

\newcommand{\ceil}[1]{\left\lceil {#1} \right\rceil}
\newcommand{\floor}[1]{\left\lfloor {#1} \right\rfloor}
\newcommand{\brc}[1]{\left\{ {#1} \right\}}
\newcommand{\pth}[2][\!]{\mleft({#2}\mright)}%
\newcommand{\brk}[1]{\mleft[{#1}\mright]}%
\newcommand{\pbrc}[2][\!\!]{#1\left[ {#2} \bigr. \right]}
\newcommand{\pbrcS}[1]{\left[ {#1} \right]}

\newcommand{\dotprod}[2]{{#1}\cdot{#2}}
\newcommand{\DotProd}[2]{\permut{{#1},{#2}}}
\newcommand{\permut}[1]{\left\langle {#1} \right\rangle}

\newcommand{\MakeBig}{}
\newcommand{\MakesBig}{}
\newcommand{\MakeSBig}{}
\newcommand{\MakeVBig}{}

\newcommand{\DistChar}{\mathsf{d}}%
\newcommand{\DistSetY}[2]{\DistChar\pth{#1,#2}}
\newcommand{\distG}[2]{\DistChar_\Graph\pth{#1, #2}}
\newcommand{\DistMtrCharX}[1]{\DistChar_{#1}}

\newcommand{\dY}[2]{\DistChar\pth{#1, #2}}%
\newcommand{\dGY}[2]{\DistChar_\Graph\pth{#1, #2}}%

\newcommand{\cset}{\mathcal{C}}%

\newcommand{\MTR}{\mathcal{M}}
\newcommand{\DistMtr}[2]{\DistChar\pth{#1,#2}}

\newcommand{\dist}[1]{\left\| {#1} \right\|}
\newcommand{\norm}[1]{\left\| {#1} \right\|}

\providecommand{\ds}{\displaystyle}
\newcommand{\sep}[1]{\,\left|\, {#1} \bigr.\right.}%
\newcommand{\sepw}[1]{\left|\, {#1} \right.}%

\newcommand{\ProbExt}[2]{\mathop{\mathbf{Pr}}_{#1}\pbrcx{#2}}
\newcommand{\ProbLTR}{\mathbb{P}}%

\newcommand{\Prob}[1]{\mathop{\ProbLTR} \mleft[ #1 \mright]}%
\newcommand{\ProbCond}[2]{\mathop{\ProbLTR}\!\left[%
       #1 \;\middle\vert\; #2 \right]}
\newcommand{\ProbY}[2]{\ProbLTR^{}_{\!#1}\pbrc{#2}}

\newcommand{\ExChar}{\mathbb{E}}%
\newcommand{\ExSym}{\mathop{\ExChar}}%
\newcommand{\Ex}[2][\!]{\ExSym#1\pbrcx{#2}}
\newcommand{\ExExt}[2]{\ExSym_{#1}\!\pbrcx{#2}}
\newcommand{\ExExtS}[2]{\ExSym_{\substack{#1}}\!\pbrcx{#2}}
\newcommand{\ExCond}[2]{\ExSym\!\left[%
       #1 \;\middle\vert\; #2 \right]}

\newcommand{\VarSymb}{\mathop{\mathbb{V}}}
\newcommand{\Var}[1]{\mathop{\mathbb{V}}\!\pbrcx{#1}}

\newcommand{\pbrcx}[1]{\left[ {#1} \right]}%
\newcommand{\nfrac}[2]{{#1}/{#2}}
\newcommand{\cardin}[1]{\left| {#1} \right|}%

\newcommand\CCancel[2][black]{{\renewcommand\CancelColor{\color{#1}}\cancel{#2}}}

\newcommand{\IncGraphPage}[4][]{%
   \IncludeGraphics[page=#4,#1]{\File{#2/#3}}%
}

\newenvironment{enumEx}{
   \begin{compactenumA}
    }{\end{compactenumA}}

\newcommand{\cl}{{\mathrm{cl}}}

\usepackage[amsmath,thmmarks]{ntheorem}
\theoremseparator{.}%

\newtheorem{theorem:unnumbered}{Theorem}[section]
\newtheorem{theorem}{Theorem}[section]

\newtheorem{observation}[theorem]{Observation}
\newtheorem{conjecture}[theorem]{Conjecture}
\newtheorem{lemma}[theorem]{Lemma}
\newtheorem{corollary}[theorem]{Corollary}
\newtheorem{claim}[theorem]{Claim}
\newtheorem{assumption}[theorem]{Assumption}
\newtheorem{tedium}[theorem]{Tedium}

\theoremstyle{plain}%
\theoremheaderfont{\sffamily}
\theorembodyfont{\upshape}%

\newtheorem{defn}[theorem]{Definition}
\newtheorem{definition}[theorem]{Definition}

\newtheorem*{remark_unnumbered}[theorem:unnumbered]{Remark}%
\newtheorem*{remarks}[theorem]{Remarks}%
\newtheorem{remark}[theorem]{Remark}%
\newtheorem{fact}[theorem]{Fact}%
\newtheorem{example}[theorem]{Example}
\newtheorem{exercise}[theorem]{Exercise}
\newtheorem{problem}[theorem]{Problem}
\newtheorem{algorithm}[theorem]{Algorithm}%
\newtheorem{xca}[theorem]{Exercise}
\newtheorem{openproblem}[theorem]{Open Problem}
\newtheorem{exercise_h}[theorem]{Exercise}

% Proof environments
\theoremheaderfont{\em}%
\theorembodyfont{\upshape}%
\theoremstyle{nonumberplain}
\theoremseparator{}
% TeX4ht's ntheorem support consumes theorem end symbols.  The web-only
% wrappers below place the duck before ntheorem closes its HTML block.
\theoremsymbol{}
\newtheorem{proof}{Proof:}
\newtheorem{proofextin}{Proof}
\newtheorem{altproof}{Alternative proof:}

\theoremsymbol{}
\newtheorem{proofnoqed}{Proof:}

\theoremsymbol{}
\theorempreskip{0pt}%
\theorempostskip{0pt}%
\newtheorem{proofof}{Proof of\!}%

\makeatletter
\let\proof@theorem\proof
\let\endproof@theorem\endproof
\renewenvironment{proof}{\proof@theorem}{\unskip\myqedsymbol\endproof@theorem}
\let\proofextin@theorem\proofextin
\let\endproofextin@theorem\endproofextin
\renewenvironment{proofextin}{\proofextin@theorem}{\unskip\myqedsymbol\endproofextin@theorem}
\let\altproof@theorem\altproof
\let\endaltproof@theorem\endaltproof
\renewenvironment{altproof}{\altproof@theorem}{\unskip\myqedsymbol\endaltproof@theorem}
\let\proofof@theorem\proofof
\let\endproofof@theorem\endproofof
\renewenvironment{proofof}{\proofof@theorem}{\unskip\myqedsymbol\endproofof@theorem}
\makeatother

\newenvironment{proofExt}[1]{{\emph{#1:}}}{\hfill{\hfill\myqedsymbol}}
\newenvironment{proofnoproof}{{}}{{\hfill{\hfill\myqedsymbol}}}

%%%%%%%%%%%%%%%%%%%%%%%%%%%%%%%%%%%%%%%%%%%%%%%%%%%%%%%%%%%%%%%%%%%%
%%%%%%%%%%%%%%%%%%%%%%%%%%%%%%%%%%%%%%%%%%%%%%%%%%%%%%%%%%%%%%%%%%%%
% Macros to embed figures in text

\newlength{\RecordedWidth}
\newcommand{\RecordWidth}[1]{\settowidth{\RecordedWidth}{#1}}%

\newlength{\msize}%

%-------------------------------------------------------------------
% \minipageWFigure{ TEXT }{ Figure }{ Figure caption/label }
%-------------------------------------------------------------------
\newcommand{\minipageWFigure}[3]{%
   \RecordWidth{#2}%
   \setlength{\msize}{\dimexpr(\textwidth-1.1\RecordedWidth)\relax}
   \noindent%
   \SaveIndent%
   \begin{minipage}{\msize}%
       \RestoreIndent%
       #1
   \end{minipage}%
   \null\hfill%
   \begin{minipage}{\RecordedWidth}
       \null\hfill%
       #2%       

       #3%
   \end{minipage}%
}

\newcommand{\EBRY}[2]{#1_{\langle#2\rangle}}%

\newcommand{\NiceFrac}[2]{%
   {\raisebox{0.2em}{\ensuremath{#1}}\!\!\!}%
   {\rotatebox[origin=c]{-20}{\Big/}}% Rotated /
   {\! \raisebox{-0.2em}{\ensuremath{#2}} }}%

% end - Macros to embed figures in text
%%%%%%%%%%%%%%%%%%%%%%%%%%%%%%%%%%%%%%%%%%%%%%%%%%%%%%%%%%%%%%%%%%%% 
%%%%%%%%%%%%%%%%%%%%%%%%%%%%%%%%%%%%%%%%%%%%%%%%%%%%%%%%%%%%%%%%%%%%

%%%%%%%%%%%%%%%%%%%%%%%%%%%%%%%%%%%%%%%%%%%%%%%%%%%%%%%%%%%%%%%%%%
% Color in math
%\providecommand{\Mh}[1]{{#1}}%
\providecommand{\IfColorMath}[2]{#2}%

\IfFileExists{.latex_printer_friendly}{\def\GenPrinterVer{1}}{}%

\ifx\GenPrinterVer\undefined
   \IfFileExists{.latex_color}{\def\GenColorMath{1}}{}
\else
   \renewcommand{\IfPrinterVer}[2]{#1}%
\fi

\ifx\GenColorMath\undefined
\else
\renewcommand{\IfColorMath}[2]{#1}%
%\renewcommand{\Mh}[1]{{\textcolor{red}{#1}}}%
\fi
% Color in math - end
%%%%%%%%%%%%%%%%%%%%%%%%%%%%%%%%%%%%%%%%%%%%%%%%%%%%%%%%%%%%%%%%%%

\newcommand{\HSetY}[2]{H\pth{#1, #2}}%
\newcommand{\CHSetY}[2]{C\pth{#1, #2}}%
\newcommand{\Holder}{H\"older\xspace}
\newcommand{\OldOm}{{\mathcal{O}}}
\newcommand{\NewOm}{{\mathcal{N}}}

%%%%%%%%%%%%%%%%%%%%%%%%%%%%%%%%%%%%%%%%%%%%%%%%%%%%%%%%%%%%%%%%%%%%%%%
%---------------------------------------------------------------
\newcommand{\IntRange}[1]{\left\llbracket #1 \right\rrbracket}
\newcommand{\IRX}[1]{\left\llbracket #1 \right\rrbracket}%
\newcommand{\IRY}[2]{\left\llbracket #1:#2 \right\rrbracket}

%---------------------------------------------------------------
%%%%%%%%%%%%%%%%%%%%%%%%%%%%%%%%%%%%%%%%%%%%%%%%%%%%%%%%%%%%%%%%%%%%%%%

\newcommand{\obj}{f}%
\newcommand{\Wg}{\omega}%
\newcommand{\TotalW}{W}%

\newcommand{\rankY}[2]{{#1}_{ \permut{#2}}}%%
\newcommand{\areaX}[1]{\mathrm{area}\pth{#1}}%%

\newcommand{\mybinom}[2]%
{\Bigl(\begin{array}{@{}c@{}}#1\\#2\end{array}\Bigr)}

\providecommand{\Set}[2]{\left\{ #1 \;\middle\vert\; #2 \right\}}
\newcommand*{\biga}[1]{\vcenter{\hbox{\scalebox{1.5}{\ensuremath#1}}}}
\newcommand{\CNF}{\TermI{CNF}\xspace}%
\newcommand{\Tree}{T}%

\DeclareSymbolFont{sfletters}{OML}{cmbrm}{m}{it}

\DeclareMathSymbol{\salpha}{\mathord}{sfletters}{"0B}

%%%%%%%%%%%%%%%%%%%%%%%%%%%%%%%%%%%%%%%%%%%%%%%%%%%%%%%%%%%%%%%%%%%%

\providecommand{\ProblemInd}[2]{\textsf{\textcolor{red}{#1}\index{Problem!#2}}}
\newcommand{\NPProblemInd}[4]{
   \begin{center}
      \vspace{0.6cm}
       \fbox{
          \begin{minipage}{0.9\linewidth}
              \vspace{-0.6cm} \hspace{-0.35cm}
              {\large \bfseries \ProblemInd{#1}{#2}}
              \vspace{+0.2cm}%

              \noindent{\textcolor{blue}{\bfseries \textsf{Instance}}:} {#3}
              
              \noindent{\textcolor{blue}{\bfseries \textsf{Question}:}} {#4}
          \end{minipage}
       }
   \end{center}
   \smallskip
   }

\newcommand{\NPProblemOpt}[4]{\NPProblemInd{#1}{#2}{#3}{#4}}

%%%%%%%%%%%%%%%%%%%%%%%%%%%%%%%%%%%%%%%%%%%%%%%%%%%%%%%%%%%%%%%%%%%%

\newcommand{\Hastad}{H{\aa}stad\xspace}

\newcommand{\ANNbr}{\Term{ANN}\index{approximate nearest neighbor}\xspace}
\newcommand{\LSH}{\TermI{L{S}H}\xspace}

\newcommand{\DS}{\EuScript{D}} % DataStructure
\newcommand{\DSC}{\mathsf{D}} % DataStructure

\newcommand{\distH}[2]{\DistChar_{H}\pth{#1, #2}}
\newcommand{\distHY}[2]{\DistChar_{H}\pth{#1, #2}}

\newcommand{\DSAprxNearNbr}{\ensuremath{\DataStructure{{\approx%
         \text{Near{}}}}}%
   \index{near neighbor!data-structure!approximate}%
}
\newcommand{\DataStructure}[1]{\ensuremath{\DS_{#1}}}
\newcommand{\DistNN}[2]{\DistChar\pth{#1,#2}}
\newcommand{\pntc}{{p}}%

\newcommand{\seq}{{M}}%
\newcommand{\cTimes}{\beta}%

\newcommand{\DSTimes}{L}%
\newcommand{\rad}{r}%

\newcommand{\HypercubeX}[1]{\Hypercube^{#1}}
\newcommand{\kFunctions}[2]{\mathrm{combine}\pth{#1,#2}}
\newcommand{\WEqualChar}{\Bumpeq}
\newcommand{\WEqual}[2]{{#1}\WEqualChar{#2}}
\newcommand{\FamilyAmp}{\FamilyC_{\WEqualChar}}
\newcommand{\PSPACE}{\ComplexityClass{PSPACE}\xspace}%

\newcommand{\PS}{P}%

\newcommand{\PSetFar}{\ensuremath{\PntSet^{}_{\geq}\xspace}}%
\newcommand{\DA}{D}%
\newcommand{\uniqX}[1]{\mathrm{uniq}\pth{#1}}%
\newcommand{\diffX}[1]{\mathop{}\!\mathrm{d}#1}

\newcommand{\Frq}{F}%
\newcommand{\frq}{f}%
\newcommand{\STRM}{\mathcal{S}}%
\newcommand{\strm}{{s}}%

\providecommand{\Turan}{Tur\'an\xspace}
\newcommand{\fGrowX}[1]{G^{}_{\Dim}\pth{#1}}
\newcommand{\CompressFunc}{\mathtt{Compress}}
\newcommand{\RSA}{\Term{RSA}\xspace}
\newcommand{\MST}{\Term{MST}\xspace}
\providecommand{\CClassPP}{\CClass{P{}P}\xspace}
\providecommand{\ID}{\mathrm{id}}

\newcommand{\EncodeFunc}{\text{\ttfamily{Enc}}}
\newcommand{\DecodeFunc}{\text{\ttfamily{Dec}}}
\newcommand{\discX}[1]{\mathrm{disc}\pth{#1}}%

\newcommand{\propertyC}{\psi}%
\newcommand{\propertyX}[1]{\propertyC\pth{#1}}%
\newcommand{\numC}{\Gamma}%

\newcommand{\degreeX}[1]{\mathrm{degree}\pth{#1}}

\newcommand{\ind}{m}

\providecommand{\nequiv}{\not\equiv}
\providecommand{\sqcupplus}{\sqcup^+}

\newcommand{\AlgorithmI}[1]{{\protect\color{RedViolet}{#1}}}

\usepackage{marginnote}
\usepackage{ragged2e} 
% Short, subtle snark macro
\setlength{\marginparwidth}{1.5in}
\setlength{\marginparsep}{10pt}
\newcommand{\snark}[1]{%
   \marginnote{\raggedright\scriptsize\sffamily\color{gray}#1}%
}

\newcommand{\BigH}{\rule[-0.4em]{0pt}{1.3em}}%

\newcommand{\binsY}[2]{\sqcup_{\geq #1}\pth{#2}}
\newcommand{\ballsY}[2]{\text{\Football}_{\geq #1}\pth{#2}}

\newcommand{\BLX}[1]{{{\sqcupplus{}}_{\geq #1}}}

% Page-breaking hints have no meaning in HTML.  Display bodies are preserved
% verbatim for MathJax, so the builder also removes occurrences found there.
\renewcommand{\allowdisplaybreaks}[1][]{}

\let\Mh\undefined


% -----------------------------------------------------------------------------
% Bibliography & Document Begin
%-----------------------------------------------------------------------------
\IfFileExists{\jobname.bib}{%
   \addbibresource{\jobname.bib}%
}{%
   \addbibresource{shortcuts.bib}%
   \addbibresource{geometry.bib}%
}
\makeindex
\begin{document}
}{\input{prefix_4ht}}%
  \else
    \IfFileExists{styles/prefix_latex.tex}{%-----------------------------------------------------------------------------
% Master Prefix File for Randomized Algorithms Class Notes & Book
% Unified self-contained configuration for LuaLaTeX and XeLaTeX with Memoir
%-----------------------------------------------------------------------------

\RequirePackage{iftex}
\ifPDFTeX
  \PackageError{engine-guard}{This document requires LuaLaTeX or XeLaTeX}{%
    Please recompile using LuaLaTeX or XeLaTeX instead of pdfLaTeX.%
  }
  \stop
\fi

\ifdefined\LWR@started
  \PackageError{latex-gate}{%
    Attempting to run lwarp with prefix_latex.tex.%
  }{Please use prefix_lwarp.tex (or set \string\warpMode) for HTML generation.}
  \stop
\fi

% Always use memoir
\documentclass[12pt,oneside,fleqn]{memoir}

%-----------------------------------------------------------------------------
% Typography, Fonts & Math Packages
%-----------------------------------------------------------------------------
\usepackage{amssymb}
\usepackage{mathtools}
\PassOptionsToPackage{warnings-off={mathtools-overbracket,mathtools-colon}}{unicode-math}
\usepackage{fontspec}
\usepackage{unicode-math}

% New Computer Modern OpenType text and math fonts
\setmainfont{NewCM10-Book.otf}[
  BoldFont          = NewCM10-Bold.otf,
  ItalicFont        = NewCM10-BookItalic.otf,
  BoldItalicFont    = NewCM10-BoldItalic.otf,
  SlantedFont       = NewCM10-Book.otf,
  SlantedFeatures   = {FakeSlant = 0.25},
  BoldSlantedFont   = NewCM10-Bold.otf,
  BoldSlantedFeatures = {FakeSlant = 0.25},
  SmallCapsFeatures = {Letters = SmallCaps}
]
\setsansfont{NewCMSans10-Book.otf}[
  BoldFont          = NewCMSans10-Bold.otf,
  ItalicFont        = NewCMSans10-BookOblique.otf,
  BoldItalicFont    = NewCMSans10-BoldOblique.otf
]
\setmonofont{NewCMMono10-Book.otf}[
  BoldFont          = NewCMMono10-Bold.otf,
  ItalicFont        = NewCMMono10-BookItalic.otf,
  BoldItalicFont    = NewCMMono10-BoldOblique.otf
]

\setmathfont{NewCMMath-Book.otf}
\setmathfont{NewCMMath-Bold.otf}[version=bold]

% OpenType math aliases and fixes
\providecommand{\square}{\mdlgwhtsquare}
\def\Box{\mdlgwhtsquare}
\def\Diamond{\mdlgwhtdiamond}
\def\blacksquare{\mdlgblksquare}
\def\leadsto{\rightsquigarrow}
\providecommand{\llbracket}{\lBrack}
\providecommand{\rrbracket}{\rBrack}
\DeclareTextCommand{\textcircled}{TU}[1]{%
  \mbox{\ooalign{\hfil\raise0.1ex\hbox{#1}\hfil\crcr\char"25EF}}%
}

\usepackage{alphabeta}
\usepackage[normalem]{ulem} % Provides \sout
\usepackage{mleftright}
\usepackage{microtype}
\usepackage{xspace}
\usepackage{xurl}
\usepackage[dvipsnames,table]{xcolor}
\usepackage{graphicx}
\usepackage{caption}
\usepackage{booktabs,array,colortbl,hhline,multirow,tabularx}
\usepackage[inline]{enumitem}
\usepackage[thicklines]{cancel}
\renewcommand\CancelColor{\color{purple}}
\usepackage{euscript}
\usepackage{marvosym}
\let\Cross\undefined
\let\Square\undefined
\usepackage{bbding}
\usepackage{epigraph}
\usepackage{makeidx}
\usepackage{perpage}
\usepackage{tikz}

%-----------------------------------------------------------------------------
% Boondox / mathcalb font declaration
%-----------------------------------------------------------------------------
\DeclareFontFamily{U}{BOONDOX-calo}{\skewchar\font=45 }
\DeclareFontShape{U}{BOONDOX-calo}{m}{n}{
  <-> s*[1.05] BOONDOX-r-calo}{}
\DeclareFontShape{U}{BOONDOX-calo}{b}{n}{
  <-> s*[1.05] BOONDOX-b-calo}{}
\DeclareMathAlphabet{\mathcalb}{U}{BOONDOX-calo}{m}{n}
\SetMathAlphabet{\mathcalb}{bold}{U}{BOONDOX-calo}{b}{n}
\DeclareMathAlphabet{\mathbcalb}{U}{BOONDOX-calo}{b}{n}

%-----------------------------------------------------------------------------
% Algorithm Environment (formerly salgorithm.sty)
%-----------------------------------------------------------------------------
\newcommand{\ProgWordX}[1]{{\small\textbf{#1}}{\ }}
\newcommand{\Do}{\ProgWordX{do}\xspace}
\newcommand{\End}{\ProgWordX{end}\xspace}
\newcommand{\Return}{\ProgWordX{return}\xspace}%
\newcommand{\Procbegin}{\ProgWordX{begin}}
\newcommand{\If}{\ProgWordX{if}\xspace}
\newcommand{\Then}{\ProgWordX{then}\xspace}
\newcommand{\Else}{\ProgWordX{else}\xspace}
\newcommand{\For}{\ProgWordX{for}\xspace}
\newcommand{\While}{{\small\bfseries while}\ }

\newcommand{\xbeginlgox}{\begin{minipage}{1in}\begin{tabbing}
           \quad\=\qquad\=\qquad\=\qquad\=\qquad\=\qquad\=\qquad\=\kill}
\newcommand{\xendlgox}{\end{tabbing}\end{minipage}}

\newenvironment{algorithmEnv}{%
   \begin{tabular}{|l|}\hline\xbeginlgox\bigskip}%
   {\xendlgox\\\hline\end{tabular}}

\newenvironment{program}{%
   \begin{minipage}{4.0in}%
   \begin{tabbing}%
       \ \ \ \ \=\ \ \ \=\ \ \ \ \=\ \ \ \ \=\ \ \ \ \=%
      \ \ \ \ \=\ \ \ \ \=\ \ \ \ \=\ \ \ \ \=%
      \ \ \ \ \=\ \ \ \ \=\ \ \ \ \=\ \ \ \ \=\kill%
}{%
   \end{tabbing}%
   \end{minipage}%
}

\newcommand{\AlgorithmBox}[1]{%
   \fbox{\begin{program} #1 \end{program}}%
}

\definecolor{AlgorithmColor}{cmyk}{0.07,0.90,0,0.34}
\providecommand{\Algorithm}[1]{{\AlgorithmI{#1}\index{algorithm!#1@{\AlgorithmI{#1}}}}}
\newcommand{\CodeComment}[1]{\textcolor{blue}{\texttt{#1}}}

%-----------------------------------------------------------------------------
% Path & Directory Resolution
%-----------------------------------------------------------------------------
\newcommand{\RootDIR}{../}
\IfFileExists{../../../book.tex}{\renewcommand{\RootDIR}{../../../}}{}
\IfFileExists{../../book.tex}{\renewcommand{\RootDIR}{../../}}{}
\IfFileExists{../book.tex}{\renewcommand{\RootDIR}{../}}{}
\IfFileExists{./book.tex}{\renewcommand{\RootDIR}{}}{}

\providecommand{\File}[1]{#1}
\providecommand{\RootFile}[1]{\RootDIR #1}

%-----------------------------------------------------------------------------
% QED Duck Symbol
%-----------------------------------------------------------------------------
\newlength{\CharHeight}
\AtBeginDocument{\setlength{\CharHeight}{\fontcharht\font`X}}

\IfFileExists{styles/qed_duck.pdf}{%
   \providecommand{\myqedsymbol}{%
      \reflectbox{\includegraphics[height=1.4\CharHeight]{styles/qed_duck.pdf}}%
   }%
}{%
   \providecommand{\myqedsymbol}{$\blacksquare$}%
}

%-----------------------------------------------------------------------------
% BibLaTeX Setup (Pure Biber, No BibTeX)
%-----------------------------------------------------------------------------
\usepackage[bibencoding=utf8,style=alphabetic,refsegment=chapter,
indexing=cite,
backend=biber]{biblatex}

\ExecuteBibliographyOptions{%
   maxnames = 40,%
   maxalphanames = 4,%
   maxcitenames = 3,%
   minalphanames = 3,%
   minnames = 1,%
   giveninits = true,%
   doi = false,%
   url = false,%
   isbn = false%
}

\renewbibmacro{in:}{}

\AtEveryBibitem{%
   \ifentrytype{inproceedings}{%
      \clearname{editor}%
      \clearlist{editor}%
      \clearfield{editor}%
      \clearlist{publisher}%
      \clearfield{series}%
      \clearlist{location}%
      \clearfield{publisher}%
   }{}%
   \ifentrytype{book}{%
      \clearname{publisher}%
   }{}%
}
\AtEveryCitekey{%
   \ifentrytype{inproceedings}{%
      \clearlist{publisher}%
   }{}%
   \clearname{editor}%
   \clearlist{location}%
}
\DeclareFieldFormat[inproceedings]{date}{}
\DeclareFieldFormat[inproceedings]{pages}{#1, \printfield{year}}
\DeclareFieldFormat[inproceedings]{booktitle}{\mkbibemph{#1,\nopunct}}

\DeclareFieldFormat{journaltitle}{\mkbibemph{#1}\isdot}

\DeclareFieldFormat[article]{number}{\mkbibparens{#1}}
\DeclareFieldFormat[article]{date}{}
\DeclareFieldFormat[article]{pages}{#1, \printfield{year}}
\DeclareFieldFormat[article]{issue}{\nopunct#1}
\renewbibmacro*{issue+date}{%
   \iffieldundef{issue}{}%
   {%
      \printtext{\unspace}%
      (\printfield{issue})%
   }%
   \iffieldundef{pages}%
   {\addcomma\printfield{year}}%
   {}%
}

\renewcommand{\bibpagespunct}{%
   \ifentrytype{article}{\addcolon}{\addcomma}\space%
}

\DeclareFieldFormat[article]{titlecase}{\MakeSentenceCase{#1}}
\DeclareFieldFormat[inproceedings]{titlecase}{\MakeSentenceCase{#1}}
\DeclareFieldFormat[incollection]{titlecase}{\MakeSentenceCase{#1}}
\DeclareFieldFormat[book]{titlecase}{\MakeSentenceCase{#1}}

\renewbibmacro*{journal}{%
   \iffieldundef{journaltitle} {} {\printtext[journaltitle]{%
         \printfield[myplain]{journaltitle}%
         \setunit{\subtitlepunct}%
         \printfield[myplain]{journalsubtitle}}}}
\DeclareFieldFormat{myplain}{#1}

\renewbibmacro*{booktitle}{%
   \ifboolexpr{ test {\iffieldundef{booktitle}} and test
      {\iffieldundef{booksubtitle}} } {} {\printtext[booktitle]{%
         \printfield[myplain]{booktitle}%
         \setunit{\subtitlepunct}%
         \printfield[myplain]{booksubtitle}}%
      \newunit}%
   \printfield{booktitleaddon}}

\renewbibmacro*{volume+number+eid}{%
   \setunit*{\addcomma\addspace}%
   \printfield{volume}%
   \printfield{number}%
   \setunit{\addcomma\space}%
   \printfield{eid}%
}

\newbibmacro{string+doi}[1]{%
   \iffieldundef{doi}{%
      \iffieldundef{url}{%
         #1%
      }{%
         \href{\thefield{url}}{#1}%
      }%
   }{%
      \href{http://dx.doi.org/\thefield{doi}}{#1}}%
}
\DeclareFieldFormat*{title}{\usebibmacro{string+doi}{\mkbibemph{#1}}}
\renewcommand*{\multicitedelim}{\addcomma\space}
\DeclareFieldFormat[inproceedings]{title}{\usebibmacro{string+doi}{{#1}}}
\DeclareFieldFormat[article]{title}{\usebibmacro{string+doi}{{#1}}}
\DeclareFieldFormat[incollection]{title}{\usebibmacro{string+doi}{{#1}}}
\DeclareFieldFormat[book]{titlecase}{{#1}}

\DeclareCiteCommand{\tcite}%
{\usebibmacro{prenote}}%
{\usebibmacro{citeindex}\usebibmacro{cite}}%
{\multicitedelim}%
{\usebibmacro{postnote}}

%-----------------------------------------------------------------------------
% Hyperref & Visual Formatting
%-----------------------------------------------------------------------------
\usepackage{hyperref}
\hypersetup{%
   breaklinks,%
   colorlinks=true,%
   linkcolor=[rgb]{0.45,0.0,0.0},%
   citecolor=[rgb]{0,0,0.45}%
}

\definecolor{DeepNavy}{HTML}{002366}
\definecolor{RoyalBlue}{HTML}{004E92}
\definecolor{SteelBlue}{HTML}{4682B4}
\definecolor{TealBlue}{HTML}{008080}
\definecolor{DarkSlate}{HTML}{2F4F4F}
\definecolor{DustyGrey}{HTML}{696969}

\makeatletter
\AtBeginDocument{\let\@elt\relax}
\makeatother

\headstyles{bringhurst}
\chapterstyle{dash}

\renewcommand{\chapnamefont}{\normalfont\huge\bfseries\color{DeepNavy}}
\renewcommand{\chapnumfont}{\normalfont\huge\bfseries\color{DeepNavy}}
\renewcommand{\chaptitlefont}{\normalfont\Huge\bfseries\color{DeepNavy}}

\counterwithin{footnote}{chapter}

\setsubparaheadstyle{\normalfont\normalsize\itshape\color{DustyGrey}}
\setlength{\beforechapskip}{-20pt} 
\setlength{\afterchapskip}{30pt}
\setsecnumformat{\csname the#1\endcsname.\quad}

%-----------------------------------------------------------------------------
% Standalone Chapter vs Master Book Modes
%-----------------------------------------------------------------------------
\newcommand{\Copyright}{This work is licensed under the Creative Commons Attribution-Noncommercial 3.0 License. To view a copy of this license, visit \url{http://creativecommons.org/licenses/by-nc/3.0/} or send a letter to Creative Commons, 171 Second Street, Suite 300, San Francisco, California, 94105, USA.}

\newcommand{\RealChapterEndInner}{%
   {%
      \begingroup%
      \emergencystretch=1em%
      \printbibliography[heading=subbibliography,segment=\therefsegment]{}%
      \endgroup%
   }%
}

\providecommand{\BookMode}[1]{}
\providecommand{\StandAloneMode}[1]{#1}

\ifx\bookMode\undefined
   % STANDALONE CHAPTER MODE
   \renewcommand{\StandAloneMode}[1]{#1}
   \renewcommand{\BookMode}[1]{}

   \newcommand{\ChapterNumPage}[2]{%
      \if\relax\detokenize{#1}\relax
         \setcounter{chapter}{0}%
      \else\ifnum#1>0
         \setcounter{chapter}{\numexpr#1-1\relax}%
      \else
         \setcounter{chapter}{0}%
      \fi\fi
      \if\relax\detokenize{#2}\relax
         \setcounter{page}{1}%
      \else\ifnum#2>0
         \setcounter{page}{#2}%
      \else
         \setcounter{page}{1}%
      \fi\fi
   }

   \DeclareRobustCommand{\setheadingcolor}{\color{RoyalBlue}}
   \setsecheadstyle{\setheadingcolor\normalfont\Large\bfseries}
   \setsubsecheadstyle{\normalfont\large\bfseries\color{SteelBlue}}
   \setsubsubsecheadstyle{\normalfont\normalsize\bfseries\color{TealBlue}}
   \setparaheadstyle{\normalfont\normalsize\bfseries\color{DarkSlate}}

   \DeclareDocumentCommand{\Chapter}{s o m}{%
      \IfBooleanTF{#1}{%
         \IfValueTF{#2}{\chapter*[]{#2}}{\chapter*[]{#3}}%
      }{%
         \IfValueTF{#2}{\chapter*[]{#2}}{\chapter[]{#3}}%
      }%
      \setcounter{footnote}{0}%
   }
   \newcommand{\ChapterPrefix}{}
   \newcommand{\ChapterGen}{\pagestyle{plain}\vspace*{-0.7cm}}
   \newcommand{\ChapterEnd}{%
      \RealChapterEndInner{}%
      \end{document}%
   }

   \renewcommand{\File}[1]{#1}
   \renewcommand{\RootFile}[1]{../#1}
\else
   % MASTER BOOK MODE
   \renewcommand{\StandAloneMode}[1]{}
   \renewcommand{\BookMode}[1]{#1}

   \newcommand{\ChapterNumPage}[2]{}

   \DeclareDocumentCommand{\Chapter}{s o m}{%
      \IfBooleanTF{#1}{%
         \IfValueTF{#2}{\chapter*[#2]{#3}}{\chapter*{#3}}%
         \ifx\Dir\undefined\else
            \typeout{CHAPTER_MAP: \Dir/\CurrentChapFile => chap: 0, page: \thepage}%
         \fi
      }{%
         \IfValueTF{#2}{\chapter[#2]{#3}}{\chapter{#3}}%
         \ifx\Dir\undefined\else
            \typeout{CHAPTER_MAP: \Dir/\CurrentChapFile => chap: \thechapter, page: \thepage}%
         \fi
      }%
   }
   \newcommand{\ChapterPrefix}{%
      \vspace{-1cm}%
      \subsection*{\small
         {CS 574: Randomized Algorithms}\\
         {Sariel Har-Peled}\\
         {\today}}%
   } \newcommand{\ChapterGen}{}
   \newcommand{\ChapterEnd}{\RealChapterEndInner{}}

   \renewcommand{\File}[1]{\Dir/#1}
   \renewcommand{\RootFile}[1]{#1}

   \newcommand{\InputChap}[2]{%
      \def\Dir{#1}%
      \def\CurrentChapFile{#2}%
      \typeout{}%
      \typeout{------------------------------------------------------------}%
      \typeout{PROCESSING FILE: #1/#2}%
      \typeout{}%
      \input{#1/#2}%
   }
\fi

%-----------------------------------------------------------------------------
% Fragments & Quotes
%-----------------------------------------------------------------------------
\newcommand{\InputExt}[1]{%
   \IfFileExists{#1}%
   {\input{#1}}%
   {\IfFileExists{../#1}%
      {\input{../#1}}%
      {\IfFileExists{../../#1}%
         {\input{../../#1}}%
         {\typeout{ERROR: Unable to find #1}}%
      }%
   }%
}
\makeatletter
\def\FragNormKey#1{%
  \frag@norm@check#1def_\relax{#1}%
}
\def\frag@norm@check#1def_#2\relax#3{%
  \if\relax\detokenize{#1}\relax
    #3%
  \else
    def_#3%
  \fi
}
\makeatother

\newcommand{\FragmentName}[1]{\RootFile{fragment/\FragNormKey{#1}.tex}}

\newenvironment{fragment}[1]{}{}

\makeatletter
\AtBeginDocument{%
   \let\real@latex@label\label
   \def\label#1{%
      \global\@namedef{all@labels@seen@#1}{1}%
      \real@latex@label{#1}%
   }%
}

\DeclareDocumentCommand{\IncFragment}{m}{%
   \begingroup
      \def\label##1{%
         \@ifundefined{all@labels@seen@##1}{%
            \global\@namedef{all@labels@seen@##1}{1}%
            \real@latex@label{##1}%
         }{}%
      }%
      \IfFileExists{\FragmentName{#1}}%
         {\InputExt{\FragmentName{#1}}}%
         {\InputExt{\RootFile{fragment/#1.tex}}}%
   \endgroup
}
\makeatother

\newcommand{\Quote}[2]{%
   \SaveIndent
   \begin{center}%
      {\footnotesize
         \begin{quote}
             \RestoreIndent
             {#1}\\
             \hspace*{2cm}-- {#2.}
         \end{quote}%
      }%
      \index{quotation!#2}%
   \end{center}%
   \bigskip
}

\newcommand{\QuotePExt}[3]{%
   \bigskip%
   \hfill%
   \begin{minipage}{0.95\linewidth}%
       \footnotesize {#1}%

       \rule{\linewidth}{0.3mm}%
       \hfill #2, #3%
   \end{minipage}%
   \bigskip%
}

\epigraphfontsize{\small\itshape}
\setlength{\epigraphwidth}{0.6\textwidth}
\newcommand{\QuoteExt}[3]{%
   \vspace*{-1.6cm}%
   \epigraph{{\footnotesize #1}}{{#2}, {#3}}%
   \vspace*{-1.2cm}%
}

\providecommand{\delX}[1]{{\color{red}\sout{#1}}}
\providecommand{\newX}[1]{{\color{blue}#1}}
\providecommand{\chgY}[2]{{\color{red}\sout{#1}}{\color{blue}#2}}
\providecommand{\remX}[1]{{\color{purple}\textbf{[Remark: #1]}}}
\newcommand{\SpellIgnore}[1]{#1}
\newcommand{\EQ}{\mathrel{\;\;=\;\;}}
\newcommand{\LEQ}{\mathrel{\;\;\leq\;\;}}
\newcommand{\LT}{\mathrel{\;\;<\;\;}}
\newcommand{\GT}{\mathrel{\;\;<\;\;}}

\newcommand{\GEQ}{\mathrel{\;\;\geq\;\;}}

%-----------------------------------------------------------------------------
% Save a symbol that we know is going to get redefined.
\def\savesymbol#1{%
  \expandafter\let\expandafter\origsym\expandafter=\csname#1\endcsname
  \expandafter\let\csname orig#1\endcsname=\origsym
  \expandafter\let\csname#1\endcsname=\relax
}

% Restore a previously saved symbol, and rename the current one.
\def\restoresymbol#1#2{%
  \expandafter\let\expandafter\newsym\expandafter=\csname#2\endcsname
  \expandafter\global\expandafter\let\csname#1#2\endcsname=\newsym
  \expandafter\let\expandafter\origsym\expandafter=\csname orig#2\endcsname
  \expandafter\global\expandafter\let\csname#2\endcsname=\origsym
}

\usepackage{pictex}%
\savesymbol{triangle}%
\usepackage{tree}%
\restoresymbol{TreePKG}{triangle}%

% Definitions absorbed from sbook.sty
\providecommand{\MdPrc}{\nu}

\providecommand{\TwoFiguresBase}[4]{%
   \begin{tabular}{cc}
      \begin{minipage}{0.47\linewidth}
          \begin{center}
              {{{#1}}}
          \end{center}
      \end{minipage}
      &
      \begin{minipage}{0.47\linewidth}
          \begin{center}
              {{#3}}
          \end{center}
      \end{minipage}
      \\
      {#2} & {#4}
   \end{tabular}%
}

\providecommand{\TwoFigures}[6]{%
   \begin{tabular}{cc}
      \begin{minipage}{0.48\linewidth}
          \begin{center}
              {\IncludeGraphics[#1]{#2}}
          \end{center}
      \end{minipage}
      &
      \begin{minipage}{0.48\linewidth}
          \centerline{\IncludeGraphics[#4]{#5}}
      \end{minipage}
      \\
      {#3} & {#6}
   \end{tabular}%
}

\providecommand{\ThreeFigures}[6]{%
   \begin{tabular}{ccc}
        \begin{minipage}{0.30\linewidth}
            \IncludeGraphics[width=5.1cm]{#1}\\
        \end{minipage}
        &
        \begin{minipage}{0.30\linewidth}
            \IncludeGraphics[width=5.1cm]{#2}\\
        \end{minipage}
        &
        \begin{minipage}{0.30\linewidth}
            \IncludeGraphics[width=5.1cm]{#3}\\
        \end{minipage}
        \vspace{-1cm}
        \\
        {#4} & {#5}& {#6}
   \end{tabular}%
}

%-----------------------------------------------------------------------------
\providecommand{\remove}[1]{}
\providecommand{\NotAMSMode}[1]{#1}
\providecommand{\Solution}[1]{}
\providecommand{\Version}[1]{\hfill{Version: #1}}
\providecommand{\Distor}{\mathrm{dist}}
\providecommand{\Hardness}[1]{(h:#1)}
\providecommand{\QuoteOptExt}[3]{\QuoteExt{#1}{#2}{#3}}
\providecommand{\Hypercube}{\EuScript{H}}
\providecommand{\CoverTime}{\EuScript{C}}
\providecommand{\Nbr}{\Gamma}
\providecommand{\CommuteTime}[2]{\mathbf{CT}_{#1#2}}
\providecommand{\Resistance}{\mathbf{R}}
\providecommand{\EffectResist}[2]{\mathbf{R}_{#1#2}}
\providecommand{\NetworkE}{\EuScript{N}}
\providecommand{\trinom}[4]{\binom{#1}{#2 \;\; #3 \;\; #4}}
\providecommand{\Ranges}{\EuScript{R}}
\providecommand{\CClass}[1]{\ensuremath{\mathbf{#1}}}
\providecommand{\coCClass}[1]{\ensuremath{\mathrm{co-}}\CClass{#1}}
\providecommand{\CX}{\EuScript{C}}
\providecommand{\packet}{\mathbf{c}}
\providecommand{\packetA}{\mathbf{d}}
\providecommand{\mopt}{m_{\mathrm{opt}}}
\providecommand{\girth}{\mathop{\mathrm{girth}}}
\providecommand{\CR}{\mathop{\mathrm{cr}}}
\providecommand{\TRUE}{{\textcolor[named]{RedViolet}{\textsf{TRUE}}}\xspace}
\providecommand{\FALSE}{{\textcolor[named]{RedViolet}{\textsf{FALSE}}}\xspace}
\providecommand{\Part}[1]{\part{#1}}
\providecommand{\Section}[1]{\section{#1}}
\providecommand{\SubSection}[1]{\subsection{#1}}
\providecommand{\SubSubSection}[1]{\subsubsection{#1}}
\providecommand{\Exercise}[5]{%
\begin{exercise_h}[#1]%
   \exelab{#3}{}%
   #4%
\end{exercise_h}%
}
\providecommand{\si}[1]{#1}
\providecommand{\myParagraph}[1]{\noindent\textbf{#1}}
\providecommand{\DefName}[1]{#1}
\providecommand{\deflabel}[1]{#1}
\providecommand{\Description}[1]{}

%-----------------------------------------------------------------------------
% Mathematical Notations & Environments (from notations.sty)
%-----------------------------------------------------------------------------
%-----------------------------------------------------------------------------
% List environments (questions & compactenum)
%-----------------------------------------------------------------------------

\newlist{compactenumA}{enumerate}{5}%
\setlist[compactenumA]{topsep=0pt,itemsep=-1ex,partopsep=1ex,parsep=1ex,%
   label=(\Alph*)}%

\newlist{compactenumI}{enumerate}{5}%
\setlist[compactenumI]{topsep=0pt,itemsep=-1ex,partopsep=1ex,parsep=1ex,%
   label=(\Roman*)}%

\newlist{compactenuma}{enumerate}{5}%
\setlist[compactenuma]{topsep=0pt,itemsep=-1ex,partopsep=1ex,parsep=1ex,%
   label=(\alph*)}%

\newlist{compactitem}{enumerate}{5}%
\setlist[compactitem]{topsep=0pt,itemsep=-1ex,partopsep=1ex,parsep=1ex,%
   label=\textbullet}%

\newlist{compactenumi}{enumerate}{5}%
\setlist[compactenumi]{topsep=0pt,itemsep=-1ex,partopsep=1ex,parsep=1ex,%
   label=(\roman*)}%

\setcounter{secnumdepth}{5}

\newcommand{\HST}{\TermI{HST}\xspace}
\newcommand{\Term}[1]{\textsf{#1}}
\newcommand{\TermI}[1]{\Term{#1}\index{#1@\Term{#1}}}

\DeclareMathAlphabet{\mathpzc}{OT1}{pzc}{m}{it}

\newcommand{\seg}{\mathsf{s}}

\newcommand{\Diam}{\mathrm{diam}}
\newcommand{\diam}{\mathrm{diam}}

\newcommand{\Copt}{C_{\mathrm{opt}}}
\newcommand{\CenterPnt}{\mathrm{cen}}

\newcommand{\parent}{\overline{\mathrm{p}}}
\newcommand{\parentX}[2][\!]{\parent\pth[#1]{#2}}

\newcommand{\FMS}{\EuScript{X}\index{metric space}}  
\newcommand{\FMSB}{\EuScript{Y}\index{metric space}}  
\newcommand{\FastCut}{\Algorithm{FastCut}\xspace}
\newcommand{\MinCut}{\Algorithm{MinCut}\xspace}
\newcommand{\MinCutRep}{\Algorithm{MinCutRep}\xspace}
\newcommand{\Contract}{\Algorithm{Contract}\xspace}

\newcommand{\Grid}{\mathrm{G}}
\newcommand{\ClosestPair}{\mathcal{CP}}
\newcommand{\SphereChar}{\ensuremath{\mathbb{S}}}
\newcommand{\Sphere}[1]{{\SphereChar^{(#1)}}}

\newcommand{\PntSet}{{\mathsf{P}}}%
\providecommand{\PntSetA}{{\mathsf{Q}}}%

\newcommand{\query}{\mathcalb{q}}

\newcommand{\pnt}{\mathsf{p}}
\newcommand{\pntA}{\mathsf{u}}
\newcommand{\pntB}{\mathsf{v}}
\newcommand{\pntC}{{\mathsf{x}}}
\newcommand{\pntD}{{\mathsf{y}}}
\newcommand{\pntE}{{\mathsf{z}}}

\newcommand{\pntCoord}{p}
\newcommand{\pntCoordA}{q}

%%%%%%%%%%%%%%%%%%%%%%%%%%%%%%%%%%%%%%%%%%%%%%%%%%%%%%%%%%%%%%%%%%%%% 
%%%%%%%%%%%%%%%%%%%%%%%%%%%%%%%%%%%%%%%%%%%%%%%%%%%%%%%%%%%%%%%%%%%%%
% Markov chain notations.

\newcommand{\MCSpc}{\mathtt{S}}
\newcommand{\MCPrb}{\mathtt{P}}
\newcommand{\MCPrbY}[2]{\MCPrb_{#1#2}}
\newcommand{\MCPrbM}{\mathbf{P}}%
\newcommand{\MCPrbT}[3]{\MCPrb_{#1#2}^{(#3)}}
\newcommand{\MCFPrbT}[3]{\mathtt{r}_{#1#2}^{(#3)}}
\newcommand{\MCVstProb}[2]{\mathtt{f}_{#1#2}}
\newcommand{\MCExSteps}[2]{\mathtt{h}_{#1#2}}
\newcommand{\MCPVec}[1]{\mathbf{q}^{(#1)}}
%%%%%%%%%%%%%%%%%%%%%%%%%%%%%%%%%%%%%%%%%%%%%%%%%%%%%%%%%%%%%%%%%%%%%

\newcommand{\Spread}{\Phi}
\newcommand{\VC}{\ensuremath{\mathsf{VC}}\xspace}
\newcommand{\MC}{\ensuremath{\mathsf{MC}}\xspace}
\newcommand{\hyperplane}{h}

\newcommand{\QuickSort}{\Algorithm{QuickSort}\xspace}
\newcommand{\Alg}{\Algorithm{Alg}\xspace}

\newcommand{\Graph}{\mathsf{G}}%
\newcommand{\GraphB}{\mathsf{H}}
\newcommand{\GraphC}{\mathsf{K}}
\newcommand{\GraphD}{\mathsf{F}}

\providecommand{\G}{}
\renewcommand{\G}{\Graph}

\newcommand{\RLP}{\ComplexityClass{RLP}\xspace}
\newcommand{\USTCON}{\ComplexityClass{USTCON}}
\newcommand{\STCON}{\overrightarrow{\ComplexityClass{STCON}}}

\newcommand{\Mtrx}{\mathsf{M}}
\newcommand{\TMtrx}{\mathsf{Q}}
\newcommand{\IMtrx}{\mathcal{I}}

\newcommand{\Expander}[3]{$\left[#1,#2,#3\right]$-expander}

\newcommand{\ChartX}[1]{\ensuremath{\chi_{#1}}}
\newcommand{\Field}{\mathbb{F}}
\newcommand{\LSpace}{\mathbb{G}}
\newcommand{\LD}[2]{\mathsf{L{D}}\pth{#1,#2}}
\newcommand{\SEVN}{\widehat{\lambda_2}}
\newcommand{\NX}[2]{#1\!\left[#2\right]}
\newcommand{\Entry}[2]{#1\!\left[#2\right]}
\newcommand{\NXExt}[3]{#1_{#2}\!\left[#3\right]}
\newcommand{\RangeI}[1]{\llbracket #1 \rrbracket}
\newcommand{\DialX}[1]{\gamma\pth{#1}}
\newcommand{\BiDialX}[1]{{\gamma}_2\pth{#1}}

\newcommand{\EV}[1]{\widehat{\lambda_{#1}}}
\newcommand{\TEV}[1]{\widehat{\lambda_{#1}}}
\newcommand{\EVal}{\widehat{\lambda}}
\newcommand{\Basis}{\mathcal{B}}

\newcommand{\ZProd}[2]{\ensuremath{#1 \, \text{\textcircled{z}} #2}}
\newcommand{\RProd}[2]{\ensuremath{#1 \, \text{\textcircled{r}} #2}}

\newcommand{\OneVec}[1]{\mathbf{1}^{#1}}
\newcommand{\DstrbV}[1]{\mathbf{\mathsf{q}}^{(#1)}}
\newcommand{\DstrbE}[2]{\mathsf{q}^{(#1)}_{#2}}

\newcommand{\DEdges}{\widetilde{\Edges}}

\newcommand{\VV}{\mathsf{V}}%
\newcommand{\EE}{\mathsf{E}}%

\newcommand{\Vertices}{\VV}%
\newcommand{\Edges}{\EE}%

\newcommand{\EdgesX}[1]{\Edges\pth{#1}}
\newcommand{\VerticesX}[1]{V\pth{#1}}

\newcommand{\StringX}[1]{\mathsf{s}\pth{#1}}

\newcommand{\BMtrx}{\mathsf{B}}
\newcommand{\WMtrx}{\mathsf{W}}
\newcommand{\NWMtrx}{\overline{\mathsf{W}}}
\newcommand{\Sequence}{\EuScript{S}}

\newcommand{\med}{\mathrm{med}}
\newcommand{\medX}[1]{\mathrm{med}\pth{#1}}

\newcommand{\medianC}{\nu}%

\newlength{\savedparindent}
\newcommand{\SaveIndent}{\setlength{\savedparindent}{\parindent}}
\newcommand{\RestoreIndent}{\setlength{\parindent}{\savedparindent}}

\newcommand{\NNN}{\mathbb{N}}%
\providecommand{\NN}{\mathbb{N}}%
\AtBeginDocument{%
   \renewcommand{\Re}{\mathbb{R}}%
}
%\shogi
\newcommand{\reals}{\Re}%
\newcommand{\ZZ}{\mathbb{Z}}%

\newcommand{\level}{\mathop{\mathrm{level}}}

\providecommand{\TPDF}[2]{\texorpdfstring{#1}{#2}}

\newcommand{\Hint}[1]{\emph{[#1]}}
\newcommand{\Arr}{\mathop{\mathrm{\EuScript{A}}}}

\newcommand{\Array}{\mathcal{T}}%

\newcommand{\dirEdge}[2]{\left({{#1} \rightarrow {#2}} \right)}

%%%%%%%%%%%%%%%%%%%%%%%%%%%%%%%%%%%%%%%%%%%%%%%%%%%%%%%%%%%%%%%%%%%%%%
% Emphasize stuff...
\definecolor{blue25}{rgb}{0, 0, 11}
\newcommand{\emphic}[2]{%
%   \BADVERYBAD%
   \textcolor{blue25}{%
      \bfseries{\emph{#1}}}%
   \index{#2}}
\newcommand{\emphi}[1]{\emphic{#1}{#1}}
\newcommand{\emphOnly}[1]{\textcolor{blue25}{\bfseries{\emph{#1}}}}
\providecommand{\emphiC}[1]{\emphic{#1}{#1}}


\definecolor{almostblack}{rgb}{0, 0, 0.3}
\providecommand{\emphw}[1]{{\textcolor{almostblack}{\emph{#1}}}}%

%
%%%%%%%%%%%%%%%%%%%%%%%%%%%%%%%%%%%%%%%%%%%%%%%%%%%%%%%%%%%%%%%%%%%%%%

\newcommand{\volX}[1]{\mathrm{vol}\pth{#1}}

\newcommand{\Interval}{\EuScript{I}}
\newcommand{\IntervalA}{\EuScript{I}}

\newcommand{\Exp}{\mathrm{Exp}}

\let\Ball\undefined
\newcommand{\Ball}{\mathbf{b}}
\newcommand{\hdim}{d}
\newcommand{\Domain}{\mathcal{D}}
\newcommand{\points}[1]{[\bfseries{#1 points}]}
\newcommand{\Event}{\EuScript{E}}%
\newcommand{\EventB}{\EuScript{B}}%
\newcommand{\EventF}{\EuScript{F}}%
\newcommand{\EventG}{\EuScript{G}}%

\newcommand{\distCmd}[1]{\left\| {#1}  \right\|}
\newcommand{\distY}[2]{\distCmd{#1 - #2}}

\newcommand{\NormalC}{\mathbf{N}}
\newcommand{\NormalY}[2]{\NormalC\pth{#1, #2}}%
\newcommand{\NormalDistOneDim}{\NormalC}
\newcommand{\NormalDistX}[1]{\NormalC^{#1}}
\DeclareMathOperator{\lca}{lca}
\newcommand{\RangeSpace}{\mathsf{S}}%
\newcommand{\RangeSpaceA}{\mathsf{T}}

\newcommand{\range}{\mathbf{r}}
\newcommand{\rangeA}{\mathbf{r}'}

\newcommand{\BadProb}{\varphi}

\newcommand{\EntropyChar}{\mathbb{H}}
\newcommand{\Entropy}[1]{\EntropyChar\pth{#1}}
\newcommand{\EntropyX}[1]{\EntropyChar\pth[]{#1}}

\newcommand{\Matrix}{\mathsf{M}}
\newcommand{\MatrixB}{\mathsf{B}}
\newcommand{\MatrixC}{\mathsf{C}}
\newcommand{\MatrixD}{\mathsf{D}}

\renewcommand{\th}{th\xspace}

\newcommand{\Family}{\mathcal{F}}
\newcommand{\FamilyB}{\mathcal{G}}
\newcommand{\FamilyC}{\mathcal{H}}
\newcommand{\FMSC}{\EuScript{Y}\index{metric space}}

\newcommand{\median}{\overline{\mathsf{m}}}
\newcommand{\QuickSelect}{\Algorithm{Quick{}Select}\xspace}

\providecommand{\row}{\mathbf{r}}%
\newcommand{\rowc}{{r}}%

\newcommand{\vecb}{\mathbf{b}}

\newcommand{\vecA}{\mathbf{v}}
\newcommand{\vecB}{\mathbf{u}}
\newcommand{\vecRand}{\mathbf{r}}

\newcommand{\vecH}{\mathbf{h}}
\newcommand{\vecV}{\mathbf{v}}

\newcommand{\vecX}{\mathbf{x}}
\newcommand{\vecY}{\mathbf{y}}

\newcommand{\DistBinom}[2]{\mathrm{Bin}\pth{#1,#2}}

\newcommand{\DistGeom}[1]{\mathrm{Geom}\pth{#1}}

\newcommand{\Sample}{\mathsf{R}}

\newcommand{\lcm}{\mathop{\mathrm{lcm}}}
\newcommand{\bdiv}{\mathop{\mathrm{div}}}%
\newcommand{\divides}{\mid}%
\newcommand{\poly}{\mathrm{poly}}%

\newcommand{\Group}{\mathcal{G}} 
\newcommand{\GroupA}{\mathcal{H}} 

\newcommand{\algGCD}{\Algorithm{EuclidGCD}\xspace}%

\newcommand{\gOp}{\times}%
\newcommand{\gIdent}{\mathsf{i}}%

\newcommand{\ordX}[2][\!]{\mathop{\mathrm{o{r}d}}\pth[#1]{#2}}

\newcommand{\lSym}[2]{%
   \left(
       \ds {#1}\divides{#2}%
   \right) %
}
\newcommand{\jSym}[2]{\left\llbracket%
       \ds {#1}\divides{#2}%
       \right\rrbracket%
}
\newcommand{\Line}{\ell}

\newcommand{\LevelChar}{\mathbb{L}}
\newcommand{\AtMostKLevelChar}{\LevelChar_{\leq k}}
\newcommand{\AtMostKLevel}[1]{\LevelChar_{\leq k} \pth{ #1}}

\newcommand{\LineSet}{\mathsf{L}}%
\newcommand{\ArrX}[1]{\EuScript{A}\pth{#1}}%

\providecommand{\Matousek}{Matou{\v s}ek\xspace}
\providecommand{\Erdos}{Erd{\H o}s\xspace}
\providecommand{\Lovasz}{Lov\'asz\xspace}

\newcommand{\ProblemC}[1]{{\textsf{\textcolor{red}{#1}\index{problem!#1}}}}
\DefineNamedColor{named}{OliveGreen} {cmyk}{0.64,0,0.95,0.40}
\providecommand{\ComplexityClass}[1]{{{\textcolor[named]{OliveGreen}{%
            \textsc{#1}}}}}
\providecommand{\NPHard}{{\ComplexityClass{NP-Hard}}%
   \index{NP!hard}\xspace}
\providecommand{\NPComplete}{\ComplexityClass{NP-Complete}%
   \index{NP!complete}\xspace}
\providecommand{\APXHard}{\ComplexityClass{APX-Hard}%
   \index{APX!Hard}\xspace}
\providecommand{\NPC}{\ComplexityClass{NPC}%
   \index{NP!complete}\xspace}
\providecommand{\NP}{\ComplexityClass{NP}%
   \index{NP}\xspace}
\providecommand{\POLYT}{\ComplexityClass{P}\xspace}

\newcommand{\envFrame}[1]{\begin{center}\fbox{\begin{minipage}[c]{0.75\textwidth}{#1}\end{minipage}}\end{center}}

\newcommand{\mysqrt}[1]{\ts\!\!\sqrt{#1}}
\newcommand{\SegSet}{\mathsf{S}}
\newcommand{\ArrVD}[1]{\Arr^{|} \pth{#1}}
\newcommand{\RIC}{\textsf{RIC}\index{random incremental
      construction}\xspace}
\newcommand{\ArrB}{\EuScript{B}}

\newcommand{\TreeOpt}{\EuScript{T}_\mathrm{opt}}
\newcommand{\TreeOptA}{\EuScript{T}_\mathrm{opt}'}

\newcommand{\DepthExt}[2]{\mathrm{depth}_{#1}\pth{#2}}
\newcommand{\depthX}[1]{\mathrm{depth}\pth{#1}}

\newcommand{\TreeA}{\EuScript{U}}
\newcommand{\TreeB}{\EuScript{V}}

\newcommand{\Trap}{\sigma}
\newcommand{\TrapA}{\sigma'}
\newcommand{\TrapB}{\sigma''}

\newcommand{\SampleA}{\mathsf{R}'}

\newcommand{\AllRegions}{\mathcal{T}}
\newcommand{\RegionsSet}[1]{\mathcal{F}\pth{#1}}
\newcommand{\RegionsHeavy}[2]{\mathcal{F}_{\!\geq#1}\pth{#2}}
\newcommand{\NumOfRegionsZero}[1]{\ExChar \!f\pth{#1}}
\newcommand{\NumOfRegions}[2]{\ExChar\! f_{\!\geq #1}\pth[]{#2}}
\newcommand{\ObjSet}{\mathsf{S}}

\newcommand{\WeightX}[2][\!]{\omega \pth[#1]{#2}}
\newcommand{\Edge}{\mathsf{e}}

\newcommand{\Face}{\mathsf{f}}

\newcommand{\DefSet}[1]{D\pth{#1}}
\newcommand{\KillSet}[1]{K\pth{#1}}

\newcommand{\CSProbFull}[4]{
   \,\rho_{\!\!\stackrel{\rule[-.005cm]{0cm}{0.1cm}}{#1,
       #4}}\pth{\MakeSBig #2,#3}}

\newcommand{\disk}{D}
\newcommand{\diskA}{D'}

\newcommand{\IntervalSet}{\widehat{\mathcal{I}}}

\newcommand{\GroundSet}{\textsf{X}}
\newcommand{\FGroundSet}{\mathsf{x}}
\newcommand{\RangeSet}{{\mathcal{R}}}
\newcommand{\RangeSetA}{\mathcal{R}'}
\newcommand{\MeasureChar}{\overline{m}}
\newcommand{\Measure}[1]{\MeasureChar\pth{#1}}
\newcommand{\sMeasure}[1]{\overline{s}\pth{#1}}

\newcommand{\SetA}{S}%
\newcommand{\SetB}{T}%
\newcommand{\SetC}{N}%

\newcommand{\VCProj}[2]{#1_{|#2}}
\newcommand{\VCDim}{\mathrm{dim}_{\VC}}
\newcommand{\Dim}{\delta}
\newcommand{\DimA}{\delta'}
\newcommand{\PntSetB}{{{\mathsf{Q}}}}
\newcommand{\PntSetC}{{{\mathsf{R}}}}
\newcommand{\Growth}[2]{{\EuScript{G}}_{#1}\pth{#2}}
\newcommand{\ShatterFunc}[1]{\pi_{#1}}
\newcommand{\Distribution}{\EuScript{D}}
\newcommand{\Error}{\EuScript{E}}

\newcommand{\bd}{{\partial}}%
\newcommand{\eps}{{\varepsilon}}%

\newcommand{\clmlab}[1]{\label{claim:#1}}
\newcommand{\clmref}[1]{\HLink{Claim}{claim:#1}}

\newcommand{\CHX}[1]{{\mathcal{CH}}\pth{#1}}

\newcommand{\ts}{\hspace{0.6pt}}

\def\qedsymbol{\myqedsymbol}%

\newcommand{\qedsign}{\hfill{\hfill\myqedsymbol}}
\newcommand{\qedsignBig}{\hfill{\hfill\rule{4mm}{4mm}}}
\newcommand{\QED}{{\hfill{\hfill\qedsymbol}}}

\newcommand{\num}{\alpha}

\newcommand{\NEW}[1]{#1}

\newcommand{\ThreeFiguresExt}[9]{%
   \begin{tabular}{ccc}
     \begin{minipage}{0.31\linewidth}
         \begin{center}
             {\IncludeGraphics[#1]{#2}}
         \end{center}
     \end{minipage}
     &
       \begin{minipage}{0.31\linewidth}
           \centerline{\IncludeGraphics[#4]{#5}}
       \end{minipage}
     &
       \begin{minipage}{0.31\linewidth}
           \centerline{\IncludeGraphics[#7]{#8}}
       \end{minipage}
     \\
     ~\\[-0.2cm]
     {#3} & {#6} & {#9}
   \end{tabular}}

\newcommand{\IncludeGraphics}[2][]{%
   \typeout{}%
   \typeout{Graphics: #2}%
   \typeout{\ includegraphics[#1]{#2}}%
   \includegraphics[#1]{#2}
   \typeout{}%
}

%    -----------------------------------------------------------------
%    -----------------------------------------------------------------
%    -----------------------------------------------------------------
%    Handling references
%    -----------------------------------------------------------------
%    -----------------------------------------------------------------
%    -----------------------------------------------------------------
\newcommand{\HLinkShort}[2]{\hyperref[#2]{#1\ref*{#2}}}
\newcommand{\HLink}[2]{\hyperref[#2]{#1~\ref*{#2}}}
\newcommand{\HLinkPage}[2]{\hyperref[#2]{#1~\ref*{#2}%
      $_\text{p\pageref{#2}}$}}

\newcommand{\HLinkSuffix}[3]{\hyperref[#2]{#1\ref*{#2}{#3}}}
\newcommand{\HLinkPageSuffix}[3]{\hyperref[#2]{#1\ref*{#2}%
      #3$_\text{p\pageref{#2}}$}}

\newcommand{\figlab}[1]{\label{fig:#1}}
\newcommand{\figref}[1]{\HLink{Figure}{fig:#1}}

\newcommand{\exelab}[1]{\label{exer:#1}}
\newcommand{\exeref}[1]{\HLink{Exercise}{exer:#1}}

\providecommand{\deflab}[1]{\label{def:#1}}
\newcommand{\defref}[1]{\HLink{Definition}{def:#1}}
\newcommand{\defrefY}[2]{\hyperref[def:#2]{#1}}
\newcommand{\defrefpage}[1]{\HLinkPage{Definition}{def:#1}}

\newcommand{\etal}{\textit{et~al.}\xspace}

\newcommand{\obslab}[1]{\label{observation:#1}}
\newcommand{\obsref}[1]{\HLink{Observation}{observation:#1}}
\newcommand{\obsrefY}[2]{\hyperref[observation:#2]{#1}}

\newcommand{\exmlab}[1]{\label{example:#1}}
\newcommand{\exmref}[1]{\HLink{Example}{example:#1}}

\newcommand{\remlab}[1]{\label{rem:#1}}
\newcommand{\remref}[1]{\HLink{Remark}{rem:#1}}%

\newcommand{\chaplab}[1]{\label{chap:#1}}

\newcommand{\itemlab}[1]{\label{item:#1}}
\newcommand{\itemref}[1]{\HLinkSuffix{}{item:#1}{}}

\newcommand{\tedlab}[1]{\label{tedium:#1}}
\newcommand{\tedref}[1]{\HLink{Tedium}{tedium:#1}}%

\newcommand{\factlab}[1]{\label{fact:#1}}%

\newcommand{\alglab}[1]{\label{algorithm:#1}}

\newcommand{\net}{\mathcal{N}}%


\newcommand{\lemlab}[1]{\label{lemma:#1}}
\newcommand{\lemref}[1]{\HLink{Lemma}{lemma:#1}}%
\newcommand{\lemrefY}[2]{\hyperref[lemma:#2]{#1}}
\newcommand{\lemrefpage}[1]{\HLinkPage{Lemma}{lemma:#1}}
\newcommand{\lemrefshort}[1]{\HLinkShort{L}{lemma:#1}}

\newcommand{\problab}[1]{\label{prob:#1}}%
\newcommand{\probref}[1]{\HLink{Problem}{prob:#1}}%

\newcommand{\eqlab}[1]{\label{equation:#1}}
\newcommand{\Eqrefbare}[1]{\HLinkSuffix{(}{equation:#1}{)}}
\newcommand{\Eqref}[1]{\HLinkSuffix{Eq.~(}{equation:#1}{)}}
\newcommand{\Eqrefpage}[1]{\HLinkPageSuffix{Eq.~(}{equation:#1}{)}}

\newcommand{\seclab}[1]{\label{sec:#1}}
\newcommand{\secref}[1]{\HLink{Section}{sec:#1}}
\newcommand{\secrefY}[2]{\hyperref[sec:#2]{#1}}
\newcommand{\secrefpage}[1]{\HLinkPage{Section}{sec:#1}}

\newcommand{\thmlab}[1]{{\label{theo:#1}}}
\newcommand{\thmref}[1]{\HLink{Theorem}{theo:#1}}
\newcommand{\thmrefY}[2]{\hyperref[theo:#2]{#1}}
\newcommand{\thmrefpage}[1]{\HLinkPage{Theorem}{theo:#1}}

\newcommand{\corlab}[1]{\label{cor:#1}}
\newcommand{\corref}[1]{\HLink{Corollary}{cor:#1}}%
\newcommand{\correfY}[2]{\hyperref[cor:#2]{#1}}

\newcommand{\Yend}{Y_{\mathrm{end}}}%


%    -----------------------------------------------------------------
%    -----------------------------------------------------------------
%    -----------------------------------------------------------------
%    Handling references
%    -----------------------------------------------------------------
%    -----------------------------------------------------------------
%    -----------------------------------------------------------------

\newcommand{\ceil}[1]{\left\lceil {#1} \right\rceil}
\newcommand{\floor}[1]{\left\lfloor {#1} \right\rfloor}
\newcommand{\brc}[1]{\left\{ {#1} \right\}}
\newcommand{\pth}[2][\!]{\mleft({#2}\mright)}%
\newcommand{\brk}[1]{\mleft[{#1}\mright]}%
\newcommand{\pbrc}[2][\!\!]{#1\left[ {#2} \bigr. \right]}
\newcommand{\pbrcS}[1]{\left[ {#1} \right]}

\newcommand{\dotprod}[2]{{#1}\cdot{#2}}
\newcommand{\DotProd}[2]{\permut{{#1},{#2}}}
\newcommand{\permut}[1]{\left\langle {#1} \right\rangle}

\newcommand{\MakeBig}{}
\newcommand{\MakesBig}{}
\newcommand{\MakeSBig}{}
\newcommand{\MakeVBig}{}

\newcommand{\DistChar}{\mathsf{d}}%
\newcommand{\DistSetY}[2]{\DistChar\pth{#1,#2}}
\newcommand{\distG}[2]{\DistChar_\Graph\pth{#1, #2}}
\newcommand{\DistMtrCharX}[1]{\DistChar_{#1}}

\newcommand{\dY}[2]{\DistChar\pth{#1, #2}}%
\newcommand{\dGY}[2]{\DistChar_\Graph\pth{#1, #2}}%

\newcommand{\cset}{\mathcal{C}}%

\newcommand{\MTR}{\mathcal{M}}
\newcommand{\DistMtr}[2]{\DistChar\pth{#1,#2}}

\newcommand{\dist}[1]{\left\| {#1} \right\|}
\newcommand{\norm}[1]{\left\| {#1} \right\|}

\providecommand{\ds}{\displaystyle}
\newcommand{\sep}[1]{\,\left|\, {#1} \bigr.\right.}%
\newcommand{\sepw}[1]{\left|\, {#1} \right.}%

\newcommand{\ProbExt}[2]{\mathop{\mathbf{Pr}}_{#1}\pbrcx{#2}}
\newcommand{\ProbLTR}{\mathbb{P}}%

\newcommand{\Prob}[1]{\mathop{\ProbLTR} \mleft[ #1 \mright]}%
\newcommand{\ProbCond}[2]{\mathop{\ProbLTR}\!\left[%
       #1 \;\middle\vert\; #2 \right]}
\newcommand{\ProbY}[2]{\ProbLTR^{}_{\!#1}\pbrc{#2}}

\newcommand{\ExChar}{\mathbb{E}}%
\newcommand{\ExSym}{\mathop{\ExChar}}%
\newcommand{\Ex}[2][\!]{\ExSym#1\pbrcx{#2}}
\newcommand{\ExExt}[2]{\ExSym_{#1}\!\pbrcx{#2}}
\newcommand{\ExExtS}[2]{\ExSym_{\substack{#1}}\!\pbrcx{#2}}
\newcommand{\ExCond}[2]{\ExSym\!\left[%
       #1 \;\middle\vert\; #2 \right]}

\newcommand{\VarSymb}{\mathop{\mathbb{V}}}
\newcommand{\Var}[1]{\mathop{\mathbb{V}}\!\pbrcx{#1}}

\newcommand{\pbrcx}[1]{\left[ {#1} \right]}%
\newcommand{\nfrac}[2]{{#1}/{#2}}
\newcommand{\cardin}[1]{\left| {#1} \right|}%

\newcommand\CCancel[2][black]{{\renewcommand\CancelColor{\color{#1}}\cancel{#2}}}

\newcommand{\IncGraphPage}[4][]{%
   \IncludeGraphics[page=#4,#1]{\File{#2/#3}}%
}

\newenvironment{enumEx}{
   \begin{compactenumA}
    }{\end{compactenumA}}

\newcommand{\cl}{{\mathrm{cl}}}

\usepackage[amsmath,thmmarks]{ntheorem}
\theoremseparator{.}%

\newtheorem{theorem:unnumbered}{Theorem}[section]
\newtheorem{theorem}{Theorem}[section]

\newtheorem{observation}[theorem]{Observation}
\newtheorem{conjecture}[theorem]{Conjecture}
\newtheorem{lemma}[theorem]{Lemma}
\newtheorem{corollary}[theorem]{Corollary}
\newtheorem{claim}[theorem]{Claim}
\newtheorem{assumption}[theorem]{Assumption}
\newtheorem{tedium}[theorem]{Tedium}

\theoremstyle{plain}%
\theoremheaderfont{\sffamily}
\theorembodyfont{\upshape}%

\newtheorem{defn}[theorem]{Definition}
\newtheorem{definition}[theorem]{Definition}

\newtheorem*{remark_unnumbered}[theorem:unnumbered]{Remark}%
\newtheorem*{remarks}[theorem]{Remarks}%
\newtheorem{remark}[theorem]{Remark}%
\newtheorem{fact}[theorem]{Fact}%
\newtheorem{example}[theorem]{Example}
\newtheorem{exercise}[theorem]{Exercise}
\newtheorem{problem}[theorem]{Problem}
\newtheorem{algorithm}[theorem]{Algorithm}%
\newtheorem{xca}[theorem]{Exercise}
\newtheorem{openproblem}[theorem]{Open Problem}
\newtheorem{exercise_h}[theorem]{Exercise}

% Proof environments
\theoremheaderfont{\em}%
\theorembodyfont{\upshape}%
\theoremstyle{nonumberplain}
\theoremseparator{}
\theoremsymbol{\myqedsymbol}
\newtheorem{proof}{Proof:}
\newtheorem{proofextin}{Proof}
\newtheorem{altproof}{Alternative proof:}

\theoremsymbol{}
\newtheorem{proofnoqed}{Proof:}

\theoremsymbol{\myqedsymbol}
\theorempreskip{0pt}%
\theorempostskip{0pt}%
\newtheorem{proofof}{Proof of\!}%

\newenvironment{proofExt}[1]{{\emph{#1:}}}{\hfill{\hfill\myqedsymbol}}
\newenvironment{proofnoproof}{{}}{{\hfill{\hfill\myqedsymbol}}}

%%%%%%%%%%%%%%%%%%%%%%%%%%%%%%%%%%%%%%%%%%%%%%%%%%%%%%%%%%%%%%%%%%%%
%%%%%%%%%%%%%%%%%%%%%%%%%%%%%%%%%%%%%%%%%%%%%%%%%%%%%%%%%%%%%%%%%%%%
% Macros to embed figures in text

\newcommand{\EBRY}[2]{#1_{\langle#2\rangle}}%

\newcommand{\NiceFrac}[2]{%
   {\raisebox{0.2em}{\ensuremath{#1}}\!\!\!}%
   {\rotatebox[origin=c]{-20}{\Big/}}% Rotated /
   {\! \raisebox{-0.2em}{\ensuremath{#2}} }}%

% end - Macros to embed figures in text
%%%%%%%%%%%%%%%%%%%%%%%%%%%%%%%%%%%%%%%%%%%%%%%%%%%%%%%%%%%%%%%%%%%% 
%%%%%%%%%%%%%%%%%%%%%%%%%%%%%%%%%%%%%%%%%%%%%%%%%%%%%%%%%%%%%%%%%%%%

%%%%%%%%%%%%%%%%%%%%%%%%%%%%%%%%%%%%%%%%%%%%%%%%%%%%%%%%%%%%%%%%%%
% Color in math
%\providecommand{\Mh}[1]{{#1}}%
\providecommand{\IfColorMath}[2]{#2}%

\IfFileExists{.latex_printer_friendly}{\def\GenPrinterVer{1}}{}%

\ifx\GenPrinterVer\undefined
   \IfFileExists{.latex_color}{\def\GenColorMath{1}}{}
\else
   \renewcommand{\IfPrinterVer}[2]{#1}%
\fi

\ifx\GenColorMath\undefined
\else
\renewcommand{\IfColorMath}[2]{#1}%
%\renewcommand{\Mh}[1]{{\textcolor{red}{#1}}}%
\fi
% Color in math - end
%%%%%%%%%%%%%%%%%%%%%%%%%%%%%%%%%%%%%%%%%%%%%%%%%%%%%%%%%%%%%%%%%%

\newcommand{\HSetY}[2]{H\pth{#1, #2}}%
\newcommand{\CHSetY}[2]{C\pth{#1, #2}}%
\newcommand{\Holder}{H\"older\xspace}
\newcommand{\OldOm}{{\mathcal{O}}}
\newcommand{\NewOm}{{\mathcal{N}}}

%%%%%%%%%%%%%%%%%%%%%%%%%%%%%%%%%%%%%%%%%%%%%%%%%%%%%%%%%%%%%%%%%%%%%%%
%---------------------------------------------------------------
\newcommand{\IntRange}[1]{\left\llbracket #1 \right\rrbracket}
\newcommand{\IRX}[1]{\left\llbracket #1 \right\rrbracket}%
\newcommand{\IRY}[2]{\left\llbracket #1:#2 \right\rrbracket}

%---------------------------------------------------------------
%%%%%%%%%%%%%%%%%%%%%%%%%%%%%%%%%%%%%%%%%%%%%%%%%%%%%%%%%%%%%%%%%%%%%%%

\newcommand{\obj}{f}%
\newcommand{\Wg}{\omega}%
\newcommand{\TotalW}{W}%

\newcommand{\rankY}[2]{{#1}_{ \permut{#2}}}%%
\newcommand{\areaX}[1]{\mathrm{area}\pth{#1}}%%

\newcommand{\mybinom}[2]%
{\Bigl(\begin{array}{@{}c@{}}#1\\#2\end{array}\Bigr)}

\providecommand{\Set}[2]{\left\{ #1 \;\middle\vert\; #2 \right\}}
\newcommand*{\biga}[1]{\vcenter{\hbox{\scalebox{1.5}{\ensuremath#1}}}}
\newcommand{\CNF}{\TermI{CNF}\xspace}%
\newcommand{\Tree}{T}%

\DeclareSymbolFont{sfletters}{OML}{cmbrm}{m}{it}

\DeclareMathSymbol{\salpha}{\mathord}{sfletters}{"0B}

%%%%%%%%%%%%%%%%%%%%%%%%%%%%%%%%%%%%%%%%%%%%%%%%%%%%%%%%%%%%%%%%%%%%

\providecommand{\ProblemInd}[2]{\textsf{\textcolor{red}{#1}\index{Problem!#2}}}
\newcommand{\NPProblemInd}[4]{
   \begin{center}
      \vspace{0.6cm}
       \fbox{
          \begin{minipage}{0.9\linewidth}
              \vspace{-0.6cm} \hspace{-0.35cm}
              {\large \bfseries \ProblemInd{#1}{#2}}
              \vspace{+0.2cm}%

              \noindent{\textcolor{blue}{\bfseries \textsf{Instance}}:} {#3}
              
              \noindent{\textcolor{blue}{\bfseries \textsf{Question}:}} {#4}
          \end{minipage}
       }
   \end{center}
   \smallskip
   }

\newcommand{\NPProblemOpt}[4]{\NPProblemInd{#1}{#2}{#3}{#4}}

%%%%%%%%%%%%%%%%%%%%%%%%%%%%%%%%%%%%%%%%%%%%%%%%%%%%%%%%%%%%%%%%%%%%

\newcommand{\Hastad}{H{\aa}stad\xspace}

\newcommand{\ANNbr}{\Term{ANN}\index{approximate nearest neighbor}\xspace}
\newcommand{\LSH}{\TermI{L{S}H}\xspace}

\newcommand{\DS}{\EuScript{D}} % DataStructure
\newcommand{\DSC}{\mathsf{D}} % DataStructure

\newcommand{\distH}[2]{\DistChar_{H}\pth{#1, #2}}
\newcommand{\distHY}[2]{\DistChar_{H}\pth{#1, #2}}

\newcommand{\DSAprxNearNbr}{\ensuremath{\DataStructure{{\approx%
         \text{Near{}}}}}%
   \index{near neighbor!data-structure!approximate}%
}
\newcommand{\DataStructure}[1]{\ensuremath{\DS_{#1}}}
\newcommand{\DistNN}[2]{\DistChar\pth{#1,#2}}
\newcommand{\pntc}{{p}}%

\newcommand{\seq}{{M}}%
\newcommand{\cTimes}{\beta}%

\newcommand{\DSTimes}{L}%
\newcommand{\rad}{r}%

\newcommand{\HypercubeX}[1]{\Hypercube^{#1}}
\newcommand{\kFunctions}[2]{\mathrm{combine}\pth{#1,#2}}
\newcommand{\WEqualChar}{\Bumpeq}
\newcommand{\WEqual}[2]{{#1}\WEqualChar{#2}}
\newcommand{\FamilyAmp}{\FamilyC_{\WEqualChar}}
\newcommand{\PSPACE}{\ComplexityClass{PSPACE}\xspace}%

\newcommand{\PS}{P}%

\newcommand{\PSetFar}{\ensuremath{\PntSet^{}_{\geq}\xspace}}%
\newcommand{\DA}{D}%
\newcommand{\uniqX}[1]{\mathrm{uniq}\pth{#1}}%
\newcommand{\diffX}[1]{\mathop{}\!\mathrm{d}#1}

\newcommand{\Frq}{F}%
\newcommand{\frq}{f}%
\newcommand{\STRM}{\mathcal{S}}%
\newcommand{\strm}{{s}}%

\providecommand{\Turan}{Tur\'an\xspace}
\newcommand{\fGrowX}[1]{G^{}_{\Dim}\pth{#1}}
\newcommand{\CompressFunc}{\mathtt{Compress}}
\newcommand{\RSA}{\Term{RSA}\xspace}
\newcommand{\MST}{\Term{MST}\xspace}
\providecommand{\CClassPP}{\CClass{P{}P}\xspace}
\providecommand{\ID}{\mathrm{id}}

\newcommand{\EncodeFunc}{\text{\ttfamily{Enc}}}
\newcommand{\DecodeFunc}{\text{\ttfamily{Dec}}}
\newcommand{\discX}[1]{\mathrm{disc}\pth{#1}}%

\newcommand{\propertyC}{\psi}%
\newcommand{\propertyX}[1]{\propertyC\pth{#1}}%
\newcommand{\numC}{\Gamma}%

\newcommand{\degreeX}[1]{\mathrm{degree}\pth{#1}}

\newcommand{\ind}{m}

\providecommand{\nequiv}{\not\equiv}
\providecommand{\sqcupplus}{\sqcup^+}

\newcommand{\AlgorithmI}[1]{{\protect\color{RedViolet}{#1}}}

\usepackage{marginnote}
\usepackage{ragged2e} 
% Short, subtle snark macro
\setlength{\marginparwidth}{1.5in}
\setlength{\marginparsep}{10pt}
\newcommand{\snark}[1]{%
   \marginnote{\raggedright\scriptsize\sffamily\color{gray}#1}%
}

\newcommand{\binsY}[2]{\sqcup_{\geq #1}\pth{#2}}
\newcommand{\ballsY}[2]{\text{\Football}_{\geq #1}\pth{#2}}

\newcommand{\BLX}[1]{{{\sqcupplus{}}_{\geq #1}}}

\MakePerPage{footnote} % Resets footnote counter on every page

\allowdisplaybreaks % Allows page breaks inside multiline math environments

\let\Mh\undefined


% 1. Create custom page style for chapter start pages
\makepagestyle{chapfirstpage}

\makefootrule{chapfirstpage}{\textwidth}{0.4pt}{1ex}

% 2. Define footers: {Left}{Center}{Right}
% For two-sided documents (recto/verso):
\makeoddfoot{chapfirstpage}{\footnotesize Sariel Har-Peled}{\thepage}{\today}
\makeevenfoot{chapfirstpage}{\thepage}{Sariel Har-Peled}{}

% (Leave headers empty on chapter start pages)
\makeoddhead{chapfirstpage}{}{}{}
\makeevenhead{chapfirstpage}{}{}{}

% 3. Alias 'chapter' style to use your new style
\aliaspagestyle{chapter}{chapfirstpage}

\newcommand{\pdf}{\textsf{pdf}\xspace}
\newcommand{\cdf}{\textsf{cdf}\xspace}

% -----------------------------------------------------------------------------
% Bibliography & Document Begin
%-----------------------------------------------------------------------------
\IfFileExists{\jobname.bib}{%
   \addbibresource{\jobname.bib}%
}{%
   \addbibresource{shortcuts.bib}%
   \addbibresource{geometry.bib}%
}
\makeindex
\begin{document}
}{\input{prefix_latex}}%
  \fi
\fi
\makeatother
%

\ChapterNumPage{23}{181}

\Chapter{The power of two choices}
\ChapterPrefix{}%
\ChapterGen{}


\QuotePExt{%
   The Peace of Olivia. How sweet and peaceful it sounds! There the
   great powers noticed for the first time that the land of the Poles
   lends itself admirably to partition.%
}%
{The tin drum}%
{G{\"u}nter Grass}%

\bigskip%
\noindent%
Consider the problem of throwing $n$ balls into $n$ bins. It is well
known that the maximum load is $\Theta( \log n / \log \log n)$ with high
probability. Here we show that if one is allowed to pick {$d \geq 2$ bins
   for each ball, and throw it into the one among them that contains the
   fewest balls,} then the maximum load of a bin decreases to
$\Theta( \log \log n/ \log d)$. A variant of this approach leads to a maximum
load of {$O(1) + (\log \log n)/ d$.}

As a concrete example, for $n=10^9$ the maximum load is $13$ with a
single choice (the regular case), compared to $4$ with two choices---see
\figref{n:bins:2:choices}.




%%%%%%%%%%%%%%%%%%%%%%%%%%%%%%%%%%%%%%%%%%%%%%%%%%%%%%%%%%%%%%%%%%%%%
%%%%%%%%%%%%%%%%%%%%%%%%%%%%%%%%%%%%%%%%%%%%%%%%%%%%%%%%%%%%%%%%%%%%%
%%%%%%%%%%%%%%%%%%%%%%%%%%%%%%%%%%%%%%%%%%%%%%%%%%%%%%%%%%%%%%%%%%%%%
\section{Balls and bins with many rows}
\seclab{balls:many:rows}%

\subsection{The game}

Consider throwing $n$ balls into $n$ bins. Every bin can contain a
single ball. As such, as we throw the balls, some balls are
rejected because their assigned bin already contains a ball. We
collect all the rejected balls, and throw them again into a second row
of $n$ bins. We repeat this process until all the balls have found a
good and loving home (i.e., an empty bin). How many rows does one need
before this process is completed?

\subsection{Analysis}

\begin{lemma}
    \lemlab{strong:conc}%
    %
    Let $m = \alpha n$ balls be thrown into $n$ bins.  Let $\Yend$ be the
    number of bins that are not empty after all $m$ balls are thrown
    (here a bin may receive several balls; there are no rejections).
    \begin{compactenumA}
        \item For $\alpha \geq 0$, we have
        \begin{equation*}
            \mu%
            =%
            \Ex{\Yend}%
            \geq %
            \alpha n - \frac{\alpha^2}{2} n.
        \end{equation*}

        \item For $n \geq 4$ and  $c > 1$, we have
        \begin{math}
            \Bigl.  \Prob{\Bigl.\cardin{\Yend - \mu} > \smash{\sqrt{3
                     c m \log n}}} %
            \leq%
            1/n^c.
        \end{math}
    \end{compactenumA}
\end{lemma}
\begin{proof}
    (A) The probability that a bin is empty is $(1-1/n)^{\alpha n}$. As
    such, we have that
    \begin{align*}
        \mu %
        &\EQ%
        \Ex{\Yend}%
        \EQ%
        n\pth{ 1-(1-1/n)^{\alpha n}}%
        \\&
        \GEQ%
        n\pth{ 1-\exp\pth{-1/n}^{\alpha n}}%
        =%
        n\pth{ 1-\exp\pth{ -\alpha }}
        \\&
        \GEQ %
        n\pth{ \alpha - \frac{\alpha^2}{2}},        
    \end{align*}
    using $1-1/n \leq \exp(-1/n)$ and
    \begin{math}
        \exp(-x)
        \leq
        1 - x + \frac{x^2}{2},
    \end{math}
    for $x \geq 0$.

    
    \medskip%
    (B) Let $X_i$ be the index of the bin the $i$\th ball picked. Let
    $Y_i = \ExCond{\Yend}{X_1, \ldots, X_i}$. This is a Doob
    martingale, with $|Y_i - Y_{i-1}| \leq 1$. As such,
    \thmrefY{Azuma's inequality}{\si{azuma:ineq:s:f}} implies, for
    $\lambda = \sqrt{ 3c m \ln n}$, that
    \begin{equation*}
        \Prob{ \bigl. \cardin{\Yend - \Ex{\Yend}} \geq \lambda}%
        \leq%
        2\exp \pth{-\lambda^2/ 2m}%
        \leq%
        1/n^c.
    \end{equation*}
\end{proof}%


\paragraph*{Back to the problem.}

Let $n_i$ be the number of balls thrown in the $i$\th round, so
$n_1 = n$. We track an idealized ratio $\alpha_i$, with $\alpha_1 = 1$, such that
$n_i \approx \alpha_i n$. Inductively, assume that for $i>1$
\begin{equation*}
    n_i = \alpha_i n + O( \sqrt{ \alpha_{i-1} n \log n} ). 
\end{equation*}
By \lemref{strong:conc}, with high probability, the number of balls
stored in the $i$\th row is
\begin{equation*}
    s_i
    \GEQ
    \alpha_i n - \frac{\alpha_i^2}{2} n - O(\sqrt{n_i  \log n})
    \GEQ
    \alpha_i n - \alpha_i^2 n.
\end{equation*}
The last inequality is correct only if $\alpha_i$ is sufficiently large
(roughly $\alpha_i \geq 1/n^{1/3}$), but lets just ignore this issue (it can
easily be resolved but it really not worth the technical tedium it
requires to resolve right now), but we will assume (for our sins) that it
is correct throughout. Thus, we have
\begin{equation*}
    n_{i+1} \EQ n_i - s_i%
    \LEQ%a
    n \alpha_i^2.
\end{equation*}
\begin{defn}
    For a number $x > 0$, we use $\lg x = \log_2 x$.
\end{defn}

\begin{observation}
    Consider the sequence $\alpha_1 =1$, $c = \alpha_2=1-1/e$, and
    $\alpha_{i+1} = \alpha_i^2$, for $i \geq 2$. We have that
    $\alpha_i = c^{2^{i-2}}$. In particular, for
    \begin{equation*}
        \Delta%
        =%
        3 + \lg \log_{1/c} n
        =%
        3 + \lg \frac{\lg n}{\lg (1/c)}
        =
        3 + \lg \lg n  - \lg\lg \frac{1}{1-1/e},
    \end{equation*}
    we have that $\alpha_\Delta = c^{2^{\Delta-2}} < 1/n$.%
\end{observation}

The above observation almost implies that $\Delta$ rows suffice. The
problem is that the above calculations (i.e., the high probability
guarantee in \lemref{strong:conc}) break down when $n_i = O(\log n)$
---that is, when $\alpha_i = O((\log n) / n)$. However, if one throws
$O( \log n)$ balls into $n$ bins, the probability of a single collision
is at most $O( (\log n)^2/n)$. Repeating this for $O(c)$ additional rows, the
probability that any ball is left is at most $1/n^c$.

Done carefully, the above argument implies the following theorem. We omit
the details, since essentially the same analysis is done next in a more
involved setting; the matching lower bound follows from the same
argument.

\begin{theorem}
    Consider the process of throwing $n$ balls into $n$ bins in several
    rounds. Here, any ball that cannot be placed in a round, because its
    chosen bin is already occupied, is promoted to the next round. The
    next round throws all the rejected balls from the previous round into
    a new row of $n$ empty bins.  With high probability, all balls are
    placed within $M = \lg \lg n +\Theta(1)$ rounds, and
    $\lg \lg n - O(1)$ rounds do not suffice.
\end{theorem}



%%%%%%%%%%%%%%%%%%%%%%%%%%%%%%%%%%%%%%%%%%%%%%%%%%%%%%%%%%%%%%%%%%%%%
%%%%%%%%%%%%%%%%%%%%%%%%%%%%%%%%%%%%%%%%%%%%%%%%%%%%%%%%%%%%%%%%%%%%%

\subsection{With only $d$ rows}


\begin{lemma}
    \lemlab{fast:dec}%
    %
    For $\alpha \in (0,1/4]$, let $\gamma_1 = \alpha$, and
    $\gamma_i = 2 \gamma_{i-1}^2$. We have that
    $\gamma_{d+1} \leq \alpha^{(2^{d} + 1)/2}$.
\end{lemma}
\begin{proof}
    The proof, minimal as it may be, is by induction on $i$, showing that
    $\gamma_i \leq \alpha^{(2^{i-1}+1)/2}$ (the case $i=1$ is trivial):
    \begin{equation*}
        \gamma_{i+1}%
        =%
        2 \gamma_i^2 %
        \leq%
        2 \pth{\alpha^{(2^{i-1} + 1)/2}}^2
        =%
        2 \alpha^{(2^{i} + 2)/2}
        \leq%
        \alpha^{(2^{i} + 1)/2},
    \end{equation*}
    since $2\sqrt{\alpha} \leq 1$.
\end{proof}
    

\begin{lemma}
    \lemlab{strong:conc:3}%
    %
    Consider a matrix of $d$ rows with $n$ bins each, and throw $m$ balls
    into it one by one. Every bin can contain only a single ball; a ball
    that finds its chosen bin in the $i$\th row occupied is thrown into
    the next row, and so on, until it finds an empty bin, or it is
    rejected because it failed on the $d$\th row. Let $Y(d,n, m)$ be the
    number of balls that did not get stored in this matrix of bins. We
    have
    \begin{compactenumA}
        \item For $m = \alpha n$, where $\alpha \leq 1/4$ is a constant,
        we have $Y(d,n, \alpha n ) \leq n \alpha^{(2^{d} + 1)/2}$, with high
        probability.

        \item \itemlab{strong:c:3:B} For $m = dn$ balls, we have
        $\Ex{\bigl. Y(d,n, dn)} = O( n \log d) $.

        \item For a constant $c>1$ and $m = c n \log d$ balls, we have
        \[
            \Ex{\bigl. Y(d,n, c n \log d ) }
            \leq%
            n/e^{d/2},
        \]
        assuming $d$ is sufficiently large (so that $D \leq d/2$ in
        the proof below).
    \end{compactenumA}
\end{lemma}

\begin{proof}
    (A) By \lemref{strong:conc}, in expectation, at least
    $s_1 = n\alpha \exp(-\alpha)$ balls are placed in the first row. As such, in
    expectation $n_2 = n \alpha (1 - \exp( -\alpha)) \leq n\alpha^2$ balls get thrown into
    the second row. Using Chernoff's inequality, we get that
    $n_2 \leq 2 \alpha^{2}n$, with high probability. (Why Chernoff's
    inequality can be applied here requires an argument, which we leave
    as an exercise for the reader.)
    Setting $\gamma_1 = \alpha$, and $\gamma_i = 2 \gamma_{i-1}^2$, we get the
    claim (the number of balls rejected by row $d$ is at most
    $\gamma_{d+1} n$) via \lemref{fast:dec}.

    
    (B) As long as we throw $\Omega( n \log d)$ balls into a row, we
    expect by \lemref{strong:conc} that at least $n(1-1/d^{O(1)})$
    balls get stored in this row. As such, let $D = O(\log d )$,
    and observe that the first $d-D$ rows in expectation contain
    $n(d-D)(1-1/d^{O(1)})$ balls. This implies that only $O(D n)$ balls are
    not stored in these first $d-D$ rows, which implies the claim.


    (C) Break the $d$ rows into two groups. The first group of size
    \begin{equation*}
        D= \ceil{(c \log d-1)/(1-1/e)} + 1 = O(\log d),
    \end{equation*}
    and the second group is the remaining rows. As long as the number of
    balls arriving at a row is larger than $n$, we expect at least
    $n(1-1/e)$ balls to be stored in that row. As such, after the first
    $D$ rows, we expect the number of remaining balls to be $\leq
    n$. Indeed, after $i$ rows, then the expected number of balls moving
    on to the $(i+1)$\th row is at most
    \begin{equation*}
        n_{i+1} = cn \log d - i n(1-1/e).
    \end{equation*}
    Solving for $n_{i+1} \leq n$, we have
    \begin{math}
        cn \log d - i n(1-1/e) \leq n \implies i (1-1/e) \geq c \log
        d-1 \implies i \geq (c \log d-1)/(1-1/e) \geq D - 1.
    \end{math}    
    As such, $n_i \leq n$, for $i \geq D$.
    
    Applying the same calculation to the second group, the number of
    balls arriving at the $(D+i)$\th row, in expectation, is at most
    $n/e^i$. In particular, we get that the number of balls that fail to
    be placed is at most $n/e^{d-D} \leq n/e^{d/2}$.
\end{proof}


%%%%%%%%%%%%%%%%%%%%%%%%%%%%%%%%%%%%%%%%%%%%%%%%%%%%%%%%%%%%%%%%%%%%%
%%%%%%%%%%%%%%%%%%%%%%%%%%%%%%%%%%%%%%%%%%%%%%%%%%%%%%%%%%%%%%%%%%%%%
%%%%%%%%%%%%%%%%%%%%%%%%%%%%%%%%%%%%%%%%%%%%%%%%%%%%%%%%%%%%%%%%%%%%%
\section{The power of two choices}
\seclab{power:2:choices}%

\paragraph*{Making $d$ choices.}
Let us throw $n$ balls into $n$ bins. For each ball, we first pick
randomly $d\geq 2$ bins, and place the ball in the bin (among these $d$
bins) that currently contains the smallest number of balls (here, a bin
might contain an arbitrary number of balls). If there are several bins
with the same minimum load, we resolve it arbitrarily.

Here, we will show that the maximum number of balls in any bin is bounded
by $O\bigl( \tfrac{\log \log n}{\log d} \bigr)$ with high probability at
the end of this process. For $d=1$, which is the regular balls into bins
setting, we have already seen that this quantity is
$\Theta\bigl( \tfrac{\log n}{ \log \log n}\bigr)$, so this result is quite
surprising.


% \subsection{Some technical lemmas}

\subsection{Upper bound}

\begin{defn}
    The \emphOnly{load}\index{load} of a bin is the number of balls in it. The
    \emphOnly{height}\index{height} of a ball is the load of the bin it was inserted
    into, just after it was inserted.
\end{defn}

%\newcommand{\UBLX}[1]{u_{\geq #1}}%

Some notations:
\begin{compactenumA}
    \smallskip%
    \item $\beta_i$: An upper bound on the number of bins that have load at
    least $i$ by the end of the process (defined concretely in
    \lemref{how:betas:behave}).
    
    \smallskip%
    \item $h(t)$: The height of the $t$\th ball.

    \smallskip%
    \item $\binsY{i}{t}$: Number of bins with load at least $i$ at
    time $t$.

    \smallskip%
    \item $\ballsY{i}{t}$: Number of balls with height at least $i$ at
    time $t$.
\end{compactenumA}

\begin{observation}
    $\binsY{i}{t} \leq \ballsY{i}{t}$.
\end{observation}

\newcommand{\BadEvent}{\mathcal{B}}%
\newcommand{\GoodEvent}{\mathcal{G}}%


Let $\BLX{i} = \binsY{i}{n}$ be the number of bins, at the end of the
process, that have load $\geq i$.

\begin{observation}
    \obslab{t:c:stupid}%
    %
    Since every bin counted in $\BLX{i}$ contains at least $i$ balls,
    and there are $n$ balls, it follows that $\BLX{i} \leq n/i$.  In
    particular, we have $\BLX{4} \leq n/4$.
\end{observation}


\begin{lemma}
    \lemlab{how:betas:behave}%
    %
    Let $\beta_1 = n, \beta_2 = n/2, \beta_3=n/3$, and
    $\beta_4 = n/4$, and let
    \begin{equation*}
        \beta_{i+1} = 2 n \pth{\beta_i/n}^d,
    \end{equation*}
    for $i\geq 4$.  Fix a constant $c > 1$, and let $I$ be the last index
    such that $\beta_{I} \geq 16c\ln n$. Then, with probability
    $\geq 1-1/n^c$, we have that
    \begin{compactenumA}
        \smallskip%
        \item $\BLX{i} \leq \beta_i$, for $i=1,\ldots, I$.
        \smallskip%
        \item $\BLX{I+1} \leq c' \log n$, for some constant $c'$. %

        \smallskip%
        \item For every integer $j > 0$, and any constant $\eps > 0$, we
        have
        \[
        \Prob{ \BLX{I+1+j} > 0} \leq O(1/n^{(d-1-\eps)j}).
        \]

        \smallskip%
        \item The maximum load of a bin is $I+O(c)$.
    \end{compactenumA}
\end{lemma}

\begin{proof}
    (A) The claim for $i=1,2,3,4$ follows from
    \obsref{t:c:stupid}.

    Let $\BadEvent_i$ be the bad event that $\BLX{i} > \beta_i$, for
    $i=1,\ldots, n$. The following analysis is conditioned on none of these
    bad events happening. Let
    $\GoodEvent_k = {\cap_{i=1}^k \overline{\BadEvent_i}}$ be the good
    event.  Fix $i < I$, and let $Y_t$ be the indicator variable of the
    event $h(t) \geq i+1$. Below, all probabilities are conditioned on
    $\GoodEvent_i$, which we omit from the notation. We have that
    \begin{equation*}
        \tau_j = \Prob{Y_j=1 } \leq p_i%
        \qquad\text{ for }\qquad%
        p_i=\pth{ \beta_i /n}^d,
    \end{equation*}
    as all $d$ probes must hit bins of load at least $i$, and there are
    at most $\beta_i$ such bins. Summing over the $n$ balls gives
    $\Ex{ \ballsY{i+1}{n} } \leq p_i n$.  The variables
    $Y_1, \ldots, Y_n$ are not independent, but consider a variable
    $Y_j'$ that is $1$ if $Y_j=1$, or if $Y_j=0$, then $Y_j'$ is $1$ with
    probability $\frac{p_i-\tau_j}{1-\tau_j}$. Clearly, the variables
    $Y_1',\ldots, Y_n'$ are independent, and
    $\sum_{j=1}^n Y_j' \geq \sum_{j=1}^n Y_j$. Importantly,
    \begin{equation*}
        \forall j
        \qquad
        \Prob{Y_j' =1}
        =
        p_i.
    \end{equation*}
    \remX{E1 in review.md: the display should be about $Y_j'$, and even then the construction gives $\Prob{Y_j'=1} = \tau_j + (1-\tau_j)(p_i - \tau_j) \neq p_i$; the extra probability must be $(p_i-\tau_j)/(1-\tau_j)$. Also $Y_j' = Y_j \vee W_j$ with $\tau_j$ the marginal probability inherits the dependence of the $Y_j$, so ``clearly independent'' is unjustified. Fix: use the conditional probability $q_j = \ProbCond{Y_j = 1}{\text{history}}$ in place of $\tau_j$; then $\Prob{Y_j'=1 \mid \text{history}} = p_i$ and the $Y_j'$ are i.i.d.}\remX{G1 in review.md: conditioning on $\GoodEvent_i$, an event about the \emph{end} of the process, invalidates the per-ball bound $\Prob{Y_j = 1} \leq p_i$ (exact toy: $n=3$, $d=2$, $i=1$, $\beta_1 = 2$ gives $\Prob{Y_3 = 1 \mid \GoodEvent} = 33/41 > 4/9 = p_1$). Standard fix (ABKU): define $Y_t = 1$ iff $h(t) \geq i+1$ \emph{and} $\binsY{i}{t-1} \leq \beta_i$; then $\Prob{Y_t = 1 \mid \text{history}} \leq p_i$ always, $\sum_t Y_t = \ballsY{i+1}{n}$ on $\GoodEvent_i$, and one bounds $\Prob{\BadEvent_{i+1} \cap \GoodEvent_i}$. Minor: in the align below, $\Prob{\BadEvent_{i+1}} = \Prob{\ballsY{i+1}{n} > \beta_{i+1}}$ should be $\leq$ (E6).}


    For
    $i < I$, setting
    \begin{equation*}
        \beta_{i+1} = 2np_i  = 2 n \pth{\beta_i/n}^d,        
    \end{equation*}
    we have, by \lemrefY{Chernoff's inequality}{chernoff:small:delta}
    applied to the independent variables $Y_1', \ldots, Y_n'$, that
    \begin{align*}
      \alpha_{i+1}%
      =%
      \Prob{\BadEvent_{i+1}}%{\cap_{k=1}^i\overline{\BadEvent_{1}}}%
      &=% 
      \Prob{\ballsY{i+1}{n} > \beta_{i+1}}%
      \\&
      =%
      \Prob{\ballsY{i+1}{n} > 2np_i}
      \\&
      \leq 
      \Prob{ \textstyle{\sum_{t=1}^n} Y_t' > (1+1)n p_i}%
      \\&
      % \\&%
      \leq%
      \exp( -n p_i/4 )
      =
      \exp(-\beta_{i+1}/8) < 1/n^{2c}.
    \end{align*}

    (B) We have $\beta_{I+1} \leq 16 c \log n$.  Setting
    $\Delta = 2e \cdot 16 c \log n$, and conditioning on the good event
    $\GoodEvent_I$,%
    consider the sequence $Y_1', \ldots, Y_n'$ as above, where $Y_t$ is the
    indicator that the $t$\th ball has height $\geq I+1$. Arguing as above
    (with $i = I$), for $Y' = \sum_{t=1}^n Y_t'$, we have
    $\Ex{Y'} \leq \beta_{I+1}$. As such, we have
    \begin{equation*}
        \Prob{\BLX{I+1} > \Delta}
        \leq%
        \Prob{\ballsY{I+1}{n} > \Delta}%
        \leq%
        \Prob{Y' > \frac{\Delta}{\Ex{Y'}} \Ex{Y'} }           
        \leq%
%        \Prob{\sum_t Y_t' > \frac{\Delta}{\beta_{I+1}  } \beta_{I+1}}%
%        \leq%      
        2^{ - \Delta}
        \leq%
        \frac{1}{n^c},        
    \end{equation*}
    by \lemref{chernoff:delta:big}, as $\Ex{ Y' } \leq \beta_{I+1}$,
    and $\Delta/\beta_{I+1} > 2e$.
    
%    if $\beta_{i+1} > 16c \ln n$.  If $\beta_{i+1}$

    It remains to bound the probability of the good event on which all of
    the above was conditioned. We have that
    \begin{equation*}
        \Prob{\GoodEvent_{I+1}}%
        =%
        \prod_{\ell=4}^{I+1} \ProbCond{\overline{\BadEvent_{\ell+1}}}{
           {\cap_{k=1}^\ell\overline{\BadEvent_k}}}
        =%
        \prod_i
        (1-\alpha_i)
        \geq
        1-1/n^{c-1},        
    \end{equation*}
    since $I \leq n$.\remX{E3 in review.md: (i) the product should run over $\ell = 4, \ldots, I-1$ (the factor for $\ell = I+1$ concerns $\BadEvent_{I+2}$, never bounded, and $\overline{\BadEvent_{I+1}} = \{\BLX{I+1} \leq \beta_{I+1}\}$ is not what (B) proves, as $\beta_{I+1}$ can be $< 1$); (ii) $1 - 1/n^{c-1}$ does not give the lemma's $1 - 1/n^c$, and the $2^{-\Delta}$ term of (B) is missing. Fix: union bound $\sum_{i=4}^{I-1} \Prob{\BadEvent_{i+1} \cap \GoodEvent_i} + \Prob{\{\ballsY{I+1}{n} > \Delta\} \cap \GoodEvent_I} \leq (I-4) n^{-2c} + 2^{-\Delta} \leq n^{-c}$, using $I = O(\log\log n)$.}

    \medskip%
    (C) \chgY{Observe that $\binsY{i+1}{n} \leq \binsY{i}{n}$. As such, for
    all $j > 0$, we have that}{Condition on the event $\BLX{I+1} \leq \Delta$ of (B). Since
    $\binsY{i+1}{n} \leq \binsY{i}{n}$, for all $j > 0$ we have that}
    \begin{math}
        \binsY{I+1+j}{n}%
        \leq \ballsY{I+1}{n}%
        \leq \Delta = 2e \cdot 16 c \log n,
    \end{math}
    by (B).  As such, we have
    \begin{equation*}
        \Ex{ \ballsY{I+1+j}{n} }%
        \leq%
        n(\Delta/n)^d %
        =%
        O( \log^d n / n^{d-1})
        =%
        O(1/n^{d-1-\eps}) \ll 1,
    \end{equation*}
    for $\eps > 0$ an arbitrary constant, and $n$ sufficiently large.
    Using Markov's inequality, we get that
    $q = \Prob{\ballsY{I+1+j}{n} \geq 1} = O(1/n^{d-1-\eps})$. The
    probability that the first $j$ such rounds fail (i.e., that
    $\ballsY{I+1+j}{n} > 0$) is at most $q^j$, as claimed.\remX{E2 in review.md: there are no independent rounds here; the events $\{\ballsY{I+1+j}{n} \geq 1\}$ are nested in a single run, and the bound $\Ex{\ballsY{I+1+j}{n}} \leq n(\Delta/n)^d$ is the same for every $j$, so nothing yields an exponent $j$. The claim is true; a correct argument: a ball reaches height $\geq I+1+j$ only if its $d$ probes hit bins of load $\geq I+j$, whose number is $\leq \min(\Delta, Z_{j-1})$ where $Z_{j-1} = \ballsY{I+j}{n}$; so $\Ex{Z_j} \leq \Delta^{d-1} n^{1-d} \Ex{Z_{j-1}}$, giving $\Ex{Z_j} \leq \Delta^{(d-1)j+1}/n^{(d-1)j}$ and Markov. See review.md.}

    \smallskip%

    (D) \chgY{This follows immediately by picking $\eps = 1/2$, and then
    using (C) with $j = O(c)$.}{Take $\eps = 1/2$ and $j = \ceil{2c}$ in (C): as $d \geq 2$, we have
    $(d-1-\eps) j \geq c$, so the probability that level $I+1+j$ is
    non-empty is at most $O(1/n^{c})$, and the maximum load is at most
    $I+1+j = I + O(c)$.}
\end{proof}


\begin{lemma}
    \lemlab{beta:i:good}%
    %
    For \chgY{$i=4,\ldots, I$}{every $i \geq 4$}, we have that $\beta_i \leq n/2^{d^{i-4}+1}$.
\end{lemma}
\begin{proof}
    The proof is by induction. For $i=4$, we have $\beta_4 \leq n/4$,
    as claimed. Otherwise, we have
    \begin{equation*}
        \beta_{i+1}%
        =%
        2 n \pth{\beta_i/n}^d%
        \leq%
        2n \pth{1/2^{d^{i-4} + 1}}^d
        =%
        n /2^{d^{i+1-4} + d - 1}
        \leq%
        n /2^{d^{i+1-4} + 1}.
    \end{equation*}
\end{proof}

\begin{theorem}
    When throwing $n$ balls into $n$ bins, with $d$ choices, with
    probability \chgY{$\geq 1-1/n^{O(1)}$}{$\geq 1-1/n^{c}$, for any constant $c>1$ (the $O(1)$ below depends on $c$)}, we have that the maximum load of
    a bin is $O(1) + \tfrac{\lg \lg n}{\lg d}$.
\end{theorem}
\begin{proof}
    By \lemref{how:betas:behave}\chgY{, with the desired probability the
    $\beta_i$s bound the load in the bins for $i \leq I$.}{ (D), with the desired probability the
    maximum load is $I + O(c)$, so it suffices to bound $I$.}  By
    \lemref{beta:i:good}, \chgY{it follows that for
    $I = O(1) + \tfrac{\lg \lg n}{\lg d}\Bigr.$, we have that
    {$\beta_{I+1} \leq 16 c \ln n$}.  ~This gives the desired bound.}{$\beta_i < 16 c \ln n$ once
    $2^{d^{i-4}} > n/(16 c\ln n)$, that is, for
    $i = 4 + \ceil{\log_d \lg n} = O(1) + \tfrac{\lg \lg n}{\lg d}$.
    Since the $\beta_i$ are decreasing, $I$ is at most this value, which
    gives the desired bound.}
\end{proof}

It is not hard to verify that the upper bounds $\beta_i$ are tight up to
constants, so the maximum load is, up to an additive constant, the
same. We state this without proof.

\begin{theorem}
    When throwing $n$ balls into $n$ bins, with $d$ choices (where the
    ball is placed into the bin with the least load), with probability
    $\geq 1-o(1/n)$, we have that the maximum load of a bin is at
    least $\tfrac{\lg \lg n}{\lg d} - O(1)$.
\end{theorem}

\subsection{Applications}

As a direct application of the power of two choices, we can use it for
open hashing (chaining), where we use two hash functions, and place an
element in the bucket of the hash table with fewer elements. By the
above, this improves\, with high probability, the worst-case search time
from $O( \log n / \log \log n)$ to $O( \log \log n)$. This comes at the
cost of doubling the time it takes to do lookup on average.


%%%%%%%%%%%%%%%%%%%%%%%%%%%%%%%%%%%%%%%%%%%%%%%%%%%%%%%%%%%%%%%%%%%%%%%%%

\subsection{The power of restricted $d$ choices: %
   Always go left}
\seclab{restricted:d:choices}%

\noindent%
\textbf{The always go left rule.} %
\chgY{}{Assume that $d$ divides $n$. }Consider throwing a ball into $n$ bins (which might already have some
balls in them) as follows -- you pick uniformly a number
$X_i \in \IRX{n/d}$, for $i=1,\ldots, d$. Next, you try locations
$Y_1, \ldots, Y_d$, where $Y_j = X_j + (j-1)(n/d)$, for $j=1,\ldots, d$. Let
$L_j$ be the load of bin $Y_j$, for $j=1,\ldots, d$, and let
$L = \min_j L_j$ be the minimum load \chgY{of any bin.}{among these $d$ bins.}\remX{W6: the random numbers are indexed $X_i$ but used as $X_j$; also $Y_j$ and $\tau$ clash with $Y(d,n,m)$ and $\tau_j$ used nearby.} Let $\tau$ be the minimum
index such that $L_j = L$. We throw the ball into $Y_\tau$.

What the above scheme does is to partition the $n$ bins into $d$
groups each of size $n/d$, placed from left to right. We pick a bin
uniformly from each group, and always throw the ball in the leftmost
location that realizes the minimum load.

The following proof is only a sketch: it reduces the process to the
row-filling game and skips the concentration details.

\begin{theorem}
    When throwing $n$ balls into $n$ bins, using the always-go-left rule,
    with $d$ groups of size $n/d$, the maximum load of a bin is
    $O(1) + \tfrac{\lg \lg n}{d}$, with high probability.
\end{theorem}
\begin{proof}[Sketch.]
    We consider each of the $d$ groups to be a row in the matrix being
    filled. So each row has $n/d$ entries, and there are $d$ rows. We
    can now think about the above algorithm as first trying to place
    the ball in the first row (if there is an empty bin), otherwise,
    trying the next row and so on. If all the $d$ locations are full,
    in the row filling game we fail to place this ball.  By
    \lemref{strong:conc:3} \itemref{strong:c:3:B}, we have that the
    number of unplaced balls is
    $\Ex{\bigl. Y\bigl(d,n/d, (n/d) d\bigr)} = O\bigl( (n/d) \log d
    \bigr) $.  Thus, we have that the number of balls that get placed
    as the first ball in their bin is
    \begin{equation*}
        \geq n\Bigl(1-\frac{O(\log d)}{d}\Bigr),
    \end{equation*}
    and the height of these balls is one.

    \chgY{Applying the same argumentation for balls of height
    2---}{Applying the same argument to balls of height $2$,}
    \lemref{strong:conc:3} (C) implies that\chgY{}{, in expectation,} at most {$dn/e^{d/2}$}
    balls have height strictly larger than $2$.\remX{U2 in review.md: (C) with $n/d$ bins per row gives $(n/d)/e^{d/2}$, not $dn/e^{d/2}$; and the reduction to the row game is only exact at height~1, see G5.}

    \lemref{strong:conc:3} (A) \chgY{implies that now we can repeat the same
    analysis as the power of two choices, the critical difference is
    that every one of the $d$ groups, behaves like a separate height.}{now lets us repeat the same analysis
    as for the power of two choices. The critical difference is that
    each of the $d$ groups acts as its own height level, so the analysis
    advances $d$ levels per unit of load.}\remX{G5 in review.md: for heights $\geq 2$ the process is not the row game of \lemref{strong:conc:3} (probes are uniform over the group, not over the bins of load $\geq i-1$), and the row game is the optimistic side. The conclusion is true (V{\"o}cking), but the lemma does not imply it; say so explicitly.}
    Since the maximum height is $O( \log \log n)$ in the regular
    analysis, this implies that we get $O( (\log \log n)/d)$ maximum
    load, with high probability.
\end{proof}

\begin{figure}
    \centering%
       \begin{tabular}{|c|c|c|c|}
         \hline
         \# balls in bin & Regular & $2$-choices & $2$-choices+go left\\ 
         \hline
         \hline
         0 & 369,899,815  &  240,525,897  &  228,976,604 \\
         1 & 365,902,266  &  528,332,061  &  546,613,797 \\
         2 & 182,901,437  &  221,765,420  &  219,842,639 \\
         3 &  61,604,865  &    9,369,389  &    4,566,915 \\
         4 &  15,760,559  &        7,233  &           45 \\
         5 &   3,262,678  &               &              \\
         6 &     568,919  &               &              \\
         7 &      86,265  &               &              \\
         8 &      11,685  &               &              \\
         9 &       1,347  &               &              \\
         10 &         143  &               &              \\
         11 &          17  &               &              \\
         12 &           2  &               &              \\
         13 &           2  &               &
         \\
         \hline
       \end{tabular}%
       \caption{Simulation of the three schemes described here. This was
          done with $n=1,000,000,000$ balls thrown into $n$ bins. Since
          $\log \log n$ is so small (i.e., $\approx 3$ in this case), there do
          not seem to be any reasonable cases where there is a
          significant difference between $2$-choices and the go-left
          variant. In the simulations, the \emph{go-left} variant always
          has a somewhat better distribution, as shown above. The randomness used in the implementation is the standard one provided by \texttt{C++}, and is thus not very random.}
 \figlab{n:bins:2:choices}
\end{figure}


\section{Avoiding terrible choices}
\seclab{a:t:c}

Interestingly, one can prove (see the bibliographical notes) that two
choices are not really necessary. Indeed, consider the variant where the
$i$\th ball \chgY{randomly chooses a random location $r_i$.}{chooses a
   uniformly random location $r_i$.} The ball is then placed in the bin
with the least load among the bins $r_i$ and $r_{i-1}$ (the first ball
inspects only a single bin -- $r_1$). It is not difficult to show that the
above analysis applies in this setting, and the maximum load is
$O( \log \log n)$ -- despite making only $n$ choices for $n$
balls. Intuitively, what is going on is that the power of two choices
lies in the ability to avoid following a horrible, no good, terrible
choice, by having an alternative. This alternative choice does not have
to be of quite the same quality as the original choice---it can be stolen
from the previous ball, etc.


%%%%%%%%%%%%%%%%%%%%%%%%%%%%%%%%%%%%%%%%%%%%%%%%%%%%%%%%%%%%%%%%%%%%%%%%%%%
%%%%%%%%%%%%%%%%%%%%%%%%%%%%%%%%%%%%%%%%%%%%%%%%%%%%%%%%%%%%%%%%%%%%%%%%%%%
%%%%%%%%%%%%%%%%%%%%%%%%%%%%%%%%%%%%%%%%%%%%%%%%%%%%%%%%%%%%%%%%%%%%%%%%%%%

\section{Escalated choices}

A variant that seems to work even better in practice\chgY{,}{ (than $2$-choices)} is the following
\emphOnly{escalated choices}\index{escalated choices} algorithm: The idea is to try more than one
bin only if you need to. To this end, try a random bin. If it is
empty, then the algorithm stores the ball in it. Otherwise, the
algorithm tries harder. In the $j$\th iteration, for $j >1$, the
algorithm picks \chgY{a random location. If any of the $j$ locations have}{one more random location (so $j$ locations have been inspected
so far). If any of these $j$ locations have}
load $< \ceil{j/2}$, then the algorithm places the ball in the
min-load bin among these. Otherwise, the algorithm continues to the
next iteration.

Experiments show that, on average, this algorithm probes only $1.96$
bins per ball (thus, making fewer probes than $2$-choices). In this
setting, the experiments show that $4$-choices with go-left do
better, but if one uses the threshold $< \ceil{j/3}$, then the average
number of probes is $2.30179$, while again having a better
performance.\remX{W7 in review.md: say in what measure ``better'' is meant (maximum load? load histogram?) and what is compared with what; also $1.96$ vs.\ $2.30179$ precision.} The intuition is that a sequence of really bad choices
is rare, and one can afford to try harder in such cases to get out of
them.

A theoretical analysis of this variant should be interesting.

(I ``invented'' this variant, but it might already be known.)

%%%%%%%%%%%%%%%%%%%%%%%%%%%%%%%%%%%%%%%%%%%%%%%%%%%%%%%%%%%%%%%%%%%%%%%%%%%
%%%%%%%%%%%%%%%%%%%%%%%%%%%%%%%%%%%%%%%%%%%%%%%%%%%%%%%%%%%%%%%%%%%%%%%%%%%
%%%%%%%%%%%%%%%%%%%%%%%%%%%%%%%%%%%%%%%%%%%%%%%%%%%%%%%%%%%%%%%%%%%%%%%%%%%


\section{Bibliographical notes}

The multi-row balls into bins (\secref{balls:many:rows}) is from the
work by Broder and Karlin \cite{bk-mah-90}. The power of two choices
(\secref{power:2:choices}) is from Azar \etal~\cite{abku-ba-90}.

The restricted $d$ choices structure, the always go-left rule, described
in \secref{restricted:d:choices}, is due to V{\"{o}}ck{}i{n}g
\cite{v-hshlb-03}.

The memory-based variant of \secref{a:t:c}, where each ball chooses
between a fresh random bin and the bin inspected by the previous ball,
was introduced and analyzed by Mitzenmacher \etal~\cite{mps-bbm-02}.



\StandAloneMode{
   \section{From previous lectures}
   
   \IncFragment{Azuma_stronger}%
   \IncFragment{chernoff_small_delta}%
   \IncFragment{d_approx_basic}%
   \IncFragment{chernoff_delta_big}%
}%

\ChapterEnd{}

% END
%%%%%%%%%%%%%%%%%%%%%%%%%%%%%%%%%%%%%%%%%%%%%%%%%%%%%%%%%%%%%%%%%%%%%%%%%%%%%% 
%%%%%%%%%%%%%%%%%%%%%%%%%%%%%%%%%%%%%%%%%%%%%%%%%%%%%%%%%%%%%%%%%%%%%%%%%%%%%% 
%%%%%%%%%%%%%%%%%%%%%%%%%%%%%%%%%%%%%%%%%%%%%%%%%%%%%%%%%%%%%%%%%%%%%%%%%%%%%%}
